% La cybercriminalité 1 : les fake news — Allemand Tle A — Cours signé OnBuch+
% Programme officiel MINESEC. Compiler avec : tectonic AL17-la-cybercriminalite-1-les-fake-news.tex
\newcommand{\DOCMATIERE}{Allemand}
\newcommand{\DOCNIVEAU}{Tle A}
\newcommand{\DOCMODULE}{Chapitre 4 : Média et communication}
\newcommand{\DOCLECON}{AL17}
\newcommand{\DOCTITRE}{La cybercriminalité 1 : les fake news}
% OnBuch+ — préambule « cours signés » ENRICHI (compilé avec tectonic / XeLaTeX)
% Système visuel « Pop » d'OnBuch+ : fond crème (jamais blanc), bordures encre
% épaisses, ombres « dures » décalées, titres Archivo Black en capitales.
% Version MATHÉMATIQUES : graphes (pgfplots/tikz), boîtes pédagogiques
% (définition, propriété, méthode, attention, le savais-tu, exercice résolu,
% exercice, TP, bilan), page de garde et signature OnBuch+ ancrée en bas.
%
% Macros attendues dans main.tex AVANT \input{preamble} :
%   \DOCMATIERE  \DOCNIVEAU  \DOCMODULE  \DOCLECON (numéro)  \DOCTITRE
\documentclass[11pt]{article}

\usepackage{fontspec}
\usepackage[ngerman,french]{babel} % français = langue principale ; l'allemand via \all{..} / textebox
\usepackage[a4paper,margin=2cm,top=3.4cm,bottom=2.8cm,headheight=1.4cm,footskip=1.3cm]{geometry}
\usepackage{amsmath,amssymb,amsfonts}
\usepackage{mathtools} % \xmapsto, \coloneqq... souvent utilisés par les modèles
\usepackage{cancel} % simplifications barrées (bilans de forces)
% \abs{z} (module, valeur absolue) et \norm{u} (norme), utilisés par les modèles
\providecommand{\abs}{}\let\abs\relax\DeclarePairedDelimiter{\abs}{\lvert}{\rvert}
\providecommand{\norm}{}\let\norm\relax\DeclarePairedDelimiter{\norm}{\lVert}{\rVert}
\usepackage[table]{xcolor} % [table] : fournit aussi \rowcolors (alternance de lignes)
\usepackage{pagecolor}
\usepackage{tikz}
\usetikzlibrary{arrows.meta,positioning,calc,shapes.geometric,shapes.misc,shapes.arrows,decorations.pathmorphing,decorations.pathreplacing,decorations.markings,patterns,shadows,mindmap,trees,fit,backgrounds,matrix,intersections,angles,quotes}
\usepackage{pgfplots}
\usepackage[version=4]{mhchem} % équations nucléaires \ce{^{235}_{92}U}
\usepackage[european,siunitx]{circuitikz} % schémas électriques
\pgfplotsset{compat=1.18}
\usepgfplotslibrary{fillbetween,groupplots} % aires entre courbes (calcul intégral)
\usepackage{enumitem}
\usepackage{fancyhdr}
% [most] charge toutes les bibliothèques tcolorbox (listings, minted, raster,
% documentation...) et gonfle inutilement la mémoire TeX sur un document long
% (~300 boîtes) : on ne charge que ce qui est réellement utilisé.
\usepackage[skins,breakable]{tcolorbox}
\usepackage{graphicx}
\usepackage{colortbl}
\usepackage{booktabs}
\usepackage{tabularx}
\usepackage{diagbox} % case d'en-tête diagonale des tableaux à double entrée
% Colonnes extensibles centrée / gauche / droite (C, L, R), souvent utilisées par les modèles
\newcolumntype{C}{>{\centering\arraybackslash}X}
\newcolumntype{L}{>{\raggedright\arraybackslash}X}
\newcolumntype{R}{>{\raggedleft\arraybackslash}X}
\usepackage{array}
\usepackage{multicol}
\usepackage{multirow}
\usepackage{siunitx}
% parse-numbers=false : accepte aussi \SI{2,5 \times 10^{-3}}{...}
\sisetup{locale=FR,output-decimal-marker={,},group-minimum-digits=4,parse-numbers=false}
\DeclareSIUnit\ohm{\ensuremath{\symup{\Omega}}} % U+2126 absent de la police
\DeclareSIUnit\division{div} % réglages d'oscilloscope (ms/div, V/div)
\DeclareSIUnit\tr{tr} % tours (vitesse de rotation, ex. \SI{1500}{\tr\per\minute})
\DeclareSIUnit\molar{M} % concentration molaire (ex. \SI{3}{\molar}, \SI{1}{\micro\molar})
\DeclareSIUnit\an{an} % année (le modèle confond parfois avec \year, un compteur TeX) — ex. \SI{5}{\kilo\meter\per\an}
\DeclareSIUnit\FCFA{FCFA} % franc CFA (ex. \SI{500}{\FCFA\per\kilo\gram})
\DeclareSIUnit\degreeBrix{\textdegree Brix} % teneur en sucre d'un sirop/jus (réfractométrie)
\DeclareSIUnit\month{mois} % le modèle utilise parfois \month, un compteur TeX, comme unité de durée
\usepackage{qrcode}
% \degree : commande fréquemment utilisée par les modèles pour un angle ou une
% température en dehors de \SI{}{} (non fournie par les paquets chargés).
\providecommand{\degree}{^{\circ}}
% \qed (fin de démonstration), utilisé par les modèles sans amsthm
\providecommand{\qed}{\hfill$\square$}
% \barre : double barre des valeurs interdites dans un tableau de variations (tkz-tab)
\providecommand{\barre}{\ensuremath{\|}}
% environnement proof (amsthm non chargé) utilisé par les modèles
\ifdefined\proof\else\newenvironment{proof}[1][Démonstration]{\par\noindent\textit{#1.}\ }{\qed\par}\fi
\usepackage{lastpage}
\usepackage{hyperref}
\hypersetup{hidelinks}

% ============================================================
% Palette « Pop » OnBuch+ — mêmes hex que la classe P (plus_ui.dart)
% ============================================================
\definecolor{poporange}{HTML}{F59321}
\definecolor{popdark}{HTML}{E07A0C}
\definecolor{popcream}{HTML}{FFF6E8}
\definecolor{popink}{HTML}{1C1714}
\definecolor{poppurple}{HTML}{7A5AE0}
\definecolor{popgreen}{HTML}{2E9E6A}
\definecolor{poppink}{HTML}{E0457B}
\definecolor{popblue}{HTML}{2F7DE1}
\definecolor{popgold}{HTML}{C98A12}
\definecolor{popmuted}{HTML}{8A7F76}
\definecolor{popline}{HTML}{EADFD2}
\definecolor{popgray}{HTML}{8A7F76}
\colorlet{popgrey}{popgray}
% Alias tolérants (le modèle nomme parfois les couleurs différemment)
\colorlet{popwhite}{white}
\colorlet{popblack}{popink}
\colorlet{popred}{poppink}
\colorlet{popyellow}{popgold}
% Teintes claires (fonds de tableaux/encadrés) — le modèle nomme parfois
% les couleurs avec un suffixe L (ex. popblueL) sans les avoir définies.
\colorlet{poporangeL}{poporange!20}
\colorlet{popdarkL}{popdark!20}
\colorlet{poppurpleL}{poppurple!20}
\colorlet{popgreenL}{popgreen!20}
\colorlet{poppinkL}{poppink!20}
\colorlet{popblueL}{popblue!20}
\colorlet{popgoldL}{popgold!20}
\colorlet{popredL}{popred!20}
\colorlet{popgrayL}{popgray!20}
% Teintes claires pour fonds de tableaux / zones
\colorlet{popblueL}{popblue!10!white}
\colorlet{poporangeL}{poporange!14!white}
\colorlet{popgreenL}{popgreen!12!white}
\colorlet{poppurpleL}{poppurple!12!white}
\colorlet{poppinkL}{poppink!12!white}
\colorlet{popgoldL}{popgold!14!white}

\pagecolor{popcream}

% ============================================================
% Typographie
% ============================================================
\defaultfontfeatures{Ligatures=TeX}
\setmainfont{PlusJakartaSans-Medium.ttf}[Path=../../../fonts/,BoldFont=PlusJakartaSans-Bold.ttf]
% Symboles, API et emojis absents de la police principale : repli sur DejaVu Sans (caractères actifs)
\newfontface\symfont{DejaVuSans.ttf}[Path=../../../fonts/]
\makeatletter
\newcommand\symactive[1]{\catcode#1=13 \begingroup\lccode`\~=#1 \lowercase{\endgroup\def~}{{\symfont\char#1}}}
\newcommand\symblank[1]{\catcode#1=13 \begingroup\lccode`\~=#1 \lowercase{\endgroup\def~}{}}
\newcommand\symhyph[1]{\catcode#1=13 \begingroup\lccode`\~=#1 \lowercase{\endgroup\def~}{\mbox{-}}}
\newcount\symc
\newcommand\symrange[2]{\symc=#1 \@whilenum\symc<#2 \do{\expandafter\symactive\expandafter{\the\symc}\advance\symc by 1 }}
\newcommand\symblankrange[2]{\symc=#1 \@whilenum\symc<#2 \do{\expandafter\symblank\expandafter{\the\symc}\advance\symc by 1 }}
\makeatother
\symrange{"0250}{"02FF}\symactive{"2713}\symactive{"2714}\symactive{"2717}\symactive{"2718}\symactive{"2610}\symactive{"2611}\symactive{"2612}
\symhyph{"2011}\symhyph{"2E1A}\symblankrange{"1F300}{"1FAFF}\symblank{"FE0F}
% Maths en sans-serif (Fira Math) avec chiffres et lettres droites de Plus
% Jakarta, pour que les formules se fondent dans le texte.
\usepackage[math-style=ISO,bold-style=ISO]{unicode-math}
\setmathfont{FiraMath-Regular.otf}[Path=../../../fonts/]
\setmathfont{PlusJakartaSans-Medium.ttf}[Path=../../../fonts/,range={up/num,up/{Latin,latin},\mathup}]
\usepackage{icomma} % virgule décimale sans espace en mode math
% Caractères mathématiques Unicode absents des polices de texte (Plus Jakarta,
% Archivo Black) : rendus par leur équivalent mathématique, en texte comme en
% formule — sinon ils s'affichent en carré vide (ex. « Arithmétique dans ℤ »).
\usepackage{newunicodechar}
\newunicodechar{ℝ}{\ensuremath{\mathbb{R}}}
\newunicodechar{ℤ}{\ensuremath{\mathbb{Z}}}
\newunicodechar{ℂ}{\ensuremath{\mathbb{C}}}
\newunicodechar{ℕ}{\ensuremath{\mathbb{N}}}
\newunicodechar{ℚ}{\ensuremath{\mathbb{Q}}}
\newunicodechar{𝔻}{\ensuremath{\mathbb{D}}}
\newunicodechar{α}{\ensuremath{\alpha}}
\newunicodechar{β}{\ensuremath{\beta}}
\newunicodechar{γ}{\ensuremath{\gamma}}
\newunicodechar{δ}{\ensuremath{\delta}}
\newunicodechar{ε}{\ensuremath{\varepsilon}}
\newunicodechar{θ}{\ensuremath{\theta}}
\newunicodechar{λ}{\ensuremath{\lambda}}
\newunicodechar{μ}{\ensuremath{\mu}}
\newunicodechar{σ}{\ensuremath{\sigma}}
\newunicodechar{τ}{\ensuremath{\tau}}
\newunicodechar{φ}{\ensuremath{\varphi}}
\newunicodechar{ψ}{\ensuremath{\psi}}
\newunicodechar{ω}{\ensuremath{\omega}}
\newunicodechar{Σ}{\ensuremath{\Sigma}}
% majuscules produites par \MakeUppercase dans les titres (α-aminés -> Α)
\newunicodechar{Α}{\ensuremath{\alpha}}
\newunicodechar{Β}{\ensuremath{\beta}}
\newunicodechar{Γ}{\ensuremath{\Gamma}}
\newunicodechar{Δ}{\ensuremath{\Delta}}
\newunicodechar{Ε}{\ensuremath{\varepsilon}}
\newunicodechar{Θ}{\ensuremath{\Theta}}
\newunicodechar{Λ}{\ensuremath{\Lambda}}
\newunicodechar{Μ}{\ensuremath{\mu}}
\newunicodechar{Τ}{\ensuremath{\tau}}
\newunicodechar{Φ}{\ensuremath{\Phi}}
\newunicodechar{Ψ}{\ensuremath{\Psi}}
\newunicodechar{Ω}{\ensuremath{\Omega}}
\newunicodechar{↦}{\ensuremath{\mapsto}}
\newunicodechar{∈}{\ensuremath{\in}}
\newunicodechar{∉}{\ensuremath{\notin}}
\newunicodechar{∧}{\ensuremath{\wedge}}
\newunicodechar{⇔}{\ensuremath{\Leftrightarrow}}
\newunicodechar{⟺}{\ensuremath{\Longleftrightarrow}}
\newunicodechar{⇒}{\ensuremath{\Rightarrow}}
\newunicodechar{⟹}{\ensuremath{\Longrightarrow}}
\newunicodechar{∘}{\ensuremath{\circ}}
\newunicodechar{∪}{\ensuremath{\cup}}
\newunicodechar{∩}{\ensuremath{\cap}}
\newunicodechar{≡}{\ensuremath{\equiv}}
\newunicodechar{⊂}{\ensuremath{\subset}}
\newunicodechar{⊥}{\ensuremath{\perp}}
\newunicodechar{∃}{\ensuremath{\exists}}
\newunicodechar{∀}{\ensuremath{\forall}}
\newunicodechar{∛}{\ensuremath{\sqrt[3]{\,}}}
\newunicodechar{∜}{\ensuremath{\sqrt[4]{\,}}}
\newunicodechar{⃗}{\ensuremath{{}^{\to}}}
\newunicodechar{ⁿ}{\ifmmode^{n}\else\textsuperscript{\ensuremath{n}}\fi}
\newunicodechar{ˣ}{\ifmmode^{x}\else\textsuperscript{\ensuremath{x}}\fi}
\newunicodechar{ᵇ}{\ifmmode^{b}\else\textsuperscript{\ensuremath{b}}\fi}
\newunicodechar{ʳ}{\ifmmode^{r}\else\textsuperscript{\ensuremath{r}}\fi}
\newunicodechar{ᵘ}{\ifmmode^{u}\else\textsuperscript{\ensuremath{u}}\fi}
\newunicodechar{ʸ}{\ifmmode^{y}\else\textsuperscript{\ensuremath{y}}\fi}
\newunicodechar{ᵃ}{\ifmmode^{a}\else\textsuperscript{\ensuremath{a}}\fi}
\newunicodechar{ᵉ}{\ifmmode^{e}\else\textsuperscript{\ensuremath{e}}\fi}
\newunicodechar{ᵏ}{\ifmmode^{k}\else\textsuperscript{\ensuremath{k}}\fi}
\newunicodechar{ᵖ}{\ifmmode^{p}\else\textsuperscript{\ensuremath{p}}\fi}
\newunicodechar{ᴺ}{\ifmmode^{N}\else\textsuperscript{\ensuremath{N}}\fi}
\newunicodechar{ᶜ}{\ifmmode^{c}\else\textsuperscript{\ensuremath{c}}\fi}
\newunicodechar{ʰ}{\ifmmode^{h}\else\textsuperscript{\ensuremath{h}}\fi}
\newunicodechar{ᵠ}{\ifmmode^{q}\else\textsuperscript{\ensuremath{q}}\fi}
\newunicodechar{ₙ}{\ifmmode_{n}\else\textsubscript{\ensuremath{n}}\fi}
\newunicodechar{ₐ}{\ifmmode_{a}\else\textsubscript{\ensuremath{a}}\fi}
\newunicodechar{ᵢ}{\ifmmode_{i}\else\textsubscript{\ensuremath{i}}\fi}
\newunicodechar{ᵣ}{\ifmmode_{r}\else\textsubscript{\ensuremath{r}}\fi}
\newunicodechar{ₖ}{\ifmmode_{k}\else\textsubscript{\ensuremath{k}}\fi}
\newunicodechar{ᵦ}{\ifmmode_{\beta}\else\textsubscript{\ensuremath{\beta}}\fi}
\newunicodechar{ₓ}{\ifmmode_{x}\else\textsubscript{\ensuremath{x}}\fi}
\newfontfamily\popdisplay{ArchivoBlack-Regular.ttf}[Path=../../../fonts/]
\newfontfamily\poplogo{SpaceGrotesk-Bold.ttf}[Path=../../../fonts/]
\newfontfamily\popsemi{PlusJakartaSans-SemiBold.ttf}[Path=../../../fonts/]
\newfontfamily\popxbold{PlusJakartaSans-ExtraBold.ttf}[Path=../../../fonts/]
\newfontfamily\popxboldls{PlusJakartaSans-ExtraBold.ttf}[Path=../../../fonts/,LetterSpace=3.0]

\setlist[enumerate]{leftmargin=1.7em,itemsep=5pt,topsep=4pt}
\setlist[itemize]{leftmargin=1.7em,itemsep=4pt,topsep=4pt,label={\color{poporange}\textbullet}}
\setlength{\parindent}{0pt}
\setlength{\parskip}{5pt}
\renewcommand{\arraystretch}{1.25}

% pgfplots : style Pop par défaut
\pgfplotsset{
  every axis/.append style={axis line style={popink,line width=1pt},tick style={popink},
    grid style={popline,line width=0.5pt},label style={font=\small},tick label style={font=\footnotesize}},
  popcurve/.style={poporange,line width=1.8pt,no markers,smooth},
}
% Même style disponible dans un \draw TikZ simple (hors pgfplots)
\tikzset{popcurve/.style={poporange,line width=1.8pt}}
\tikzset{popgrid/.style={popline,very thin}}
\colorlet{popcurve}{poporange} % le nom du style est parfois utilisé comme couleur

% ============================================================
% Pill / squiggle
% ============================================================
\newcommand{\poppill}[3]{%
  \tikz[baseline=-0.55ex]\node[draw=popink,line width=1.2pt,rounded corners=2.6pt,%
    fill=#2,inner xsep=6.5pt,inner ysep=3pt]{%
    \popxboldls\fontsize{7}{7}\selectfont\color{#3}\MakeUppercase{#1}};%
}
\newcommand{\popsquiggle}[1]{%
  \tikz[baseline=-0.4ex]{
    \draw[#1,line width=2.6pt,line cap=round]
      plot [smooth] coordinates {(0,0.09) (0.34,0.01) (0.68,0.21) (1.02,0.01) (1.36,0.10)};
  }%
}
% Mot-clé surligné (marqueur orange sous le texte)
\newcommand{\cle}[1]{{\popxbold\color{popdark}#1}}

% ============================================================
% En-tête / pied
% ============================================================
\pagestyle{fancy}
\fancyhf{}
\renewcommand{\headrulewidth}{0pt}
\renewcommand{\footrulewidth}{0pt}
\lhead{\raisebox{-0.15ex}{\poplogo\fontsize{13}{13}\selectfont\color{popink}OnBuch\color{poporange}+}}
\rhead{\poppill{\DOCMATIERE\ \textbullet\ \DOCNIVEAU\ \textbullet\ Leçon \DOCLECON}{white}{popink}}
\lfoot{\popxbold\fontsize{7.6}{7.6}\selectfont\color{popmuted}\MakeUppercase{Cours signé OnBuch+}\ \textbullet\ {\popsemi\DOCTITRE}}
\rfoot{\tikz[baseline=-0.6ex]\node[rounded corners=8.5pt,fill=popink,minimum height=17pt,minimum width=17pt,inner xsep=5pt,inner sep=0]{\popxbold\fontsize{8}{8}\selectfont\color{popcream}\thepage\,/\,\pageref*{LastPage}};}

% ============================================================
% Titres
% ============================================================
\newcommand{\chaptitre}[2]{%
  \begin{tcolorbox}[enhanced,colback=poporange,colframe=popink,boxrule=2.6pt,arc=11pt,
      drop shadow=popink,left=15pt,right=15pt,top=12pt,bottom=11pt,before skip=3pt,after skip=18pt]
    \poppill{#2}{popink}{popcream}\par\vspace{9pt}
    {\raggedright\hyphenpenalty=10000\exhyphenpenalty=10000\popdisplay\fontsize{22}{24}\selectfont\color{popcream}\MakeUppercase{#1}\par}
  \end{tcolorbox}
}
% Section numérotée : pastille orange avec le numéro + titre Archivo
\newcounter{popsec}
\newcommand{\coursec}[1]{%
  \stepcounter{popsec}\setcounter{popsub}{0}%
  \par\vspace{16pt}\noindent
  \begin{minipage}{\linewidth}
  \tikz[baseline=-0.7ex]\node[draw=popink,line width=1.6pt,fill=poporange,rounded corners=4pt,minimum size=22pt,inner sep=0]{\popdisplay\fontsize{12}{12}\selectfont\color{popcream}\thepopsec};%
  \hspace{8pt}{\popdisplay\fontsize{15}{17}\selectfont\color{popink}\MakeUppercase{#1}}\par
  \vspace{4pt}\noindent{\color{poporange}\rule{36pt}{3.6pt}}
  \end{minipage}\par\nopagebreak\vspace{8pt}
}
\newcounter{popsub}
\newcommand{\courssub}[1]{%
  \stepcounter{popsub}%
  \par\vspace{9pt}\noindent{\popxbold\fontsize{11.5}{13}\selectfont\color{popink}{\color{poporange}\thepopsec.\thepopsub}\hspace{6pt}#1}\par\nopagebreak\vspace{4pt}
}

% ============================================================
% Boîtes pédagogiques — même recette : fond blanc, bordure épaisse colorée,
% ombre dure encre, badge plein. Titre optionnel [#1].
% ============================================================
\tcbset{popbase/.style={breakable,enhanced,colback=white,boxrule=2.3pt,arc=9pt,
  shadow={3pt}{-3pt}{0pt}{popink},left=14pt,right=14pt,top=11pt,bottom=11pt,before skip=11pt,after skip=15pt,
  fontupper=\popsemi\fontsize{10.6}{14.5}\selectfont\color{popink},
  fontlower=\popsemi\fontsize{10.6}{14.5}\selectfont\color{popink}}}
% Coche dessinée (le glyphe ✓ n'existe pas dans les polices du document)
\newcommand{\popcheck}[1]{\tikz[baseline=-0.5ex]\draw[#1,line width=1.6pt,line cap=round,line join=round](-0.13,0.0)--(-0.04,-0.09)--(0.13,0.1);}
\providecommand{\coche}{\popcheck{popgreen}}
\providecommand{\resultat}[1]{\par\smallskip\noindent\fcolorbox{popgreen}{popgreenL}{\parbox{\dimexpr\linewidth-2\fboxsep-2\fboxrule\relax}{\textbf{Résultat.} #1}}\par\smallskip}
\newcommand{\popboxhead}[3]{\poppill{#1}{#2}{white}\ifx\relax#3\relax\else\hspace{6pt}{\popxbold\fontsize{10.6}{12}\selectfont\color{popink}#3}\fi\par\vspace{7pt}}

\newtcolorbox{aretenir}[1][]{popbase,colframe=popgreen,before upper={\popboxhead{À retenir}{popgreen}{#1}}}
% Texte en allemand (document de lecture, dialogue, extrait) : cadre
% d'étude. Un titre optionnel (source, auteur) s'affiche dans l'en-tête.
\newtcolorbox{textebox}[1][]{popbase,colframe=popblue,before upper={\popboxhead{Text}{popblue}{#1}\begin{otherlanguage}{ngerman}},after upper={\end{otherlanguage}}}
% Marques ✓ / ✗ (forme correcte / fautive) souvent utilisées par les modèles
\providecommand{\cross}{\textcolor{poppink}{\ensuremath{\times}}}
\providecommand{\xmark}{\cross}\providecommand{\crossmark}{\cross}\providecommand{\cmark}{\ensuremath{\checkmark}}
% Ligne de réponse / justification à compléter (inventée par les modèles)
\providecommand{\justification}[1]{\par\noindent\textit{Justification~:} \dotfill\par}
% \textipa (package tipa non chargé) : la police ne couvre pas l'API -> simple passage
\providecommand{\textipa}[1]{#1}
% Soulignement (ulem non chargé) : \uline, \ul
\providecommand{\uline}[1]{\underline{#1}}\providecommand{\ul}[1]{\underline{#1}}
% \glossaire{\item ...} inventée par les modèles : liste de glossaire
\providecommand{\glossaire}[1]{\begin{itemize}#1\end{itemize}}
% \alert (beamer) inventée par les modèles : texte mis en évidence
\providecommand{\alert}[1]{\textbf{\textcolor{poppink}{#1}}}
% variantes ASCII d'ordinaux écrites en macro
\providecommand{\iere}{\textsuperscript{re}}\providecommand{\ieme}{\textsuperscript{e}}\providecommand{\ere}{\textsuperscript{re}}\providecommand{\eme}{\textsuperscript{e}}\providecommand{\er}{\textsuperscript{er}}\providecommand{\ie}{\textsuperscript{e}}
% Ordinaux français écrits en macro par les modèles : 1\ière, 3\ième
\newcommand{\ière}{\textsuperscript{re}}\newcommand{\ième}{\textsuperscript{e}}
% \exemple (inventée par les modèles) : étiquette « Exemple : »
\providecommand{\exemple}{\textbf{Exemple~:}\ }
% \square utilisable aussi hors mode math (cases à cocher)
\AtBeginDocument{\let\squareMath\square\renewcommand{\square}{\ensuremath{\squareMath}}}
% Mot ou phrase en allemand à l'intérieur d'un paragraphe français
\newcommand{\all}[1]{\ifmmode\text{\textit{#1}}\else\begingroup\selectlanguage{ngerman}\itshape #1\endgroup\fi}
\let\esp\all % alias toléré
\newtcolorbox{exemplebox}[1][]{popbase,colframe=poporange,before upper={\popboxhead{Exemple}{poporange}{#1}}}
\newtcolorbox{definition}[1][]{popbase,colframe=popblue,before upper={\popboxhead{Définition}{popblue}{#1}}}
\newtcolorbox{propriete}[1][]{popbase,colframe=poppurple,before upper={\popboxhead{Propriété}{poppurple}{#1}}}
\newtcolorbox{methode}[1][]{popbase,colframe=popgold,before upper={\popboxhead{Méthode}{popgold}{#1}}}
\newtcolorbox{attention}[1][]{popbase,colframe=poppink,before upper={\popboxhead{Attention}{poppink}{#1}}}
\newtcolorbox{experience}[1][]{popbase,colframe=popblue,colback=popblueL,before upper={\popboxhead{Expérience}{popblue}{#1}}}
% « Le savais-tu ? » : carte violette pleine (contexte, culture, Cameroun)
\newtcolorbox{savaistu}[1][]{popbase,colback=poppurple,colframe=popink,
  fontupper=\popsemi\fontsize{10.4}{14.2}\selectfont\color{white},
  before upper={\poppill{Le savais-tu\,?}{white}{poppurple}\ifx\relax#1\relax\else\hspace{6pt}{\popxbold\color{white}#1}\fi\par\vspace{7pt}}}
% Exercice résolu : énoncé en haut, \tcblower puis la solution sur fond crème
\newcounter{popexo}\newcounter{popexr}
\newtcolorbox{exoresolu}[1][]{popbase,colframe=poporange,colbacklower=popcream,
  segmentation style={popink,line width=1.2pt,dashed},
  before upper={\stepcounter{popexr}\popboxhead{Exercice résolu \thepopexr}{poporange}{#1}},
  before lower={\poppill{Solution}{popink}{popcream}\par\vspace{6pt}}}
% Exercice (non résolu en place) — la correction est dans la partie Corrigés
\newtcolorbox{exercice}[1][]{popbase,colframe=popink,
  before upper={\stepcounter{popexo}\popboxhead{Exercice \thepopexo}{popink}{#1}}}
% Corrigé
\newtcolorbox{corrige}[1][]{popbase,colframe=popgreen,colback=popgreenL,
  before upper={\popboxhead{Corrigé}{popgreen}{#1}}}
% Objectifs / prérequis / bilan
\newtcolorbox{objectifs}{popbase,colframe=popink,colback=poporangeL,
  before upper={\popboxhead{Objectifs de la leçon}{popink}{}}}
\newtcolorbox{prerequis}{popbase,colframe=popink,colback=popblueL,
  before upper={\popboxhead{Prérequis}{popblue}{}}}
\newtcolorbox{bilan}{popbase,colframe=popink,colback=white,boxrule=2.8pt,
  before upper={\popboxhead{Fiche bilan}{poporange}{L'essentiel à connaître pour l'examen}}}
% Figure encadrée avec légende
\newtcolorbox{popfigure}[1][]{enhanced,breakable=false,colback=white,colframe=popline,boxrule=1.4pt,arc=8pt,
  left=10pt,right=10pt,top=10pt,bottom=8pt,before skip=10pt,after skip=12pt,halign=center,
  lower separated=false,halign lower=center,
  fontlower=\popsemi\fontsize{9}{11.5}\selectfont\color{popmuted},
  #1}
\newcommand{\legende}[1]{\par\vspace{6pt}{\popsemi\fontsize{9}{11.5}\selectfont\color{popmuted}#1}}

% ============================================================
% Signature OnBuch+
% ============================================================
\newcommand{\poplogoinline}[1]{{\poplogo\fontsize{#1}{#1}\selectfont\color{popink}OnBuch\color{poporange}+}}
\newcommand{\signatureonbuch}{%
  \begin{tcolorbox}[enhanced,colback=popink,colframe=popink,boxrule=2.2pt,arc=10pt,
      drop shadow=poporange,left=14pt,right=14pt,top=10pt,bottom=10pt,before skip=0pt,after skip=0pt]
    \tikz[baseline=-0.7ex]\node[circle,fill=popgreen,draw=popcream,line width=1.3pt,minimum size=22pt,inner sep=0]{\popcheck{white}};%
    \hspace{9pt}\begin{minipage}[c]{0.62\linewidth}
      {\popxbold\fontsize{12}{13}\selectfont\color{popcream}Signé\ \textbullet\ {\poplogo OnBuch\color{poporange}+}}\par\vspace{2pt}
      {\raggedright\popsemi\fontsize{8}{10.5}\selectfont\color{popline}\MakeUppercase{Équipe pédagogique OnBuch+}\\\MakeUppercase{Conforme au programme officiel MINESEC}\par}
    \end{minipage}\hfill
    {\poplogo\fontsize{22}{22}\selectfont\color{popcream}OnBuch\color{poporange}+}
  \end{tcolorbox}%
}

% ============================================================
% Page de garde
% #1 titre · #2 matière · #3 niveau · #4 module · #5 n° leçon · #6 durée ·
% #7 description / objectifs (texte ou itemize)
% ============================================================
\newcommand{\pagegarde}[7]{%
  \thispagestyle{empty}
  \newgeometry{a4paper,margin=1.7cm,top=1.6cm,bottom=1.4cm}%
  \begin{tikzpicture}[remember picture,overlay]
    \fill[poporange!18] ([xshift=-3.2cm,yshift=-3.6cm]current page.north east) circle (4.6cm);
    \fill[poppurple!14] ([xshift=2.4cm,yshift=7.6cm]current page.south west) circle (3.8cm);
  \end{tikzpicture}%
  {\poplogo\fontsize{30}{30}\selectfont\color{popink}OnBuch\color{poporange}+}\hfill
  \raisebox{0.6ex}{\poppill{Cours signé \textbullet\ Premium}{popink}{popcream}}\par
  \vspace{1pt}\popsquiggle{poppurple}\par
  \vspace{8pt}
  \begin{tcolorbox}[enhanced,colback=poporange,colframe=popink,boxrule=2.8pt,arc=13pt,
      drop shadow=popink,left=18pt,right=18pt,top=12pt,bottom=13pt,after skip=10pt]
    \poppill{#2 \textbullet\ #3}{popink}{popcream}\hspace{5pt}\poppill{#4}{white}{popink}\par\vspace{8pt}
    {\popdisplay\fontsize{14}{14}\selectfont\color{popink}LEÇON #5}\par\vspace{4pt}
    {\raggedright\hyphenpenalty=10000\exhyphenpenalty=10000\popdisplay\fontsize{25}{27}\selectfont\color{popcream}\MakeUppercase{#1}\par}
    \vspace{8pt}
    {\popxbold\fontsize{9.5}{9.5}\selectfont\color{popink}\MakeUppercase{Durée conseillée : #6}}
  \end{tcolorbox}
  \begin{tcolorbox}[enhanced,colback=white,colframe=popink,boxrule=2.2pt,arc=10pt,
      drop shadow=popink,left=16pt,right=16pt,top=10pt,bottom=10pt,after skip=8pt]
    \poppill{Dans cette leçon}{poppurple}{white}\par\vspace{5pt}
    \popsemi\fontsize{9.3}{12.6}\selectfont\color{popink}#7
  \end{tcolorbox}
  \noindent\begin{minipage}[t]{0.487\linewidth}
    \begin{tcolorbox}[enhanced,colback=popgreen,colframe=popink,boxrule=2.2pt,arc=10pt,
        drop shadow=popink,left=11pt,right=11pt,top=10pt,bottom=10pt,height=3.1cm,valign=center]
      \tcbox[colback=white,colframe=popink,boxrule=1.3pt,arc=5pt,boxsep=0pt,nobeforeafter,
        left=3pt,right=3pt,top=3pt,bottom=3pt,valign=center]{\qrcode[height=1.9cm]{https://play.google.com/store/apps/details?id=com.luvvix.onbuch&hl=fr}}%
      \hspace{8pt}\begin{minipage}[c]{0.5\linewidth}
        {\popdisplay\fontsize{11}{12.5}\selectfont\color{white}\MakeUppercase{Télécharge OnBuch}}\par\vspace{4pt}
        {\raggedright\popsemi\fontsize{8.4}{11}\selectfont\color{white}Gratuit \textbullet\ Google Play\\Tuteur IA, annales, concours, résultats\par}
      \end{minipage}
    \end{tcolorbox}
  \end{minipage}\hfill
  \begin{minipage}[t]{0.487\linewidth}
    \begin{tcolorbox}[enhanced,colback=poppurple,colframe=popink,boxrule=2.2pt,arc=10pt,
        drop shadow=popink,left=11pt,right=11pt,top=10pt,bottom=10pt,height=3.1cm,valign=center]
      \tikz[baseline=-0.7ex]\node[circle,fill=white,draw=popink,line width=1.3pt,minimum size=1.6cm,inner sep=0]{\popdisplay\fontsize{24}{24}\selectfont\color{poppurple}+};%
      \hspace{8pt}\begin{minipage}[c]{0.6\linewidth}
        {\popdisplay\fontsize{11}{12.5}\selectfont\color{white}\MakeUppercase{Passe à OnBuch+}}\par\vspace{4pt}
        {\raggedright\popsemi\fontsize{8.4}{11}\selectfont\color{white}Tous les cours signés, Tuteur IA illimité, fiches PDF — onglet OnBuch+ de l'app\par}
      \end{minipage}
    \end{tcolorbox}
  \end{minipage}
  \vspace{8pt}
  \popcontenu
  \vfill
  \signatureonbuch
  \restoregeometry
  \newpage
}

% Rangée « Ce cours contient » (page de garde)
\newcommand{\poptile}[3]{% #1 couleur #2 gros texte #3 légende
  \begin{tcolorbox}[enhanced,colback=#1,colframe=popink,boxrule=2pt,arc=9pt,shadow={2.5pt}{-2.5pt}{0pt}{popink},
      left=6pt,right=6pt,top=9pt,bottom=9pt,halign=center,valign=center,height=1.75cm,width=\linewidth,nobeforeafter]
    {\popdisplay\fontsize{12.5}{13}\selectfont\color{white}\MakeUppercase{#2}}\par\vspace{3pt}
    {\centering\popsemi\fontsize{7.8}{9.5}\selectfont\color{white}#3\par}
  \end{tcolorbox}}
\newcommand{\popcontenu}{%
  \par\noindent{\popxboldls\fontsize{7.5}{7.5}\selectfont\color{popmuted}CE COURS SIGNÉ CONTIENT}\par\vspace{7pt}
  \noindent\begin{minipage}[t]{0.235\linewidth}\poptile{popblue}{Cours}{complet, illustré, pas à pas}\end{minipage}\hfill
  \begin{minipage}[t]{0.235\linewidth}\poptile{poporange}{Méthodes}{et exercices résolus}\end{minipage}\hfill
  \begin{minipage}[t]{0.235\linewidth}\poptile{poppink}{Exercices}{type examen + corrigés}\end{minipage}\hfill
  \begin{minipage}[t]{0.235\linewidth}\poptile{popgreen}{Fiche bilan}{l'essentiel à retenir}\end{minipage}\par}

% Sommaire
\newtcolorbox{sommaire}{enhanced,colback=white,colframe=popink,boxrule=2.2pt,arc=10pt,
  drop shadow=popink,left=18pt,right=18pt,top=13pt,bottom=13pt,before skip=6pt,after skip=20pt,
  before upper={\poppill{Sommaire}{poppurple}{white}\par\vspace{9pt}\popsemi\fontsize{10.6}{14.5}\selectfont\color{popink}}}

% Signature de fin de cours — poussée tout en bas de la dernière page
\newcommand{\findecours}{%
  \par\vfill\nopagebreak
  \begin{center}\popsquiggle{poporange}\end{center}\vspace{4pt}
  \signatureonbuch
}
\AtBeginDocument{\def\llbracket{\mathopen{[\mkern-3.5mu[}}\def\rrbracket{\mathclose{]\mkern-3.5mu]}}} % intervalles d'entiers (glyphe ⟦ absent de la police)
% Glyphes absents de Fira Math : redéfinis à partir de glyphes disponibles
\AtBeginDocument{%
  \def\setminus{\mathbin{\backslash}}\let\smallsetminus\setminus
  \def\vdots{\mathord{\vbox{\baselineskip3.2pt\lineskiplimit0pt\kern4pt\hbox{.}\hbox{.}\hbox{.}}}}%
  \def\ddots{\mathinner{\mkern1mu\raise6.5pt\hbox{.}\mkern2mu\raise3.6pt\hbox{.}\mkern2mu\raise0.7pt\hbox{.}\mkern1mu}}%
  \def\triangle{\mathord{\scalebox{0.9}{$\symup{\Delta}$}}}%
  \def\nabla{\mathord{\raisebox{\depth}{\scalebox{1}[-1]{$\symup{\Delta}$}}}}%
  \def\checkmark{\popcheck{popdark}}%
  \def\star{\mathbin{\ast}}\def\bigstar{\mathord{\ast}}%
  \def\diamond{\mathbin{\scalebox{0.6}{$\Diamond$}}}%
  \def\bigcup{\mathop{\mathchoice{\raisebox{-0.3em}{\scalebox{1.7}{$\cup$}}}{\scalebox{1.25}{$\cup$}}{\cup}{\cup}}\displaylimits}%
  \def\bigcap{\mathop{\mathchoice{\raisebox{-0.3em}{\scalebox{1.7}{$\cap$}}}{\scalebox{1.25}{$\cap$}}{\cap}{\cap}}\displaylimits}%
  \def\lesseqqgtr{\mathrel{\lessgtr}}\def\bigoplus{\mathop{\oplus}\displaylimits}%
}
\providecommand{\ssi}{si et seulement si} % abréviation fréquente
\AtBeginDocument{\providecommand{\wideparen}{\overparen}} % arc de cercle

%% FIGFIX-BEGIN v4 — correctifs globaux des figures (docs/onbuch-generation/tools/figfix)
\usepackage{adjustbox}\usepackage{etoolbox}
% toute figure TikZ/pgfplots plus large que la ligne est réduite à la largeur disponible
\BeforeBeginEnvironment{tikzpicture}{\begin{adjustbox}{max width=\linewidth}}
\AfterEndEnvironment{tikzpicture}{\end{adjustbox}}
% légendes pgfplots lisibles par-dessus les courbes
\pgfplotsset{every axis/.append style={legend style={fill=white,fill opacity=0.92,text opacity=1,draw=black!15,inner sep=3pt}}}
% étiquettes d'arêtes (syntaxe edge["..."]) avec fond blanc
\tikzset{every edge quotes/.append style={fill=white,inner sep=1.2pt}}
%% FIGFIX-END
\begin{document}

\pagegarde{La cybercriminalité 1 : les fake news}{Allemand}{Tle A}{Chapitre 4 : Média et communication}{AL17}{2 h}{Cette leçon de type « texte » propose la lecture et le commentaire d'un article allemand sur les fake news et la cybercriminalité. Elle permet d'acquérir le vocabulaire de l'information et de la manipulation, et de maîtriser le discours indirect au subjonctif I ainsi que la subordonnée en « dass ».
\vspace{4pt}
\begin{multicols}{2}\small
\begin{itemize}
\item À la fin de cette leçon, je sais comprendre globalement puis en détail un texte allemand sur les fake news et la cybercriminalité.
\item Je sais employer correctement le vocabulaire thématique (die Falschmeldung, die Lüge, die Wahrheit, verbreiten, glauben, der Beweis) avec articles et pluriels.
\item Je sais former et conjuguer le subjonctif I à tous les temps utiles (présent, passé, futur).
\item Je sais rapporter une parole au discours indirect en allemand, avec ou sans « dass ».
\item Je sais construire et traduire une subordonnée complétive en « dass ».
\item Je sais comparer le discours indirect allemand et français et éviter les pièges de traduction.
\end{itemize}\end{multicols}}

\begin{sommaire}
\begin{multicols}{2}
\begin{enumerate}
\item Texte support : Fake News im Netz
\item Compréhension du texte
\item Lexique thématique : Wahrheit und Lüge
\item Grammaire 1 : le discours indirect (subjonctif I)
\item Grammaire 2 : la subordonnée en « dass »
\item Méthode du commentaire et commentaire modèle
\item Activité d'intégration
\item Méthodes et exercices résolus
\item Exercices
\item Corrigés des exercices
\item Fiche bilan
\end{enumerate}
\end{multicols}
\end{sommaire}

\begin{objectifs}
\begin{itemize}
\item À la fin de cette leçon, je sais comprendre globalement puis en détail un texte allemand sur les fake news et la cybercriminalité.
\item Je sais employer correctement le vocabulaire thématique (die Falschmeldung, die Lüge, die Wahrheit, verbreiten, glauben, der Beweis) avec articles et pluriels.
\item Je sais former et conjuguer le subjonctif I à tous les temps utiles (présent, passé, futur).
\item Je sais rapporter une parole au discours indirect en allemand, avec ou sans « dass ».
\item Je sais construire et traduire une subordonnée complétive en « dass ».
\item Je sais comparer le discours indirect allemand et français et éviter les pièges de traduction.
\item Je sais répondre à des questions de compréhension et rédiger un résumé ou un court paragraphe argumenté sur les fake news.
\end{itemize}
\end{objectifs}

\begin{prerequis}
\begin{itemize}
\item Les verbes de parole et d'opinion (sagen, meinen, glauben, behaupten, berichten) vus en Première
\item La conjugaison du présent de l'indicatif et du parfait (Perfekt)
\item La place du verbe dans la subordonnée (verbe conjugué en fin de proposition)
\item Le discours indirect au présent en français (il dit que... / il a dit que...)
\item Le vocabulaire des médias vu en Première (die Zeitung, der Fernseher, das Internet)
\end{itemize}
\end{prerequis}

\newpage

\coursec{Texte support : Fake News im Netz}

\courssub{Situation d'accroche et problématique}

Pendant les élections ou une épidémie, des messages WhatsApp circulent au Cameroun : „Le gouvernement a décidé…“, „Un remède miracle existe…“. Beaucoup les croient et les partagent sans vérifier. À Yaoundé comme à Douala, un simple transfert peut semer la panique dans un quartier entier. En Allemagne aussi, les \all{Fake News} se répandent sur Facebook, TikTok et Twitter (désormais X), et le mot \all{Lügenpresse} („presse mensongère“) est même devenu un débat national.

Deux questions vont guider notre lecture : comment \cle{reconnaître} une fausse information ? Et comment \cle{rapporter} ce que quelqu'un a dit sans le trahir ? C'est exactement ce que permet le discours indirect allemand, que nous étudierons à partir de ce texte de presse.

\courssub{Lecture globale du texte}

\begin{textebox}[Fake News – die neue Gefahr im Internet]
Immer mehr Menschen informieren sich nur noch über soziale Medien. Ein Medienexperte erklärte gestern in Berlin, viele Jugendliche könnten wahre und falsche Nachrichten kaum unterscheiden. Falschmeldungen verbreiteten sich viel schneller als die Wahrheit, sagte er. Oft gebe es keinen Absender und keine Quelle. Der Journalist Thomas M. behauptete, manche Fake News seien sogar von Regierungen bezahlt worden. Er riet, jede Nachricht zu überprüfen, bevor man sie teile. Denn wer Lügen verbreitet, macht sich strafbar.

Nach einer Studie lesen fast siebzig Prozent der jungen Deutschen ihre Nachrichten auf dem Handy, meistens bei Instagram, TikTok oder WhatsApp. Dort erscheinen jeden Tag Tausende von Meldungen – manche wahr, manche erfunden. Eine Falschmeldung über einen angeblichen Skandal im Gesundheitsministerium wurde zum Beispiel in nur zwei Stunden mehr als hunderttausend Mal geteilt, bevor die Behörden sie dementieren konnten.

Warum glauben so viele Menschen an Fake News? Der Experte nannte drei Gründe: Erstens seien solche Meldungen oft emotional und aufregend, zweitens würden sie von Freunden weitergeleitet, denen man vertraue, und drittens fehle vielen Nutzern die Zeit, die Quelle zu prüfen. Er warnte auch vor sogenannten „Deepfakes“, also Videos, in denen Politiker Dinge zu sagen scheinen, die sie nie gesagt haben.

Die Politik reagiert inzwischen. Das Netzwerkdurchsetzungsgesetz verpflichtet die großen Plattformen, illegale Inhalte schnell zu löschen. Schulen bieten Kurse in Medienkompetenz an, damit Schüler lernen, Quellen zu vergleichen und Beweise zu suchen. „Wer eine Nachricht teilt, ohne sie zu prüfen, wird selbst zum Problem“, sagte der Experte am Ende des Interviews. Sein Rat an alle Nutzer: Erst denken, dann klicken.
\end{textebox}

\begin{definition}[Glossaire des mots difficiles du texte]
\begin{itemize}
\item \all{die Falschmeldung, -en} : la fausse information, la « fake news »
\item \all{die Lüge, -n} : le mensonge
\item \all{die Wahrheit, -en} : la vérité
\item \all{verbreiten (verbreitete, hat verbreitet)} : répandre, propager (transitif) ; \all{sich verbreiten} : se répandre, se propager (réfléchi)
\item \all{glauben (glaubte, hat geglaubt)} : croire
\item \all{der Beweis, -e} : la preuve
\item \all{hetzen (hetzte, hat gehetzt)} : inciter à la haine, harceler ; \all{die Hetze} : la campagne de haine, le harcèlement
\item \all{der Absender, -} : l'expéditeur (d'un message)
\item \all{überprüfen (überprüfte, hat überprüft)} : vérifier, contrôler
\item \all{die Quelle, -n} : la source
\item \all{täuschen (täuschte, hat getäuscht)} : tromper ; \all{sich täuschen} : se tromper
\item \all{dementieren (dementierte, hat dementiert)} : démentir
\item \all{erfunden} : inventé (participe II de \all{erfinden})
\item \all{strafbar} : punissable par la loi ; \all{sich strafbar machen} : commettre un délit
\item \all{die Medienkompetenz} : l'éducation aux médias, la compétence médiatique
\item \all{der Skandal, -e} : le scandale
\item \all{angeblich} : soi-disant, prétendu(e)
\item \all{vertrauen (vertraute, hat vertraut) + Dat.} : faire confiance à qqn
\end{itemize}
\end{definition}

\begin{methode}[Lire un texte journalistique]
\begin{enumerate}
\item Lire le \cle{titre} et la \cle{source} : de quoi parle-t-on ? Qui écrit, pour qui ?
\item Faire une \cle{première lecture rapide} pour saisir l'idée générale, sans s'arrêter sur les mots inconnus.
\item Faire une \cle{deuxième lecture} avec le glossaire, paragraphe par paragraphe.
\item \cle{Souligner} les mots-clés (thème, chiffres, noms propres) et les \cle{verbes de parole} (\all{sagen, erklären, behaupten, raten, warnen, nennen}) : ils annoncent le discours indirect.
\end{enumerate}
\end{methode}

\courssub{Repérage : thème, type de texte, source, destinataire}

\begin{center}
\begin{tabularx}{\linewidth}{|X|X|}\hline
\rowcolor{popblueL} \textbf{Question} & \textbf{Réponse} \\ \hline
Thème (\all{das Thema}) & Les fake news et leur danger sur Internet \\ \hline
Type de texte (\all{die Textsorte}) & Article de presse (registre journalistique, titre, faits, citations d'experts) \\ \hline
Source (\all{die Quelle}) & Un journal allemand ; un média d'information générale \\ \hline
Destinataire (\all{der Leser}) & Le grand public, en particulier les jeunes utilisateurs des réseaux sociaux \\ \hline
Intention (\all{die Absicht}) & Informer, mettre en garde et conseiller (\all{Erst denken, dann klicken}) \\ \hline
\end{tabularx}
\end{center}

\courssub{Lecture détaillée paragraphe par paragraphe}

\textbf{Paragraphe 1 : le constat.} L'expert berlinois constate que les jeunes ne distinguent presque plus le vrai du faux. Questions de repérage : \all{Wo sprach der Experte? Was können viele Jugendliche kaum unterscheiden?} (Où l'expert a-t-il parlé ? Que les jeunes arrivent-ils à peine à distinguer ?)

\textbf{Paragraphe 2 : l'ampleur du phénomène.} Près de 70 \% des jeunes Allemands s'informent sur le téléphone ; une fausse information a été partagée plus de cent mille fois en deux heures. Questions : \all{Wo lesen junge Deutsche ihre Nachrichten? Wie schnell wurde die Falschmeldung geteilt?}

\textbf{Paragraphe 3 : les causes.} Trois raisons expliquent le succès des mensonges : l'émotion, la confiance dans les amis, le manque de temps. Questions : \all{Welche drei Gründe nennt der Experte? Was sind „Deepfakes“?}

\textbf{Paragraphe 4 : les réponses.} La loi oblige les plateformes à supprimer les contenus illégaux ; les écoles enseignent l'éducation aux médias. Question : \all{Was lernen die Schüler in den Kursen?}

\begin{popfigure}
\begin{tikzpicture}[node distance=0.55cm, every node/.style={font=\small}]
\node[draw=popblue, fill=popblueL, rounded corners, align=center, minimum width=2.4cm] (a) {\all{Absender}\\ expéditeur};
\node[draw=poporange, fill=poporangeL, rounded corners, align=center, minimum width=2.4cm, right=of a] (b) {\all{soziale Medien}\\ réseaux sociaux};
\node[draw=poppurple, fill=poppurpleL, rounded corners, align=center, minimum width=2.4cm, right=of b] (c) {\all{Teilen}\\ partage};
\node[draw=popgreen, fill=popgreenL, rounded corners, align=center, minimum width=2.4cm, right=of c] (d) {\all{Verbreitung}\\ propagation};
\node[draw=poppink, fill=poppinkL, rounded corners, align=center, minimum width=2.4cm, right=of d] (e) {\all{Folgen}\\ conséquences};
\draw[-{Stealth}, thick, popdark] (a) -- (b);
\draw[-{Stealth}, thick, popdark] (b) -- (c);
\draw[-{Stealth}, thick, popdark] (c) -- (d);
\draw[-{Stealth}, thick, popdark] (d) -- (e);
\end{tikzpicture}
\legende{La chaîne de propagation d'une fake news : \all{Absender} $\rightarrow$ \all{soziale Medien} $\rightarrow$ \all{Teilen} $\rightarrow$ \all{Verbreitung} $\rightarrow$ \all{Folgen}}
\end{popfigure}

\courssub{Grammaire 1 : Le discours indirect (Konjunktiv I)}

Le journaliste ne cite pas l'expert mot à mot : il \cle{rapporte} ses paroles. En allemand, cela se fait avec le \cle{Konjunktiv I} (Subjonctif I). Relevons les exemples du texte (verbe en \textbf{gras}) :

\begin{center}
\begin{tabularx}{\linewidth}{|X|X|}\hline
\rowcolor{popblueL} \textbf{Discours indirect (texte)} & \textbf{Discours direct (paroles réelles)} \\ \hline
\all{viele Jugendliche \textbf{könnten} wahre und falsche Nachrichten kaum unterscheiden} & \all{Viele Jugendliche \textbf{können} wahre und falsche Nachrichten kaum unterscheiden.} \\ \hline
\all{Falschmeldungen \textbf{verbreiteten} sich viel schneller als die Wahrheit} & \all{Falschmeldungen \textbf{verbreiten} sich viel schneller als die Wahrheit.} \\ \hline
\all{Oft \textbf{gebe} es keinen Absender und keine Quelle} & \all{Oft \textbf{gibt} es keinen Absender und keine Quelle.} \\ \hline
\all{manche Fake News \textbf{seien} sogar von Regierungen bezahlt worden} & \all{Manche Fake News \textbf{sind} sogar von Regierungen bezahlt worden.} \\ \hline
\all{drittens \textbf{fehle} vielen Nutzern die Zeit} & \all{Drittens \textbf{fehlt} vielen Nutzern die Zeit.} \\ \hline
\all{man \textbf{teile} (avant \all{teile})} & \all{man \textbf{teilt}} \\ \hline
\end{tabularx}
\end{center}

\begin{propriete}[Formation du Konjunktiv I (Présent)]
Le Konjunktiv I se forme à partir du \cle{radical de l'infinitif présent} + \cle{terminaisons spécifiques}. Il n'a qu'un temps : le présent (valeur de simultanéité ou d'antériorité selon le verbe introducteur).

\begin{center}
\begin{tabular}{|l|c|c|c|c|c|c|}\hline
\rowcolor{popblueL} \textbf{Personne} & \textbf{sein} & \textbf{haben} & \textbf{werden} & \textbf{können} & \textbf{machen (régulier)} & \textbf{Terminaison} \\ \hline
ich & sei & habe & werde & könne & mache & \textbf{-e} \\ \hline
du & seiest & habest & werdest & könnest & machest & \textbf{-est} \\ \hline
er/sie/es & \textbf{sei} & \textbf{habe} & \textbf{werde} & \textbf{könne} & \textbf{mache} & \textbf{-e} \\ \hline
wir & seien & haben & werden & können & machen & \textbf{-en} \\ \hline
ihr & seiet & habet & werdet & könnet & machet & \textbf{-et} \\ \hline
sie/Sie & seien & haben & werden & können & machen & \textbf{-en} \\ \hline
\end{tabular}
\end{center}

\end{propriete}

\begin{aretenir}[Règles clés et exceptions]
\begin{itemize}
\item \cle{3e personne du singulier (er/sie/es)} : c'est la forme la plus utilisée au discours indirect. Elle se termine toujours par \textbf{-e} (pas de \textbf{-t} comme à l'indicatif : \all{er macht} $\rightarrow$ \all{er mache}).
\item \cle{Formes identiques à l'indicatif} : \all{ich, wir, sie (pl.)} ont les mêmes formes qu'à l'indicatif présent. \textbf{On les remplace par le Konjunktiv II (forme de remplacement / \textit{Ersatzform})}. Ex. : \all{wir sind} (Indik.) $\rightarrow$ \all{wir \textbf{wären}} (Konj. II) ; \all{ich mache} $\rightarrow$ \all{ich \textbf{machte}}.
\item \cle{Verbes modaux} : pas d'Umlaut au Konjunktiv I (\all{könne}, \all{müsse}, \all{solle}, \all{wolle}, \all{dürfe}, \all{möge}).
\item \cle{Valeur} : Le Konjunktiv I marque la \cle{distance} : le locuteur rapporte sans s'engager sur la vérité du propos.
\end{itemize}
\end{aretenir}

\begin{exoresolu}[Repérage et analyse dans le texte]
Énoncé : 1) Relevez dans le texte trois verbes introducteurs de parole et indiquez qui parle. 2) Transformez les phrases directes suivantes en discours indirect avec le verbe entre parenthèses au Konjunktiv I. \tcblower
Solution :
1) \all{erklärte} (a expliqué) — le média-expert à Berlin ; \all{behauptete} (a affirmé) — le journaliste Thomas M. ; \all{riet} (a conseillé) — le journaliste Thomas M. ; \all{sagte} (a dit) — l'expert / le journaliste ; \all{warnte} (a mis en garde) — l'expert.
2) a) \all{„Die Nachricht ist falsch.“} $\rightarrow$ Er sagte, die Nachricht \textbf{sei} falsch. \\
   b) \all{„Wir prüfen die Quelle.“} $\rightarrow$ Sie erklärten, sie \textbf{prüften} die Quelle. (Ersatzform : \all{prüfen} identique indicatif $\rightarrow$ Konj. II \all{prüften}) \\
   c) \all{„Ich habe keine Zeit.“} $\rightarrow$ Er behauptete, er \textbf{habe} keine Zeit.
\end{exoresolu}

\courssub{Grammaire 2 : La subordonnée en « dass »}

Après les verbes de parole, de pensée ou de sentiment (\all{sagen, erklären, behaupten, glauben, denken, meinen, wissen}), le discours indirect s'introduit souvent par la conjonction \cle{\all{dass}}.

\begin{propriete}[Structure de la subordonnée en \all{dass}]
\begin{enumerate}
\item \all{dass} place le \cle{verbe conjugué à la fin} de la proposition.
\item Le verbe de la subordonnée est au \cle{Konjunktiv I} (discours indirect) ou à l'indicatif (discours direct rapporté / opinion personnelle).
\item L'ordre des mots dans la subordonnée : \cle{Sujet – Compléments – Verbe (Konj. I)}.
\end{enumerate}
\end{propriete}

\begin{exemplebox}[Exemples tirés du texte (reconstruits avec \all{dass})]
\begin{itemize}
\item Texte : \all{...erklärte, viele Jugendliche könnten...} $\rightarrow$ Reconstruit : \all{Der Experte erklärte, \textbf{dass} viele Jugendliche wahre und falsche Nachrichten kaum \textbf{unterscheiden könnten}.} (Verbe \all{könnten} à la fin).
\item Texte : \all{...behauptete, manche Fake News seien...} $\rightarrow$ Reconstruit : \all{Er behauptete, \textbf{dass} manche Fake News sogar von Regierungen bezahlt worden \textbf{seien}.} (Verbe \all{seien} à la fin).
\item Texte : \all{...riet, jede Nachricht zu überprüfen...} $\rightarrow$ Ici, \all{raten} régit l'\textbf{infinitif} (\all{zu} + infinitif), pas \all{dass}.
\end{itemize}
\end{exemplebox}

\begin{attention}[Pièges pour les francophones : \all{dass} vs \all{das} / Ordre des mots]
\begin{itemize}
\item \cle{\all{dass} (conjonction)} = « que » (subordonnée). S'écrit avec \textbf{ss}. Le verbe est \textbf{à la fin}.
\item \cle{\all{das} (article neutre / pronom relatif)} = « le / ce qui / qui ». S'écrit avec \textbf{un s}.
\item \textbf{Erreur fréquente} : *\all{Er sagte, dass er ist müde.} $\rightarrow$ Correct : \all{Er sagte, \textbf{dass} er müde \textbf{ist}.} (Verbe à la fin).
\item \textbf{Style journalistique} : Dans le texte original, \all{dass} est souvent \textbf{omis} (\all{erklärte, viele Jugendliche könnten...}). C'est le \cle{paratactischer Konjunktiv} (verbe en 2e position). À l'écrit scolaire et à l'oral standard, on \textbf{garde \all{dass}} + verbe à la fin.
\end{itemize}
\end{attention}

\begin{exercice}[Entraînement : \all{dass} + Konjunktiv I]
Transformez en discours indirect avec \all{dass} + Konjunktiv I (verbe à la fin).
1. \all{Der Experte: „Die Jugendlichen können die Nachrichten nicht unterscheiden.“} \\
2. \all{Thomas M.: „Manche Fake News sind von Regierungen bezahlt.“} \\
3. \all{Die Studie: „70\% der Jugendlichen lesen Nachrichten auf dem Handy.“}
\end{exercice}

\begin{corrige}[Exercice Grammaire 2]
1. \all{Der Experte erklärte, \textbf{dass} die Jugendlichen die Nachrichten nicht \textbf{unterscheiden könnten}.} \\
2. \all{Thomas M. behauptete, \textbf{dass} manche Fake News von Regierungen bezahlt \textbf{seien}.} \\
3. \all{Die Studie ergab, \textbf{dass} 70\% der Jugendlichen Nachrichten auf dem Handy \textbf{lesen}.} (Indicatif car fait objectif / étude, ou Konj. I \all{lesen} — forme identique, donc Konj. II \all{läsen} possible mais lourd ; ici indicatif accepté pour fait avéré).
\end{corrige}

\courssub{Méthode du commentaire et commentaire modèle}

\begin{methode}[Commentaire de texte journalistique (LV2 Allemand — Bac)]
\begin{enumerate}
\item \textbf{Introduction} (5-6 lignes) : \textbf{Accroche} (sujet d'actualité) $\rightarrow$ \textbf{Présentation} (titre, source, date, auteur, thème) $\rightarrow$ \textbf{Problématique} (question à laquelle le texte répond) $\rightarrow$ \textbf{Annonce du plan} (2 ou 3 parties).
\item \textbf{Développement} (2 parties équilibrées) : Chaque partie = \textbf{Idée directrice} (phrase) + \textbf{Argument} (explication) + \textbf{Citation/Exemple précis du texte} (en allemand, entre guillemets „…“) + \textbf{Analyse linguistique} (vocabulaire, grammaire, figures de style).
\item \textbf{Conclusion} (3-4 lignes) : \textbf{Bilan} (réponse à la problématique) $\rightarrow$ \textbf{Ouverture} (lien avec le contexte camerounais, autre texte, avis personnel).
\end{enumerate}
\end{methode}

\begin{exemplebox}[Commentaire modèle : « Fake News – die neue Gefahr im Internet »]
\textbf{Introduction}\\
Avec la numérisation croissante, la désinformation est devenue un fléau mondial. L'article „Fake News – die neue Gefahr im Internet“, publié dans un journal allemand, traite de la propagation des fausses informations sur les réseaux sociaux et de leurs conséquences. Comment les fake news se propagent-elles et quelles réponses la société apporte-t-elle ? Nous analyserons d'abord le constat alarmant dressé par l'expert, puis les causes et les solutions proposées.

\textbf{Développement}\\
\textbf{1. Un constat alarmant : la vitesse et l'anonymat de la désinformation.}\\
Le texte ouvre sur une observation sans appel : \all{„viele Jugendliche könnten wahre und falsche Nachrichten kaum unterscheiden“}. L'usage du \cle{Konjunktiv I} (\all{könnten, verbreiteten, gebe}) instaure une distance journalistique. L'anonymat (\all{„keinen Absender und keine Quelle“}) et la viralité (\all{„mehr als hunderttausend Mal geteilt“ en deux heures}) créent un cercle vicieux où le démenti (\all{dementieren}) arrive trop tard.

\textbf{2. Causes psychologiques et réponses institutionnelles.}\\
L'expert identifie trois facteurs : l'émotion (\all{„emotional und aufregend“}), la confiance aveugle (\all{„von Freunden weitergeleitet, denen man vertraue“}) et la paresse cognitive (\all{„fehle vielen Nutzern die Zeit“}). Face à cela, la réponse est double : juridique avec le \all{Netzwerkdurchsetzungsgesetz} (suppression rapide) et éducative avec la \all{Medienkompetenz} à l'école (\all{„Quellen zu vergleichen und Beweise zu suchen“}). La conclusion de l'expert, \all{„Erst denken, dann klicken“}, résume l'impératif de vigilance.

\textbf{Conclusion}\\
Le texte démontre que la lutte contre les fake news exige une responsabilité partagée : plateformes, État et citoyens. Cette analyse fait écho à la situation camerounaise où WhatsApp et Facebook sont vecteurs de rumeurs ; l'éducation aux médias y devient urgente pour former des citoyens critiques.
\end{exemplebox}

\begin{savaistu}[Le NetzDG : la loi allemande contre la haine en ligne]
En Allemagne, la loi NetzDG (\all{Netzwerkdurchsetzungsgesetz}), adoptée en 2017, oblige les réseaux sociaux à supprimer rapidement les contenus haineux et les fake news, sous peine d'amendes pouvant atteindre 50 millions d'euros. C'était l'une des premières lois de ce type en Europe, et elle a inspiré des débats dans de nombreux pays, y compris en Afrique. Au Cameroun, la loi n°2010/012 du 21 décembre 2010 relative à la cybercriminalité sanctionne aussi la propagation de fausses nouvelles (\all{Art. 78}).
\end{savaistu}

\begin{exercice}[À vous — Synthèse de la leçon]
\begin{enumerate}
\item \textbf{Compréhension / Expression orale :} Répondez en allemand par des phrases complètes :
      \begin{enumerate}
      \item \all{Was sind Fake News?}
      \item \all{Warum verbreiten sich Falschmeldungen so schnell?}
      \item \all{Was rät der Experte den Nutzern?}
      \item Nommez deux exemples de Konjunktiv I dans le dernier paragraphe du texte.
      \end{enumerate}
\item \textbf{Traduction (Thème) :} Traduisez en allemand en utilisant le discours indirect (\all{dass} + Konjunktiv I).
      \begin{enumerate}
      \item Le journaliste affirme que les fake news sont dangereuses. (\all{behaupten / gefährlich})
      \item L'expert explique que beaucoup de jeunes ne vérifient pas les sources. (\all{erklären / überprüfen})
      \item Il a dit que la loi oblige les plateformes à supprimer les contenus illégaux. (\all{sagen / verpflichten / löschen})
      \end{enumerate}
\item \textbf{Production écrite :} Rédigez un court paragraphe (60-80 mots) en allemand : \all{„Wie kann man sich vor Fake News schützen?“} Utilisez au moins trois verbes de parole au discours indirect (\all{meiner Meinung nach, man sollte, es ist wichtig, dass...}).
\end{enumerate}
\end{exercice}

\begin{corrige}[Exercice « À vous »]
\begin{enumerate}
\item \textbf{Compréhension :}
      1. \all{Fake News sind Falschmeldungen, die sich im Internet verbreiten.} \\
      2. \all{Sie verbreiten sich schnell, weil sie emotional sind, von Freunden geteilt werden und viele Nutzer keine Zeit zur Prüfung haben.} \\
      3. \all{Der Experte rät, jede Nachricht zu überprüfen, bevor man sie teilt. Erst denken, dann klicken.} \\
      4. \all{seien} (\all{sind}), \all{fehle} (\all{fehlt}), \all{vertraue} (\all{vertraut}), \all{scheinen} (indicatif ici car subordonnée relative \all{in denen... scheinen} — attention, pas de discours indirect direct ici), \all{werde} (\all{wird} — \all{werden} au passif ? Non, \all{wird ... gelöscht} indicatif. Le Konj. I du dernier paragraphe : \all{seien, fehle, vertraue}). 
\item \textbf{Traduction :}
      1. \all{Der Journalist behauptet, dass Fake News gefährlich \textbf{seien}.} \\
      2. \all{Der Experte erklärt, dass viele Jugendliche die Quellen nicht \textbf{überprüften}.} (Ersatzform Konj. II car \all{überprüfen} indicatif = \all{überprüfen} Konj. I pour \all{wir/sie}). \\
      3. \all{Er sagte, dass das Gesetz die Plattformen verpflichte, illegale Inhalte zu löschen.}
\item \textbf{Production écrite (proposition) :} \\
\all{Meiner Meinung nach ist Medienkompetenz der beste Schutz. Experten erklären, dass man Quellen vergleichen solle. Man sollte nicht alles glauben, was man liest. Es ist wichtig, dass man vor dem Teilen nachdenkt. Wer vorsichtig ist, wird nicht zum Problem.}
\end{enumerate}
\end{corrige}

\begin{attention}[Bilan des erreurs fréquentes à éviter pour le Bac]
\begin{itemize}
\item \textbf{Majuscules :} \all{die Wahrheit, der Beweis, die Quelle, das Gesetz, die Plattform} — \textbf{Tous les noms} en majuscule.
\item \textbf{Konjunktiv I :} \all{er macht} $\rightarrow$ \all{er mache} (pas *\all{er macht}); \all{er ist} $\rightarrow$ \all{er sei}; \all{er hat} $\rightarrow$ \all{er habe}.
\item \textbf{Formes identiques :} \all{wir haben} (Indik.) $\rightarrow$ \all{wir \textbf{hätten}} (Konj. II Ersatzform) ; \all{sie lesen} $\rightarrow$ \all{sie \textbf{liesen}}.
\item \textbf{\all{dass}-Satz :} Verbe \textbf{à la fin} ! \all{...dass er die Nachricht \textbf{liest}.}
\item \textbf{Faux amis :} \all{bekommen} = recevoir (pas devenir) ; \all{werden} = devenir / auxiliaire futur/passif.
\item \textbf{Verbes à particule :} \all{weiterleiten} (séparable) $\rightarrow$ \all{er leitet weiter} / \all{weitergeleitet} ; \all{überprüfen} (inséparable) $\rightarrow$ \all{er überprüft} / \all{überprüft}.
\end{itemize}
\end{attention}

\coursec{Compréhension du texte}

Dans la section précédente, nous avons découvert le texte support \emph{Fake News im Netz}, un dialogue entre une journaliste et un expert en médias sur le phénomène des fausses informations. Nous avons lu le texte une première fois pour le plaisir et repéré les mots nouveaux. Il est maintenant temps de passer à la \cle{compréhension} proprement dite : d'abord une lecture globale (de quoi parle le texte ?), puis une lecture détaillée (que dit exactement chaque personnage ?), enfin des exercices de vérification (vrai/faux, équivalents, reformulation). C'est exactement la démarche exigée à l'épreuve d'allemand du Baccalauréat A.

\courssub{La méthode : comment répondre à une question de compréhension}

\begin{methode}[Répondre à une question de compréhension en quatre étapes]
\begin{enumerate}
\item \textbf{Lire la question attentivement} et souligner le mot interrogatif : \all{wer} (qui), \all{was} (quoi), \all{wo} (où), \all{wann} (quand), \all{warum} (pourquoi), \all{wie} (comment), \all{worum} (de quoi).
\item \textbf{Repérer le passage du texte} qui contient la réponse. Souligner la phrase ou le groupe de mots utile.
\item \textbf{Rédiger une phrase complète en allemand} qui reprend les mots de la question. On ne recopie jamais tout le paragraphe : on sélectionne l'information demandée.
\item \textbf{Citer le texte entre guillemets allemands} „…“ si la consigne le demande (\all{Begründen Sie mit einer Textstelle} : justifiez par un passage du texte).
\end{enumerate}
\end{methode}

\begin{attention}[Le piège du « ja/nein » sec]
Au Baccalauréat, ne réponds \textbf{jamais} par \all{ja} ou \all{nein} seul, même si la question appelle oui ou non. Il faut \textbf{toujours justifier par le texte}. Mauvaise réponse : \all{Ja.} Bonne réponse : \all{Ja, denn der Experte sagt: „Falschmeldungen verbreiten sich schneller als die Wahrheit.“} (Oui, car l'expert dit : « Les fausses informations se répandent plus vite que la vérité. »)
\end{attention}

\courssub{Compréhension globale : de quoi parle le texte ?}

Avant de répondre, rappelons brièvement le texte support sur lequel portent toutes les questions de cette section.

\begin{textebox}[Fake News im Netz (rappel du texte support)]
\textbf{Journalistin:} Guten Tag, Herr Dr.~Weber. Sie sind Medienexperte. Was genau sind „Fake News“?\\
\textbf{Experte:} Fake News sind Falschmeldungen, also Nachrichten, die nicht wahr sind. Sie sehen oft aus wie echte Artikel, aber sie lügen.\\
\textbf{Journalistin:} Warum verbreiten sich diese Lügen so schnell?\\
\textbf{Experte:} Weil viele Menschen Nachrichten teilen, ohne sie zu prüfen. Eine aufregende Schlagzeile reist in wenigen Minuten um die Welt. Falschmeldungen verbreiten sich oft schneller als die Wahrheit.\\
\textbf{Journalistin:} Wer steht hinter diesen Falschmeldungen?\\
\textbf{Experte:} Manchmal wollen Leute Geld verdienen, manchmal wollen sie Politiker oder Firmen schaden. Es gibt auch Jugendliche, die nur Spaß haben wollen.\\
\textbf{Journalistin:} Wie kann man eine Falschmeldung erkennen?\\
\textbf{Experte:} Man muss zuerst die Quelle prüfen: Wer hat den Text geschrieben? Dann sucht man Beweise: Gibt es Fotos, Zahlen, Namen? Und man vergleicht mit anderen seriösen Medien.\\
\textbf{Journalistin:} Was raten Sie den Jugendlichen in Kamerun, zum Beispiel in Yaoundé oder Douala?\\
\textbf{Experte:} Glaubt nicht alles, was ihr auf WhatsApp oder Facebook lest! Prüft die Quelle, sucht Beweise und fragt eure Lehrer. Wer die Wahrheit liebt, prüft zuerst und teilt dann.
\end{textebox}

\textbf{Glossaire :} die Falschmeldung, -en (fausse information) ; die Nachricht, -en (information, nouvelle) ; die Lüge, -n (mensonge) ; lügen (mentir) ; die Wahrheit, -en (vérité) ; sich verbreiten (se répandre) ; teilen (partager) ; prüfen (vérifier) ; die Quelle, -n (source) ; der Beweis, -e (preuve) ; schaden + datif (nuire à) ; erkennen (reconnaître) ; raten + datif (conseiller à) ; seriös (sérieux, fiable) ; die Schlagzeile, -n (gros titre) ; aufregend (passionnant, sensationnel).

\begin{exoresolu}[Questions de compréhension globale]
\textbf{Énoncé.} Réponds par des phrases complètes en allemand.
\begin{enumerate}
\item \all{Worum geht es in diesem Text?} (De quoi parle ce texte ?)
\item \all{Wer spricht im Text?} (Qui parle dans le texte ?)
\item \all{Wo und wann spielt die Szene?} (Où et quand se déroule la scène ?)
\end{enumerate}
\tcblower
\textbf{Solution détaillée.}
\begin{enumerate}
\item \all{Der Text handelt von Fake News, also von Falschmeldungen im Internet, und davon, wie man sie erkennen kann.} — Le texte parle des fake news, c'est-à-dire des fausses informations sur Internet, et de la manière de les reconnaître. \emph{Méthode : le mot interrogatif \all{worum} (de quoi) appelle la formule \all{der Text handelt von + datif} (le texte traite de).}
\item \all{Eine Journalistin und ein Medienexperte, Herr Dr. Weber, sprechen im Text. Die Journalistin stellt Fragen, der Experte antwortet.} — Une journaliste et un expert en médias, M. le Dr Weber, parlent dans le texte. La journaliste pose des questions, l'expert répond.
\item \all{Der Text nennt keinen genauen Ort und keine genaue Zeit; es ist ein Interview, das heute stattfinden könnte. Der Experte erwähnt aber Kamerun, Yaoundé und Douala.} — Le texte n'indique ni lieu précis ni moment précis ; c'est une interview qui pourrait avoir lieu aujourd'hui. L'expert mentionne cependant le Cameroun, Yaoundé et Douala. \emph{Attention : quand le texte ne dit pas quelque chose, il faut le dire honnêtement (\all{der Text nennt keinen Ort}) et non l'inventer.}
\end{enumerate}
\end{exoresolu}

\courssub{Compréhension détaillée : l'expert et la journaliste}

Passons maintenant aux questions de détail. Pour chaque question, nous indiquons le passage du texte (\all{die Textstelle}) qui justifie la réponse.

\begin{center}
\begin{tabularx}{\linewidth}{|X|X|}
\hline
\rowcolor{popblueL} \textbf{Frage (question)} & \textbf{Textstelle (passage justificatif)} \\
\hline
\all{Was sind Fake News?} & \all{„Fake News sind Falschmeldungen, also Nachrichten, die nicht wahr sind.“} \\
\hline
\all{Warum verbreiten sich die Lügen so schnell?} & \all{„Weil viele Menschen Nachrichten teilen, ohne sie zu prüfen.“} \\
\hline
\all{Wer steht hinter den Falschmeldungen?} & \all{„Manchmal wollen Leute Geld verdienen, manchmal wollen sie Politiker oder Firmen schaden.“} \\
\hline
\all{Wie erkennt man eine Falschmeldung?} & \all{„Man muss zuerst die Quelle prüfen … Dann sucht man Beweise …“} \\
\hline
\all{Was rät der Experte den Jugendlichen?} & \all{„Prüft die Quelle, sucht Beweise und fragt eure Lehrer.“} \\
\hline
\end{tabularx}
\end{center}

\begin{exemplebox}[Réponses rédigées et traduites]
\begin{itemize}
\item \all{Der Experte sagt, dass Fake News Nachrichten sind, die nicht wahr sind, aber wie echte Artikel aussehen.} « L'expert dit que les fake news sont des informations qui ne sont pas vraies mais qui ressemblent à de vrais articles. »
\item \all{Die Lügen verbreiten sich schnell, weil viele Menschen Nachrichten teilen, ohne sie zu prüfen.} « Les mensonges se répandent vite parce que beaucoup de gens partagent des informations sans les vérifier. »
\item \all{Der Experte rät den Jugendlichen, nicht alles zu glauben und zuerst die Quelle zu prüfen.} « L'expert conseille aux jeunes de ne pas tout croire et de vérifier d'abord la source. »
\end{itemize}
Remarque la grammaire : après \all{sagen, dass}, le verbe conjugué va à la fin de la subordonnée (\all{… nicht wahr \textbf{sind}}). Remarque aussi le verbe à particule \all{aussehen} (ressembler) : la particule \all{aus} se détache et va en fin de proposition (\all{… wie echte Artikel \textbf{aussehen}}). Nous étudierons la subordonnée en \all{dass} en détail dans la section 5.
\end{exemplebox}

\courssub{Vrai ou faux ? Justifier par une citation}

\begin{exercice}[Richtig oder falsch ?]
Lis les phrases suivantes. Écris \all{richtig} ou \all{falsch} et justifie ta réponse par une citation du texte entre guillemets „…“.
\begin{enumerate}
\item \all{Fake News sehen nie aus wie echte Artikel.}
\item \all{Falschmeldungen verbreiten sich oft schneller als die Wahrheit.}
\item \all{Hinter den Falschmeldungen stehen immer Politiker.}
\item \all{Man soll zuerst die Quelle prüfen.}
\item \all{Der Experte rät den Jugendlichen, alles auf Facebook zu glauben.}
\end{enumerate}
\end{exercice}

\begin{corrige}[Exercice « Richtig oder falsch ? »]
\begin{enumerate}
\item \textbf{Falsch.} \all{Der Experte sagt: „Sie sehen oft aus wie echte Artikel, aber sie lügen.“} (Faux. L'expert dit : « Elles ressemblent souvent à de vrais articles, mais elles mentent. »)
\item \textbf{Richtig.} \all{„Falschmeldungen verbreiten sich oft schneller als die Wahrheit.“} (Vrai. « Les fausses informations se répandent souvent plus vite que la vérité. »)
\item \textbf{Falsch.} \all{Der Text sagt: „Manchmal wollen Leute Geld verdienen … Es gibt auch Jugendliche, die nur Spaß haben wollen.“} (Faux : le texte dit « parfois », pas « toujours », et mentionne aussi des jeunes qui veulent s'amuser.)
\item \textbf{Richtig.} \all{„Man muss zuerst die Quelle prüfen.“} (Vrai. « Il faut d'abord vérifier la source. »)
\item \textbf{Falsch.} \all{Der Experte sagt das Gegenteil: „Glaubt nicht alles, was ihr auf WhatsApp oder Facebook lest!“} (Faux : l'expert dit le contraire — « Ne croyez pas tout ce que vous lisez sur WhatsApp ou Facebook ! »)
\end{enumerate}
\end{corrige}

\begin{attention}[Les mots-pièges : \all{immer}, \all{nie}, \all{alle}]
Dans les exercices vrai/faux, méfie-toi des phrases contenant \all{immer} (toujours), \all{nie} (jamais), \all{alle} (tous) : elles exagèrent presque toujours le texte, qui dit souvent \all{manchmal} (parfois), \all{oft} (souvent), \all{viele} (beaucoup). Compare mot à mot la phrase de l'exercice avec la phrase du texte.
\end{attention}

\courssub{Recherche d'équivalents dans le texte}

Exercice classique du Bac : on te donne un mot français (ou une définition allemande) et tu dois retrouver le mot du texte qui correspond.

\begin{center}
\begin{tabularx}{\linewidth}{|X|X|X|}
\hline
\rowcolor{popblueL} \textbf{Français} & \textbf{Mot du texte} & \textbf{Phrase du texte} \\
\hline
mensonge & die Lüge, -n & \all{„aber sie lügen“} / \all{„diese Lügen“} \\
\hline
vérifier & prüfen & \all{„ohne sie zu prüfen“} \\
\hline
source & die Quelle, -n & \all{„Man muss zuerst die Quelle prüfen“} \\
\hline
preuve & der Beweis, -e & \all{„Dann sucht man Beweise“} \\
\hline
se répandre & sich verbreiten & \all{„verbreiten sich so schnell“} \\
\hline
partager & teilen & \all{„viele Menschen Nachrichten teilen“} \\
\hline
conseiller & raten (+ dat.) & \all{„Was raten Sie den Jugendlichen?“} \\
\hline
\end{tabularx}
\end{center}

\begin{attention}[Recopier sans faute]
Quand tu recopies un mot du texte, respecte la \textbf{majuscule des noms} (\all{die Quelle}, pas \all{die quelle}) et les trémas (\all{prüfen}, pas \all{prufen}). Un mot mal orthographié peut coûter des points, même si l'idée est bonne. Attention aussi au verbe \all{sich verbreiten} : il est \textbf{pronominal} ; dans la réponse, n'oublie pas le pronom réfléchi (\all{die Lügen verbreiten \textbf{sich}}).
\end{attention}

\courssub{Reformulation : dire l'idée principale avec ses propres mots}

Reformuler, c'est exprimer l'idée d'un passage \textbf{sans recopier les phrases du texte}. C'est une compétence essentielle pour le résumé (\all{die Zusammenfassung}).

\begin{exemplebox}[Reformulation réplique par réplique]
\begin{itemize}
\item \textbf{Répliques 1–2 (définition) :} \all{Fake News sind Lügen, die wie wahre Nachrichten aussehen.} « Les fake news sont des mensonges qui ont l'apparence d'informations vraies. »
\item \textbf{Répliques 3–4 (rapidité) :} \all{Viele Leute teilen Nachrichten zu schnell, deshalb reisen Lügen schneller als die Wahrheit.} « Beaucoup de gens partagent trop vite, c'est pourquoi les mensonges voyagent plus vite que la vérité. »
\item \textbf{Répliques 5–6 (auteurs) :} \all{Die Urheber der Fake News wollen Geld verdienen, anderen schaden oder nur Spaß haben.} « Les auteurs des fake news veulent gagner de l'argent, nuire aux autres ou simplement s'amuser. »
\item \textbf{Répliques 7–8 (détection) :} \all{Um eine Lüge zu erkennen, kontrolliert man die Quelle und sucht Beweise.} « Pour reconnaître un mensonge, on contrôle la source et on cherche des preuves. »
\item \textbf{Répliques 9–10 (conseil) :} \all{Jugendliche sollen kritisch sein: erst prüfen, dann teilen.} « Les jeunes doivent être critiques : d'abord vérifier, ensuite partager. »
\end{itemize}
\end{exemplebox}

\begin{propriete}[Technique de reformulation]
Pour reformuler sans plagier : (1) remplace les mots par des \textbf{synonymes} (\all{Falschmeldung} $\rightarrow$ \all{Lüge} ; \all{prüfen} $\rightarrow$ \all{kontrollieren}) ; (2) change la \textbf{structure de la phrase} (deux phrases $\rightarrow$ une phrase avec \all{um … zu} ou \all{deshalb} ; tournure personnelle $\rightarrow$ tournure avec \all{man}) ; (3) garde uniquement \textbf{l'idée essentielle}, supprime les exemples et les détails.
\end{propriete}

\begin{popfigure}
\begin{tikzpicture}[mindmap, grow cyclic, every node/.style={concept, align=center}, root concept/.append style={concept color=popblue, font=\small\bfseries}, level 1/.append style={level distance=3.2cm, sibling angle=90}, level 1 concept/.append style={concept color=poporange!70, font=\small}]
\node[root concept]{Fake News im Netz}
  child { node[align=center]{Was?\\ Falschmeldungen,\\ die lügen} }
  child { node[align=center]{Warum so schnell?\\ Teilen ohne\\ Prüfen} }
  child { node[align=center]{Wer?\\ Geld, Schaden,\\ Spaß} }
  child { node[align=center]{Was tun?\\ Quelle prüfen,\\ Beweise suchen} };
\end{tikzpicture}
\legende{Carte mentale de l'idée principale de chaque partie du texte : un bon résumé reprend ces quatre idées.}
\end{popfigure}

\courssub{Question d'opinion personnelle}

La dernière question d'une épreuve de compréhension invite souvent à donner ton avis. Ici, tu peux répondre librement, mais toujours en allemand, avec des phrases complètes et un argument.

\begin{exemplebox}[Verbreitest du manchmal Nachrichten, ohne sie zu prüfen ?]
Réponse possible 1 (honnête) : \all{Ja, manchmal habe ich eine Nachricht auf WhatsApp geteilt, ohne sie zu prüfen. Aber jetzt weiß ich, dass das gefährlich ist. Ich will zuerst die Quelle prüfen.} « Oui, il m'est arrivé de partager une information sur WhatsApp sans la vérifier. Mais maintenant je sais que c'est dangereux. Je veux d'abord vérifier la source. »

Réponse possible 2 : \all{Nein, ich bin sehr vorsichtig. Ich glaube nicht alles, was ich im Internet lese, und ich vergleiche immer mit anderen Medien.} « Non, je suis très prudent(e). Je ne crois pas tout ce que je lis sur Internet et je compare toujours avec d'autres médias. »
\end{exemplebox}

\begin{methode}[Donner son opinion en trois phrases]
\begin{enumerate}
\item \textbf{Position} : \all{Ich finde …} / \all{Meiner Meinung nach …} (à mon avis).
\item \textbf{Argument} : \all{weil / denn + proposition} (\all{…, weil Falschmeldungen gefährlich sind.}).
\item \textbf{Exemple personnel} : \all{Zum Beispiel habe ich …} (par exemple, j'ai …).
\end{enumerate}
Attention à la place du verbe : après \all{weil}, le verbe conjugué va à la fin ; après \all{denn}, il reste en deuxième position.
\end{methode}

\begin{experience}[Débat en classe : Erst prüfen, dann teilen]
Par groupes de quatre : deux élèves défendent la thèse \all{„Man darf Nachrichten sofort teilen“}, deux élèves défendent \all{„Man muss immer zuerst prüfen“}. Chaque groupe prépare trois arguments avec le vocabulaire de la leçon (\all{die Quelle, der Beweis, die Wahrheit, verbreiten, glauben}), puis le débat a lieu devant la classe. Le professeur note les erreurs de place du verbe et y revient à la fin.
\end{experience}

\begin{aretenir}[L'essentiel de la compréhension de texte]
\begin{itemize}
\item Réponds toujours par une \textbf{phrase complète en allemand}, jamais par \all{ja}/\all{nein} seul.
\item Justifie par une \textbf{citation entre guillemets} „…“ quand la consigne le demande.
\item Méfie-toi des mots absolus \all{immer}, \all{nie}, \all{alle} dans les vrai/faux.
\item Pour reformuler : synonymes + nouvelle structure + idée essentielle seule.
\item Pour l'opinion : position + argument (\all{weil}/\all{denn}) + exemple personnel.
\end{itemize}
\end{aretenir}

Nous savons désormais lire et exploiter le texte en profondeur. La section suivante, \emph{Lexique thématique : Wahrheit und Lüge}, nous permettra d'élargir et de structurer le vocabulaire rencontré ici.

\coursec{Lexique thématique : Wahrheit und Lüge}

Après avoir travaillé la compréhension globale et détaillée du texte \emph{Fake News im Netz}, nous allons maintenant fixer et structurer le vocabulaire essentiel pour parler de la désinformation, de la vérité et du mensonge en allemand. Ce lexique thématique est indispensable pour réussir l'expression écrite (résumé, commentaire) et orale (débat, exposé) sur ce thème d'actualité brûlant.

\courssub{Champ 1 : L'information et la vérité (Die Information und die Wahrheit)}

Commençons par les mots qui désignent le fait vrai, la donnée vérifiable et sa source. Observez ces exemples tirés du texte ou construits sur le même modèle :

\begin{exemplebox}[Observation]
\all{Die \textbf{Nachricht} über den Wahlsieg kam um 20 Uhr.} — « La nouvelle de la victoire électorale est arrivée à 20 h. »\\
\all{Wir brauchen verlässliche \textbf{Informationen}, keine Gerüchte.} — « Nous avons besoin d'informations fiables, pas de rumeurs. »\\
\all{Die \textbf{Wahrheit} kommt oft erst später ans Licht.} — « La vérité finit souvent par éclater plus tard. » (litt. : vient à la lumière) \\
\all{Es ist eine \textbf{Tatsache}, dass das Video gefälscht ist.} — « C'est un fait que la vidéo est truquée. »\\
\all{Der \textbf{Beweis} fehlt noch.} — « La preuve manque encore. »\\
\all{Die \textbf{Quelle} ist seriös und unabhängig.} — « La source est sérieuse et indépendante. »
\end{exemplebox}

\begin{definition}[Famille \cle{wahr} — Vérité et vraisemblance]
Le radical \cle{wahr} (vrai) structure tout le champ de la véracité.
\begin{itemize}
    \item \all{das Adjektiv: wahr} (vrai, exact) — \all{unwahr} (faux, inexact).
    \item \all{das Substantiv: die Wahrheit} (la vérité) — Pl. \all{die Wahrheiten}.
    \item \all{das Adjektiv: wahr\-schein\-lich} (probable, vraisemblable) — \all{unwahrscheinlich} (improbable).
    \item \all{das Substantiv: die Wahrscheinlichkeit} (la probabilité).
    \item \all{das Verb: wahrnehmen} (percevoir, constater) — \emph{verbe à particule séparable :} \all{er nimmt wahr}.
\end{itemize}
\end{definition}

\begin{propriete}[Genre, pluriel et emploi — Noms du champ 1]
\begin{center}
\begin{tabularx}{\linewidth}{|X|X|X|X|}
\hline
\rowcolor{popblueL}
\textbf{Deutsch (Artikel + Plural)} & \textbf{Français} & \textbf{Remarque / Collocation} & \textbf{Exemple} \\
\hline
die Nachricht, die \cle{-en} & la nouvelle, l'info & \all{eine Nachricht verbreiten / erhalten} & \all{Eine gute Nachricht.} \\
\hline
die Information, die \cle{-en} & l'information & Souvent au pluriel. \all{Informationen sammeln} & \all{Weitere Informationen auf der Website.} \\
\hline
die Wahrheit, die \cle{-en} & la vérité & \all{die Wahrheit sagen / suchen / vertuschen} & \all{Die ganze Wahrheit.} \\
\hline
die Tatsache, die \cle{-n} & le fait, la réalité & \all{eine Tatsache feststellen / ignorieren} ; \all{in der Tat} (en effet) & \all{Das ist eine Tatsache.} \\
\hline
der Beweis, die \cle{-e} & la preuve & \all{einen Beweis liefern / erbringen / antreten} & \all{Ohne Beweis keine Verurteilung.} \\
\hline
die Quelle, die \cle{-n} & la source & \all{eine Quelle zitieren / prüfen / angeben} ; \all{Quellenangabe} & \all{Die Quelle ist vertrauenswürdig.} \\
\hline
\end{tabularx}
\end{center}
\end{propriete}

\courssub{Champ 2 : Le mensonge et la manipulation (Die Lüge und die Manipulation)}

Voici le vocabulaire pour nommer la fausse information, qu'elle soit intentionnelle (mensonge, propagande) ou non (rumeur), et les mécanismes de tromperie.

\begin{textebox}[Mini-texte d'observation : Die Mechanismen der Täuschung]
\all{Eine \textbf{Falschmeldung} verbreitet sich oft schneller als die Richtigstellung. Wer eine \textbf{Lüge} erzählt, muss sich oft neue Lügen ausdenken, um die alte zu decken. \textbf{Propaganda} nutzt Emotionen, nicht Fakten. Hinter jeder \textbf{Täuschung} steckt eine Absicht: der \textbf{Betrug} um Geld, Macht oder Aufmerksamkeit. Ein \textbf{Gerücht} dagegen entsteht oft aus Unwissenheit, wirkt aber genauso zerstörerisch.}
\glossaire{
    \item \all{die Falschmeldung (-en)} : la fausse info, la \emph{fake news} (mot composé : \cle{falsch} + \cle{die Meldung}).
    \item \all{die Richtigstellung (-en)} : la rectification, le démenti.
    \item \all{sich ausdenken} : inventer (imparfait \all{dachte aus}, parfait \all{hat ausgedacht}).
    \item \all{decken} : couvrir, cacher.
    \item \all{zerstörerisch} : destructeur.
}
\end{textebox}

\begin{definition}[\cle{die Falschmeldung} (die \cle{-en})]
\emph{Information fausse diffusée volontairement (ou par négligence grave) dans les médias ou sur les réseaux pour tromper le public, manipuler l'opinion ou générer du clic (clickbait).} C'est l'équivalent allemand le plus précis de \emph{Fake News} (terme anglais aussi utilisé en allemand). Composition : \all{falsch} (faux) + \all{die Meldung} (le signalement, la dépêche, la nouvelle officielle).
\end{definition}

\begin{propriete}[Genre, pluriel et collocations — Noms du champ 2]
\begin{center}
\begin{tabularx}{\linewidth}{|X|X|X|X|}
\hline
\rowcolor{poporangeL}
\textbf{Deutsch (Artikel + Plural)} & \textbf{Français} & \textbf{Collocations types} & \textbf{Exemple} \\
\hline
die Lüge, die \cle{-n} & le mensonge & \all{eine Lüge erzählen / verbreiten / entlarven} & \all{Das war eine glatte Lüge.} \\
\hline
die Falschmeldung, die \cle{-en} & la fausse nouvelle & \all{eine Falschmeldung starten / korrigieren / melden} & \all{Die Falschmeldung ging viral.} \\
\hline
das Gerücht, die \cle{-e} & la rumeur & \all{ein Gerücht streuen / widerlegen / bestätigen} & \all{Es sind nur Gerüchte.} \\
\hline
die Propaganda (kein Pl.) & la propagande & \all{Propaganda betreiben / machen / verbreiten} & \all{Kriegspropaganda.} \\
\hline
die Täuschung, die \cle{-en} & la tromperie, la duperie & \all{jemanden täuschen / eine Täuschung aufdecken} & \all{Eine bewusste Täuschung.} \\
\hline
der Betrug (kein Pl.) & la fraude, l'arnaque & \all{Betrug begehen / melden / aufdecken} ; \all{der Betrüger} & \all{Internetbetrug nimmt zu.} \\
\hline
\end{tabularx}
\end{center}
\end{propriete}

\begin{attention}[Faux amis et pièges de traduction]
\begin{itemize}
    \item \all{die Lüge} $\neq$ \emph{la ligue}. C'est « le mensonge ». Verbe : \all{lügen} (fort : \all{lügt, log, hat gelogen}).
    \item \all{das Gerücht} $\neq$ \emph{le germe}. C'est « la rumeur ». Pas de lien avec \all{der Keim} (le germe).
    \item \all{die Propaganda} : en allemand, le mot est \emph{féminin} et \emph{sans pluriel} (comme en français), mais attention à la prononciation : [pʁoˈɡanda] (accent sur le 2e \emph{a}), pas « pro-pa-gan-da » à la française.
    \item \all{der Betrug} = la fraude/escroquerie (pénal). Ne pas confondre avec \all{das Betrügen} (substantivé : le fait de tricher, ex. à un jeu).
\end{itemize}
\end{attention}

\courssub{Champ 3 : Les actions — Verbes clés (Die Handlungen — Verben)}

Ces verbes sont le moteur de vos phrases. Maîtrisez leur conjugaison (Présent, Parfait, Prétérit) et leur régime (accusatif, datif, prépositionnel).

\begin{propriete}[Tableau des verbes thématiques — Conjugaison et régime]
\begin{center}
\begin{tabularx}{\linewidth}{|X|X|X|X|X|}
\hline
\rowcolor{popgreenL}
\textbf{Infinitif} & \textbf{Présent (3e sg.)} & \textbf{Präteritum (3e sg.)} & \textbf{Perfekt (Part. II)} & \textbf{Régime / Construction} \\
\hline
\all{verbreiten} (faible, inséparable) & \all{verbreitet} & \all{verbreitete} & \all{hat verbreitet} & \all{etw. (Akk) verbreiten} \\
\hline
\all{glauben} (faible) & \all{glaubt} & \all{glaubte} & \all{hat geglaubt} & \all{jm (Dat) / etw (Akk) glauben} \\
\hline
\all{lügen} (\textbf{fort}) & \all{lügt} & \all{log} & \all{hat gelogen} & \all{jm (Dat) / etw (Akk) lügen} \\
\hline
\all{täuschen} (faible) & \all{täuscht} & \all{täuschte} & \all{hat getäuscht} & \all{jm (Akk) täuschen} \\
\hline
\all{behaupten} (faible) & \all{behauptet} & \all{behauptete} & \all{hat behauptet} & \all{etw (Akk) / dass-Satz behaupten} \\
\hline
\all{überprüfen} (faible, \textbf{inséparable}) & \all{überprüft} & \all{überprüfte} & \all{hat überprüft} & \all{etw (Akk) überprüfen} \\
\hline
\all{teilen} (\textbf{fort}) & \all{teilt} & \all{teilte} & \all{hat geteilt} & \all{etw (Akk) teilen} (partager / poster) \\
\hline
\all{melden} (faible) & \all{meldet} & \all{meldete} & \all{hat gemeldet} & \all{etw (Akk) melden} (signaler) \\
\hline
\all{zugeben} (\textbf{fort}, séparable) & \all{gibt ... zu} & \all{gab ... zu} & \all{hat zugegeben} & \all{etw (Akk) / dass-Satz zugeben} \\
\hline
\all{leugnen} (faible) & \all{leugnet} & \all{leugnete} & \all{hat geleugnet} & \all{etw (Akk) leugnen} \\
\hline
\end{tabularx}
\end{center}
\end{propriete}

\begin{exemplebox}[Verbes en contexte — Phrases types]
\all{Der Täter \textbf{verbreitete} die Falschmeldung auf Twitter.} (Il a diffusé la fake news sur Twitter.)\\
\all{Viele Nutzer \textbf{glauben} alles, was sie lesen.} (Beaucoup d'utilisateurs croient tout ce qu'ils lisent.)\\
\all{Er \textbf{log} die Polizei an.} (Il a menti à la police. — \emph{Datif !} \all{jemandem lügen}.)\\
\all{Das Video \textbf{täuscht} den Zuschauer.} (La vidéo trompe le spectateur.)\\
\all{Der Politiker \textbf{behauptete}, er wisse von nichts.} (Le politicien a affirmé qu'il ne savait de rien. — \emph{Discours indirect !})\\
\all{Wir müssen die Quelle \textbf{überprüfen}.} (Nous devons vérifier la source.)\\
\all{Sie \textbf{teilte} den Beitrag ohne zu lesen.} (Elle a partagé le post sans lire.)\\
\all{Der Nutzer \textbf{meldete} den Account als Spam.} (L'utilisateur a signalé le compte comme spam.)\\
\all{Er \textbf{gab} schließlich zu, dass er gelogen hatte.} (Il a fini par admettre qu'il avait menti.)\\
\all{Der Verdächtige \textbf{leugnet} die Tat.} (Le suspect nie les faits.)
\end{exemplebox}

\begin{attention}[Erreurs fréquentes des francophones — Verbes]
\begin{itemize}
    \item \textbf{\all{lügen} est un verbe FORT} : \all{er lügt} (présent), \all{er log} (prétérit), \all{er hat gelogen} (parfait). \textbf{Jamais} \all{lügete} ou \all{hat lügt} !
    \item \textbf{\all{teilen} est FORT} : \all{er teilt, er teilte, er hat geteilt}. Participe passé \all{geteilt} (pas \all{teilt}).
    \item \textbf{\all{glauben} + personne = DATIF} : \all{Ich glaube \uline{ihm} nicht.} (Je ne le crois pas.) $\leftarrow$ Piège : en français « croire quelqu'un » est COD. En allemand : \all{jemandem (Dat) glauben}.
    \item \textbf{\all{zugeben} est séparable} : \all{Er gibt es zu.} / \all{Er gab es zu.} Particule \all{zu} en fin de proposition principale.
    \item \textbf{\all{melden} vs \all{sich anmelden}} : \all{melden} = signaler, déclarer ; \all{sich anmelden} = s'inscrire.
    \item \textbf{\all{überprüfen} est INSÉPARABLE} (préfixe \all{über-}) : \all{er überprüft, er überprüfte, hat überprüft}. \textbf{Jamais} \all{prüft ... über} !
\end{itemize}
\end{attention}

\courssub{Champ 4 : Les personnes (Die Personen)}

Qui agit ? Qui produit, qui relaie, qui vérifie, qui est victime ?

\begin{propriete}[Noms d'agents — Genre et pluriel]
\begin{center}
\begin{tabularx}{\linewidth}{|X|X|X|X|}
\hline
\rowcolor{poppurpleL}
\textbf{Deutsch (Artikel + Plural)} & \textbf{Français} & \textbf{Féminin / Dérivé} & \textbf{Exemple} \\
\hline
der Absender, die \cle{-} & l'expéditeur, l'auteur (du message) & die Absenderin & \all{Der Absender ist unbekannt.} \\
\hline
der Journalist, die \cle{-en} & le journaliste & die Journalistin & \all{Journalisten recherchieren Fakten.} \\
\hline
der Experte, die \cle{-n} & l'expert & die Expertin & \all{Ein Experte für Cybersicherheit.} \\
\hline
der Zeuge, die \cle{-n} & le témoin & die Zeugin & \all{Der Zeuge sagte aus.} \\
\hline
der Betrüger, die \cle{-} & l'escroc, le fraudeur & die Betrügerin & \all{Der Betrüger wurde gefasst.} \\
\hline
der Nutzer, die \cle{-} & l'utilisateur & die Nutzerin & \all{Die Nutzer melden den Post.} \\
\hline
der Verfasser, die \cle{-} & l'auteur (d'un texte) & die Verfasserin & \all{Der Verfasser des Artikels.} \\
\hline
\end{tabularx}
\end{center}
\end{propriete}

\begin{savaistu}[Le mot \emph{Fake News} en allemand]
L'anglicisme \all{die Fake News} (pluriel invariable, genre féminin par attraction de \all{die News}) est entré dans le Duden et dans l'usage courant (médias, politique, école). Il a été élu \all{Wort des Jahres} (Mot de l'année) en Allemagne en \textbf{2017} par la Gesellschaft für deutsche Sprache (GfdS). On trouve aussi \all{alternative Fakten} (« faits alternatifs »), calque de l'anglais \emph{alternative facts}, très utilisé depuis 2017 dans le débat public. Le terme purement germanique officiel reste \all{die Falschmeldung} (ou \all{die Desinformation} pour le processus global).
\end{savaistu}

\courssub{Familles de mots (Wortfamilien) — Construire son vocabulaire}

L'allemand est une langue à composition et dérivation puissantes. Apprendre un radical, c'est débloquer 5 à 10 mots.

\begin{definition}[Famille \cle{wahr} — Déclinaison complète]
\begin{center}
\begin{tabularx}{\linewidth}{|l|X|X|}
\hline
\rowcolor{popblueL}
\textbf{Mot de base} & \textbf{Dérivés / Composés} & \textbf{Sens / Traduction} \\
\hline
\all{wahr} (adj.) & \all{unwahr, wahrhaftig, wahrheitsgemäß} & faux / véridique / conforme à la vérité \\
\hline
\all{die Wahrheit} & \all{die Halbwahrheit, die Unwahrheit} & la demi-vérité / le mensonge (formel) \\
\hline
\all{wahrscheinlich} & \all{die Wahrscheinlichkeit, unwahrscheinlich} & probable / la probabilité / improbable \\
\hline
\all{wahrnehmen} (v. sép.) & \all{die Wahrnehmung, wahrnehmbar} & percevoir / la perception / perceptible \\
\hline
\end{tabularx}
\end{center}
\end{definition}

\begin{definition}[Famille \cle{lügen} — Du verbe au nom d'agent]
\begin{center}
\begin{tabularx}{\linewidth}{|l|X|X|}
\hline
\rowcolor{poporangeL}
\textbf{Mot de base} & \textbf{Dérivés / Composés} & \textbf{Sens / Traduction} \\
\hline
\all{lügen} (v. fort) & \all{die Lüge, der Lügner, die Lügnerin} & mentir / le mensonge / le menteur / la menteuse \\
\hline
 & \all{belügen, anlügen, verlügen} & mentir à qqn / mentir à qqn / mentir (résultat) \\
\hline
 & \all{die Lügengeschichte, das Lügenmärchen} & histoire mensongère / conte mensonger \\
\hline
\end{tabularx}
\end{center}
\end{definition}

\begin{definition}[Famille \cle{glauben} — Croire, confiance, crédibilité]
\begin{center}
\begin{tabularx}{\linewidth}{|l|X|X|}
\hline
\rowcolor{popgreenL}
\textbf{Mot de base} & \textbf{Dérivés / Composés} & \textbf{Sens / Traduction} \\
\hline
\all{glauben} (v. faible) & \all{der Glaube, der Glaubende} & croire / la foi, la croyance / celui qui croit \\
\hline
 & \all{glaubwürdig, die Glaubwürdigkeit} & crédible / la crédibilité \\
\hline
 & \all{unglaubwürdig, das Missvertrauen} & pas crédible / la méfiance \\
\hline
 & \all{der Glaubenssatz, der Glaubenskrieg} & article de foi / guerre de religion \\
\hline
\end{tabularx}
\end{center}
\end{definition}

\courssub{Expressions utiles (Redemittel) — Pour l'oral et l'écrit}

Ces blocs figés (chunks) vous permettent de produire des phrases correctes immédiatement.

\begin{center}
\begin{tabularx}{\linewidth}{|X|X|}
\hline
\rowcolor{poppinkL}
\textbf{Deutsch} & \textbf{Français} \\
\hline
\all{eine Nachricht / Falschmeldung / Gerücht verbreiten} & diffuser une nouvelle / une fake news / une rumeur \\
\hline
\all{die Wahrheit sagen / suchen / herausfinden} & dire la vérité / chercher la vérité / découvrir la vérité \\
\hline
\all{jemandem glauben / nicht glauben} & croire quelqu'un / ne pas croire quelqu'un (\textbf{Datif !}) \\
\hline
\all{Beweise liefern / erbringen / vorlegen} & apporter des preuves / fournir des preuves / produire des preuves \\
\hline
\all{einer Lüge / Falschmeldung auf den Grund gehen} & aller au fond d'un mensonge / d'une fake news (enquête) \\
\hline
\all{etwas als Falschmeldung entlarven / widerlegen} & démasquer qqch comme fake news / réfuter qqch \\
\hline
\all{die Quelle überprüfen / anzweifeln / angeben} & vérifier / mettre en doute / citer la source \\
\hline
\all{einen Account / Beitrag melden / blockieren} & signaler / bloquer un compte / un post \\
\hline
\all{sich gegen Desinformation wehren / schützen} & se défendre / se protéger contre la désinformation \\
\hline
\end{tabularx}
\end{center}

\courssub{Entraînement : Texte à trous et production guidée}

\begin{exoresolu}[Compléter le texte avec les mots du lexique (forme correcte)]
\textbf{Énoncé :} Complétez le texte suivant en conjuguant les verbes entre parenthèses au temps qui convient (Présent, Prétérit ou Parfait) et en mettant les noms au bon cas (Nominatif, Accusatif, Datif). Utilisez le vocabulaire des champs 1 à 4.

\all{Der \_\_\_\_\_\_\_\_\_\_ (Absender) der Nachricht blieb anonym. Er \_\_\_\_\_\_\_\_\_\_ (behaupten), er \_\_\_\_\_\_\_\_\_\_ (haben) Beweise für einen Skandal. Viele Nutzer \_\_\_\_\_\_\_\_\_\_ (glauben) ihm sofort und \_\_\_\_\_\_\_\_\_\_ (teilen) den Post tausendfach. Doch ein \_\_\_\_\_\_\_\_\_\_ (Experte) \_\_\_\_\_\_\_\_\_\_ (überprüfen) die \_\_\_\_\_\_\_\_\_\_ (Quelle) und \_\_\_\_\_\_\_\_\_\_ (entdecken), dass es sich um eine \_\_\_\_\_\_\_\_\_\_ (Falschmeldung) handelte. Der \_\_\_\_\_\_\_\_\_\_ (Betrüger) \_\_\_\_\_\_\_\_\_\_ (zusammenstellen) alte Fotos. Am Ende \_\_\_\_\_\_\_\_\_\_ (müssen) die Plattform den Account \_\_\_\_\_\_\_\_\_\_ (sperren).}

\tcblower
\textbf{Solution détaillée :}
\begin{enumerate}
    \item \all{Der \textbf{Absender} der Nachricht} (Nominatif singulier, sujet).
    \item \all{\textbf{behauptete}} (Prétérit, récit au passé).
    \item \all{\textbf{hatte}} (Prétérit de \all{haben}, subordonnée en \all{dass} implicite ou proposition principale coordonnée — ici Prétérit standard).
    \item \all{\textbf{glaubten}} (Prétérit, pluriel). \emph{Attention :} \all{ihm} (Datif) est implicite après \all{glauben}.
    \item \all{\textbf{teilten}} (Prétérit, pluriel, verbe fort \all{teilen}).
    \item \all{\textbf{Experte}} (Nominatif singulier, sujet de la phrase suivante).
    \item \all{\textbf{überprüfte}} (Prétérit, verbe \textbf{inséparable} \all{überprüfen} $\rightarrow$ \all{überprüfte}).
    \item \all{\textbf{die Quelle}} (Accusatif singulier féminin, COD de \all{überprüfte}).
    \item \all{\textbf{entdeckte}} (Prétérit, faible) ou \all{hat ... entdeckt} (Parfait à la fin de phrase). Ici Prétérit pour le style narratif.
    \item \all{\textbf{dass es sich um eine Falschmeldung handelte}} (Subordonnée \all{dass} + Prétérit de \all{handeln} + \all{sich um etw. (Akk) handeln}).
    \item \all{\textbf{Betrüger}} (Nominatif singulier, sujet).
    \item \all{\textbf{hatte ... zusammengestellt}} (Plus-que-parfait : \all{hatte} + Part. II \all{zusammengestellt}, verbe séparable \all{zusammenstellen}). Ou Prétérit \all{stellte ... zusammen}.
    \item \all{\textbf{musste}} (Prétérit de \all{müssen}, modal).
    \item \all{\textbf{sperren}} (Infinitif après modal \all{müssen}). Objet : \all{den Account} (Accusatif).
\end{enumerate}
\textbf{Traduction :} « L'expéditeur du message est resté anonyme. Il affirmait avoir des preuves d'un scandale. Beaucoup d'utilisateurs l'ont cru immédiatement et ont partagé le post des milliers de fois. Mais un expert a vérifié la source et a découvert qu'il s'agissait d'une fausse nouvelle. L'escroc avait monté de vieilles photos. À la fin, la plateforme a dû bloquer le compte. »
\end{exoresolu}

\begin{exercice}[Production écrite : Courte prise de position]
Rédigez en allemand un paragraphe de 6 à 8 lignes (environ 60--80 mots) sur le thème : \emph{„Wie kann man sich vor Fake News schützen?“} (Comment se protéger des fake news ?).
Utilisez au moins \textbf{5 noms} du lexique (champs 1 à 4), \textbf{3 verbes} du champ 3 (conjugués), et \textbf{une subordonnée en \all{dass}}.
\emph{Indications :} Pensez à : vérifier la source (\all{Quelle prüfen}), ne pas tout croire (\all{nicht alles glauben}), ne pas partager trop vite (\all{nicht sofort teilen}), signaler (\all{melden}), demander des preuves (\all{Beweise fordern}).
\end{exercice}

\begin{corrige}[Exercice : Production écrite]
\textbf{Proposition de correction (niveau B1/B2 visé) :}
\all{Um sich vor Fake News zu schützen, muss man kritisch sein. Man sollte nicht alles glauben, was man im Netz liest. Wichtig ist, die Quelle einer Nachricht zu überprüfen, bevor man sie teilt. Wenn ein Beitrag keine Beweise liefert, ist er oft eine Falschmeldung. Man muss auch darauf achten, wer der Absender ist: Ein seriöser Journalist nennt seine Quellen. Wenn man eine Lüge entdeckt, sollte man den Beitrag melden. So können wir alle die Verbreitung von Desinformation stoppen.}

\textbf{Analyse lexicale et grammaticale :}
\begin{itemize}
    \item Noms : \all{die Quelle, die Nachricht, der Beweis, die Falschmeldung, der Absender, der Journalist, die Quelle, die Lüge, der Beitrag, die Desinformation, die Verbreitung} (11 noms ciblés).
    \item Verbes conjugués : \all{muss ... sein, sollte ... glauben, liest, ist ... überprüfen, teilt, liefert, ist, muss ... achten, nennt, entdeckt, sollte ... melden, können ... stoppen}.
    \item Subordonnée \all{dass} : (ici structure infinitive \all{ohne ... zu} ou relative, mais on peut forcer \all{dass} : \all{Man muss sicherstellen, \textbf{dass} die Quelle stimmt.}).
\end{itemize}
\end{corrige}

\courssub{Synthèse visuelle : Carte mentale « Fake News »}

La figure ci-dessous résume les quatre champs lexicaux autour du concept central. Utilisez-la pour réviser ou pour préparer un exposé oral.

\begin{popfigure}
\begin{tikzpicture}[
    mindmap,
    concept color=popblue,
    level 1 concept/.append style={level distance=4.5cm, sibling angle=90, font=\large\bfseries},
    level 2 concept/.append style={level distance=3.2cm, sibling angle=45, font=\normalsize},
    every node/.style={concept, execute at begin node=\hskip0pt, align=center},
    text=white,
    grow cyclic
]
\node[concept, align=center]{Fake News\\ \all{Die Falschmeldung}}
    child[concept color=popgreen] {
        node[concept, align=center]{Wahrheit\\ \all{Information}}
        child { node[concept, scale=0.7] {\all{die Nachricht}} }
        child { node[concept, scale=0.7] {\all{die Wahrheit}} }
        child { node[concept, scale=0.7] {\all{die Tatsache}} }
        child { node[concept, scale=0.7] {\all{der Beweis}} }
        child { node[concept, scale=0.7] {\all{die Quelle}} }
    }
    child[concept color=poporange] {
        node[concept, align=center]{Lüge\\ \all{Manipulation}}
        child { node[concept, scale=0.7] {\all{die Lüge}} }
        child { node[concept, scale=0.7] {\all{die Falschmeldung}} }
        child { node[concept, scale=0.7] {\all{das Gerücht}} }
        child { node[concept, scale=0.7] {\all{die Propaganda}} }
        child { node[concept, scale=0.7] {\all{der Betrug}} }
    }
    child[concept color=poppink] {
        node[concept, align=center]{Handlungen\\ \all{Verben}}
        child { node[concept, scale=0.7] {\all{verbreiten}} }
        child { node[concept, scale=0.7] {\all{glauben}} }
        child { node[concept, scale=0.7] {\all{lügen}} }
        child { node[concept, scale=0.7] {\all{behaupten}} }
        child { node[concept, scale=0.7] {\all{überprüfen}} }
        child { node[concept, scale=0.7] {\all{melden}} }
    }
    child[concept color=poppurple] {
        node[concept, align=center]{Personen\\ \all{Akteure}}
        child { node[concept, scale=0.7] {\all{der Absender}} }
        child { node[concept, scale=0.7] {\all{der Journalist}} }
        child { node[concept, scale=0.7] {\all{der Experte}} }
        child { node[concept, scale=0.7] {\all{der Betrüger}} }
        child { node[concept, scale=0.7] {\all{der Nutzer}} }
    };
\end{tikzpicture}
\legende{Carte mentale du lexique « Fake News » : 4 branches (Vérité, Mensonge, Actions, Personnes) avec mots-clés allemands.}
\end{popfigure}

\begin{methode}[Comment réviser ce lexique efficacement ?]
\begin{enumerate}
    \item \textbf{Cartes mentales :} Reproduisez la carte ci-dessus à la main en ajoutant les articles (\all{der/die/das}) et les pluriels.
    \item \textbf{Familles de mots :} Pour chaque radical (\cle{wahr}, \cle{lügen}, \cle{glauben}, \cle{trügen}, \cle{prüfen}), écrivez 3 dérivés (nom, adjectif, verbe composé).
    \item \textbf{Phrases personnelles :} Écrivez 5 vraies phrases sur votre expérience (WhatsApp, Facebook, actualité camerounaise) en utilisant \all{verbreiten, glauben, überprüfen, melden, leugnen}.
    \item \textbf{Traduction inversée :} Prenez les 10 expressions de la boîte « Expressions utiles », cachez l'allemand, traduisez le français vers l'allemand à l'oral.
    \item \textbf{Attention aux verbes forts :} Récitez le triplet \all{lügen -- log -- gelogen} et \all{teilen -- teilte -- geteilt} avant de dormir.
\end{enumerate}
\end{methode}

\begin{attention}[Check-list des erreurs à ne plus commettre]
\begin{itemize}
    \item \textbf{Genre :} \all{die Lüge, die Falschmeldung, die Wahrheit, die Quelle, die Tatsache, die Information, die Nachricht} (tous \textbf{féminins} !). \all{der Beweis, der Betrug, der Absender, der Betrüger, der Experte, der Journalist, der Zeuge, der Nutzer} (tous \textbf{masculins} !). \all{das Gerücht} (neutre).
    \item \textbf{Pluriel :} \all{die Lügen, die Falschmeldungen, die Wahrheiten, die Tatsachen, die Beweise, die Quellen, die Gerüchte, die Informationen, die Nachrichten, die Experten, die Journalisten, die Betrüger, die Nutzer}.
    \item \textbf{Verbes forts :} \all{lügen (log, gelogen)}, \all{teilen (teilte, geteilt)}. Pas de \all{-te} au prétérit, pas de \all{-t} au participe (sauf \all{geteilt}).
    \item \textbf{Régime :} \all{glauben + DATIF} (\all{ihm glauben}), \all{täuschen + AKKUSATIF} (\all{mich täuschen}), \all{zugeben + AKKUSATIF / DASS-SATZ}.
    \item \textbf{Séparables / Inséparables :} \all{überprüfen} (inséparable \all{über-}), \all{zugeben} (séparable \all{zu-}), \all{verbreiten} (inséparable \all{ver-}), \all{zusammenstellen} (séparable \all{zusammen-}).
\end{itemize}
\end{attention}

\begin{savaistu}[Contexte camerounais : « Fake News » et élections]
Au Cameroun, comme en Allemagne ou en France, les périodes électorales voient exploser les \all{Falschmeldungen} sur WhatsApp et Facebook : fausses annonces de report du scrutin, rumeurs sur la santé des candidats, photos truquées de meetings. La loi n°2010/012 du 21 décembre 2010 relative à la cybercriminalité et la loi n°2019/020 du 24 décembre 2019 portant code pénal sanctionnent la \all{Verbreitung} de fausses nouvelles (\all{die Verbreitung falscher Nachrichten}) pouvant troubler l'ordre public. Le vocabulaire appris ici (\all{melden, überprüfen, die Quelle prüfen, der Betrüger}) est donc directement utilisable pour analyser l'actualité locale en allemand.
\end{savaistu}

\coursec{Grammaire 1 : le discours indirect (subjonctif I)}

Dans la section précédente, nous avons appris à nommer la vérité et le mensonge : \all{die Wahrheit}, \all{die Lüge}, \all{die Falschmeldung}, \all{verbreiten}, \all{glauben}. Mais comment rapporter les paroles de quelqu'un — d'un politicien, d'un témoin, d'un internaute — sans les affirmer soi-même ? Quand un journaliste allemand écrit qu'un ministre « aurait menti », il utilise un outil grammatical précis : le \cle{discours indirect} au \cle{subjonctif I} (\all{der indirekte Rede im Konjunktiv I}). C'est la grammaire de la prudence journalistique, et c'est l'une des grandes nouveautés de la classe de Terminale.

\courssub{Observation : du discours direct au discours indirect}

Reprenons des phrases inspirées de notre texte support \emph{Fake News im Netz} et observons ce qui change quand on rapporte une parole.

\begin{exemplebox}[Observer : discours direct et discours indirect]
\begin{enumerate}
\item \all{Der Minister sagt: „Ich habe die Falschmeldung nicht verbreitet.“} « Le ministre dit : “Je n'ai pas diffusé la fausse information.” »\\
$\rightarrow$ \all{Der Minister sagt, er habe die Falschmeldung nicht verbreitet.} « Le ministre dit qu'il n'a pas diffusé la fausse information. »
\item \all{Die Zeugin erklärt: „Ich war nicht im Büro.“} « La témoin déclare : “Je n'étais pas au bureau.” »\\
$\rightarrow$ \all{Die Zeugin erklärt, sie sei nicht im Büro gewesen.} « La témoin déclare qu'elle n'était pas au bureau. »
\item \all{Der Blogger behauptet: „Ich werde die Wahrheit veröffentlichen.“} « Le blogueur affirme : “Je publierai la vérité.” »\\
$\rightarrow$ \all{Der Blogger behauptet, er werde die Wahrheit veröffentlichen.} « Le blogueur affirme qu'il publiera la vérité. »
\end{enumerate}
\end{exemplebox}

\begin{center}
\begin{tabularx}{\linewidth}{|X|X|X|}
\hline
\rowcolor{popblueL} Discours direct & Discours indirect & Transformation \\
\hline
\all{„Ich \textbf{habe} gelogen.“} & \all{Er sagte, er \textbf{habe} gelogen.} & indicatif $\rightarrow$ subjonctif I passé (forme identique pour \textit{er}, mais personne changée) \\
\hline
\all{„Ich \textbf{bin} unschuldig.“} & \all{Er sagt, er \textbf{sei} unschuldig.} & \all{ist} $\rightarrow$ \all{sei} (subjonctif I présent) \\
\hline
\all{„Ich \textbf{werde} sprechen.“} & \all{Sie sagt, sie \textbf{werde} sprechen.} & \all{wird} $\rightarrow$ \all{werde} (subjonctif I futur) \\
\hline
\end{tabularx}
\end{center}

\textbf{Que remarque-t-on ?} Au discours indirect : (1) les guillemets disparaissent ; (2) les pronoms s'adaptent (\all{ich} $\rightarrow$ \all{er/sie}) ; (3) surtout, le verbe rapporté prend une forme nouvelle : \all{sei}, \all{habe}, \all{werde}. C'est le \cle{subjonctif I}.

\courssub{La règle générale : pourquoi un subjonctif ?}

\begin{propriete}[Fonction du discours indirect]
Le \cle{discours indirect} (\all{die indirekte Rede}) rapporte la parole, la pensée ou l'affirmation de quelqu'un \textbf{sans la garantir}. Celui qui rapporte dit : « voici ce qu'on m'a dit — à vous de juger si c'est vrai ».
\begin{itemize}
\item En \textbf{allemand}, cette distance est marquée grammaticalement par le \cle{subjonctif I} : \all{Er sagt, er \textbf{sei} unschuldig.}
\item En \textbf{français}, on utilise le plus souvent l'\textbf{indicatif} (il dit qu'il \emph{est} innocent) ou, dans la presse, le \textbf{conditionnel journalistique} (il \emph{serait} innocent, selon ses déclarations).
\end{itemize}
Le subjonctif I est donc le temps de la \textbf{neutralité} : le journaliste rapporte sans prendre parti. C'est exactement l'outil dont on a besoin face aux fake news : distinguer le fait vérifié de la simple rumeur rapportée.
\end{propriete}

\begin{savaistu}[Le subjonctif I, la langue des journalistes]
Le subjonctif I est très vivant dans la presse germanophone. Ouvrez un article de la \emph{Süddeutschen Zeitung}, de la \emph{Neue Zürcher Zeitung} ou du site de la Deutsche Welle : vous y lirez des phrases comme \all{Der Sprecher erklärte, die Regierung \textbf{prüfe} die Vorwürfe.} („ Le porte-parole a déclaré que le gouvernement examinait les accusations. “). Grâce au subjonctif I, le journaliste montre qu'il rapporte des propos sans les affirmer lui-même. C'est une marque d'honnêteté professionnelle — et une arme contre la désinformation. À l'oral, les Allemands utilisent plus souvent le subjonctif II (\all{er hätte gesagt...}), mais à l'écrit journalistique, le subjonctif I reste la norme.
\end{savaistu}

\courssub{Formation du subjonctif I présent}

\begin{propriete}[Formation : infinitif sans -n + désinences]
On prend le \textbf{radical de l'infinitif} (l'infinitif sans le \all{-n} final : \all{sage-}, \all{habe-}, \all{komme-}) et on ajoute les désinences :

\begin{center}
\begin{tabular}{|l|l|l|l|}
\hline
\rowcolor{popblueL} Personne & Désinence & \all{sagen} & \all{glauben} \\
\hline
ich & -e & \all{ich sage} & \all{ich glaube} \\
\hline
du & -est & \all{du sagest} & \all{du glaubest} \\
\hline
er/sie/es & -e & \all{er sage} & \all{er glaube} \\
\hline
wir & -en & \all{wir sagen} & \all{wir glauben} \\
\hline
ihr & -et & \all{ihr saget} & \all{ihr glaubet} \\
\hline
sie/Sie & -en & \all{sie sagen} & \all{sie glauben} \\
\hline
\end{tabular}
\end{center}

\textbf{Pas d'irrégularité de radical !} Même les verbes forts restent réguliers au subjonctif I présent : \all{gehen} $\rightarrow$ \all{er gehe} (jamais d'Umlaut ni de changement de voyelle), \all{kommen} $\rightarrow$ \all{er komme}, \all{wissen} $\rightarrow$ \all{er wisse}.
\end{propriete}

\begin{propriete}[Le seul verbe irrégulier : sein]
Seul le verbe \all{sein} échappe à la règle : il est formé sur le radical \all{sei-} et \textbf{sans -e} à la 1\textsuperscript{re} et à la 3\textsuperscript{e} personne du singulier :

\begin{center}
\begin{tabular}{|l|l|}
\hline
\rowcolor{popblueL} Personne & \all{sein} au subjonctif I \\
\hline
ich & \all{sei} \\
\hline
du & \all{seist} (ou \all{seiest}) \\
\hline
er/sie/es & \all{sei} \\
\hline
wir & \all{seien} \\
\hline
ihr & \all{seiet} \\
\hline
sie/Sie & \all{seien} \\
\hline
\end{tabular}
\end{center}

\all{Er sagt, er \textbf{sei} unschuldig.} « Il dit qu'il est innocent. »
\end{propriete}

\begin{exemplebox}[Exemples au subjonctif I présent]
\begin{itemize}
\item \all{Er sagt: „Ich \textbf{bin} unschuldig.“} $\rightarrow$ \all{Er sagt, er \textbf{sei} unschuldig.} « Il dit : “Je suis innocent.” $\rightarrow$ Il dit qu'il est innocent. »
\item \all{Sie behauptet: „Die Nachricht \textbf{stimmt}.“} $\rightarrow$ \all{Sie behauptet, die Nachricht \textbf{stimme}.} « Elle affirme que l'information est exacte. »
\item \all{Der Zeuge erklärt: „Ich \textbf{habe} nichts gesehen.“} $\rightarrow$ \all{Der Zeuge erklärt, er \textbf{habe} nichts gesehen.} « Le témoin déclare qu'il n'a rien vu. »
\item \all{Die Jugendlichen meinen: „Wir \textbf{glauben} nicht alles im Netz.“} $\rightarrow$ \all{Die Jugendlichen meinen, sie \textbf{glaubten} nicht alles im Netz.} « Les jeunes pensent qu'ils ne croient pas tout sur Internet. » (subjonctif II de substitution, voir plus bas)
\end{itemize}
\end{exemplebox}

\courssub{Le subjonctif I passé et futur}

\begin{propriete}[Subjonctif I passé : habe/sei + participe passé]
Pour rapporter un fait \textbf{passé}, on utilise l'auxiliaire \all{haben} ou \all{sein} au \textbf{subjonctif I présent} + le \textbf{participe passé} du verbe. Le choix de l'auxiliaire suit les mêmes règles qu'au Perfekt (\all{sein} pour les verbes de mouvement et de changement d'état) :
\begin{itemize}
\item \all{Er sagte, er \textbf{habe} gelogen.} « Il a dit qu'il avait menti. »
\item \all{Sie berichtete, sie \textbf{sei} spät gekommen.} « Elle a raconté qu'elle était arrivée tard. »
\item \all{Der Hacker gab zu, er \textbf{habe} die Daten gestohlen.} « Le pirate a admis qu'il avait volé les données. »
\end{itemize}
\end{propriete}

\begin{propriete}[Subjonctif I futur : werde + infinitif]
Pour rapporter un fait \textbf{futur}, on utilise \all{werden} au subjonctif I + l'infinitif :
\begin{itemize}
\item \all{Der Minister erklärte, er \textbf{werde} die Wahrheit veröffentlichen.} « Le ministre a déclaré qu'il publierait la vérité. »
\item \all{Sie sagte, sie \textbf{werde} die Falschmeldung löschen.} « Elle a dit qu'elle supprimerait la fausse information. »
\end{itemize}
\end{propriete}

\begin{center}
\begin{tabularx}{\linewidth}{|X|X|X|}
\hline
\rowcolor{popblueL} Temps & Formation & Exemple (3\textsuperscript{e} pers. sg.) \\
\hline
Présent & radical + -e/-est/-e/-en/-et/-en & \all{er \textbf{komme}} \\
\hline
Passé & \all{habe}/\all{sei} + participe passé & \all{er \textbf{habe gelogen}} / \all{er \textbf{sei gekommen}} \\
\hline
Futur & \all{werde} + infinitif & \all{er \textbf{werde kommen}} \\
\hline
\end{tabularx}
\end{center}

\begin{popfigure}
\begin{center}
\begin{tikzpicture}[font=\small]
% Tableau de conjugaison : subjonctif I aux trois temps
\node[draw=popblue, fill=popblueL, rounded corners, align=center, minimum width=4.2cm, minimum height=0.9cm] (titre) at (0,0) {\textbf{Subjonctif I : trois temps}};
\node[draw=popgreen, fill=popgreenL, rounded corners, align=center, minimum width=4.6cm, text width=4.4cm] (pres) at (-5.4,-2.2) {\textbf{Présent}\\er \textbf{sage}\\er \textbf{habe}\\er \textbf{sei}\\er \textbf{werde}};
\node[draw=poporange, fill=poporangeL, rounded corners, align=center, minimum width=4.6cm, text width=4.4cm] (past) at (0,-2.2) {\textbf{Passé}\\er \textbf{habe gesagt}\\er \textbf{habe gehabt}\\er \textbf{sei gewesen}\\er \textbf{sei geworden}};
\node[draw=poppurple, fill=poppurpleL, rounded corners, align=center, minimum width=4.6cm, text width=4.4cm] (fut) at (5.4,-2.2) {\textbf{Futur}\\er \textbf{werde sagen}\\er \textbf{werde haben}\\er \textbf{werde sein}\\er \textbf{werde werden}};
\draw[-{Stealth}, thick, popblue] (titre) -- (pres);
\draw[-{Stealth}, thick, popblue] (titre) -- (past);
\draw[-{Stealth}, thick, popblue] (titre) -- (fut);
\end{tikzpicture}
\legende{Le subjonctif I de \all{sagen}, \all{haben}, \all{sein}, \all{werden} aux trois temps (3\textsuperscript{e} personne du singulier).}
\end{center}
\end{popfigure}

\courssub{La règle de substitution : quand le subjonctif I ressemble à l'indicatif}

\begin{propriete}[Substitution par le subjonctif II]
Quand la forme du subjonctif I est \textbf{identique à l'indicatif présent}, on ne perçoit plus la distance du discours rapporté. On utilise alors le \cle{subjonctif II} à sa place :
\begin{itemize}
\item \all{ich sage} (subjonctif I = indicatif) $\rightarrow$ \all{Er sagte, ich \textbf{sähe} das anders.} (subjonctif II de \all{sehen})
\item \all{wir sagen} (identique à l'indicatif) $\rightarrow$ \all{Sie meinten, wir \textbf{kämen} zu spät.} « Ils pensaient que nous arrivions trop tard. »
\item \all{sie haben} (identique) $\rightarrow$ \all{Er behauptete, sie \textbf{hätten} die Nachricht verbreitet.} « Il affirmait qu'ils avaient diffusé l'information. »
\end{itemize}
En pratique, la substitution concerne surtout \textbf{ich}, \textbf{wir} et \textbf{sie} (pluriel). À la 3\textsuperscript{e} personne du singulier (\all{er sage}, \all{sie sei}), le subjonctif I est toujours distinct de l'indicatif, donc aucune substitution n'est nécessaire.
\end{propriete}

\begin{exemplebox}[Subjonctif I ou subjonctif II de substitution ?]
\begin{itemize}
\item \all{Er sagt, er \textbf{habe} keine Zeit.} (subjonctif I : \all{habe} $\neq$ indicatif \all{hat}) « Il dit qu'il n'a pas le temps. »
\item \all{Sie sagen, sie \textbf{hätten} keine Zeit.} (subjonctif II : \all{haben} serait identique à l'indicatif) « Ils disent qu'ils n'ont pas le temps. »
\item \all{Die Schüler berichten, sie \textbf{seien} müde.} (subjonctif I de \all{sein}, toujours distinct) « Les élèves racontent qu'ils sont fatigués. »
\end{itemize}
\end{exemplebox}

\courssub{Emploi : les verbes introducteurs, avec ou sans « dass »}

\begin{propriete}[Verbes de parole et de pensée + subjonctif I]
Le discours indirect est introduit par des verbes comme \all{sagen} (dire), \all{behaupten} (affirmer, prétendre), \all{meinen} (penser, estimer), \all{berichten} (rapporter), \all{erklären} (déclarer), \all{fragen} (demander), \all{antworten} (répondre), \all{glauben} (croire), \all{hinzufügen} (ajouter). Deux constructions sont possibles :
\begin{enumerate}
\item \textbf{Avec \all{dass}} : le verbe rapporté se place en \textbf{fin} de subordonnée : \all{Er sagt, \textbf{dass} er unschuldig \textbf{sei}.}
\item \textbf{Sans \all{dass}} : le verbe rapporté se place en \textbf{deuxième position}, comme dans une principale : \all{Er sagt, er \textbf{sei} unschuldig.}
\end{enumerate}
Les deux constructions sont correctes ; la construction sans \all{dass} est très fréquente dans la presse. Pour une question rapportée, on utilise \all{ob} (si) ou un mot interrogatif : \all{Sie fragte, \textbf{ob} er die Nachricht gesehen \textbf{habe}.} « Elle demanda s'il avait vu l'information. »
\end{propriete}

\begin{popfigure}
\begin{center}
\begin{tikzpicture}[font=\small]
\node[draw=popgreen, fill=popgreenL, rounded corners, align=center, text width=5.2cm, minimum height=1.6cm] (direct) at (-4.6,0) {\textbf{Discours direct}\\„Ich \textbf{habe} nicht\\gelogen“, sagte er.};
\node[draw=poporange, fill=poporangeL, rounded corners, align=center, text width=5.2cm, minimum height=1.6cm] (indirect) at (4.6,0) {\textbf{Discours indirect}\\Er sagte, er \textbf{habe}\\nicht gelogen.};
\draw[-{Stealth}, very thick, popblue] (direct) -- (indirect) node[midway, above, align=center, text width=3.4cm] {1. guillemets supprimés\\2. \all{ich} $\rightarrow$ \all{er}\\3. indicatif $\rightarrow$ Konj. I};
\end{tikzpicture}
\legende{Transformation du discours direct en discours indirect.}
\end{center}
\end{popfigure}

\courssub{Mini-texte d'application : la rumeur au lycée}

\begin{textebox}[Ein Gerücht am Gymnasium — Une rumeur au lycée]
Am Montagmorgen verbreitete sich am Gymnasium in Yaoundé eine seltsame Nachricht: Das Fußballturnier \textbf{sei} abgesagt, \textbf{behaupteten} einige Schüler. Der Direktor \textbf{habe} es persönlich angekündigt, \textbf{sagten} sie. Ein Schüler namens Eric \textbf{erklärte}, er \textbf{habe} die Information auf einer Internetseite gelesen. Die Seite \textbf{gehöre} einem bekannten Journalisten, \textbf{meinte} er, deshalb \textbf{sei} die Nachricht sicher wahr.

Die Sportlehrerin, Frau Mbarga, \textbf{fragte} sofort nach: Wer \textbf{habe} diese Nachricht geschrieben? Gebe es einen \textbf{Beweis}? Eric \textbf{antwortete}, er \textbf{wisse} das nicht genau, aber viele Schüler \textbf{hätten} die Nachricht schon geteilt.

Am Nachmittag stellte sich heraus: Alles war eine \textbf{Falschmeldung}. Der Direktor \textbf{erklärte} vor der ganzen Schule, er \textbf{habe} nie etwas abgesagt. Das Turnier \textbf{werde} wie geplant stattfinden. Er \textbf{fügte hinzu}, man \textbf{müsse} immer die Quelle prüfen, bevor man eine Nachricht \textbf{glaube} und \textbf{verbreite}.

Eric \textbf{gab zu}, er \textbf{sei} zu schnell gewesen. Er \textbf{habe} die Überschrift gelesen, aber nicht den ganzen Artikel. Die Schüler \textbf{lernten} an diesem Tag eine wichtige Lektion: Wer eine Nachricht teilt, ohne sie zu prüfen, \textbf{verbreite} vielleicht eine \textbf{Lüge} — auch wenn er die \textbf{Wahrheit} sagen will.

\emph{Glossaire :} das Gerücht, -e (la rumeur) ; absagen (annuler) ; der Beweis, -e (la preuve) ; herausstellen (sich) (se révéler) ; die Quelle, -n (la source) ; prüfen (vérifier) ; zugeben (gab zu, zugegeben) (admettre) ; die Überschrift, -en (le titre) ; teilen (partager).
\end{textebox}

\textbf{Questions de compréhension :}
\begin{enumerate}
\item \all{Was behaupteten einige Schüler am Montagmorgen?}
\item \all{Woher hatte Eric die Information angeblich?}
\item \all{Warum war die Nachricht eine Falschmeldung?}
\item \all{Welche Lektion lernten die Schüler an diesem Tag?}
\end{enumerate}

\courssub{Méthode : mettre une phrase au discours indirect}

\begin{methode}[Du discours direct au discours indirect en 4 étapes]
\begin{enumerate}
\item \textbf{Repérer le verbe de parole} (\all{sagen}, \all{behaupten}, \all{erklären}...) et le locuteur : c'est lui qui détermine le nouveau pronom.
\item \textbf{Choisir la construction} : avec \all{dass} (verbe en fin de phrase) ou sans \all{dass} (verbe en 2\textsuperscript{e} position).
\item \textbf{Mettre le verbe rapporté au subjonctif I} : présent (\all{er sage}), passé (\all{er habe gesagt} / \all{er sei gekommen}) ou futur (\all{er werde sagen}) selon le sens. Si la forme est identique à l'indicatif, passer au subjonctif II (\all{sie hätten}).
\item \textbf{Adapter les pronoms et les possessifs} : \all{ich} $\rightarrow$ \all{er/sie}, \all{mein} $\rightarrow$ \all{sein/ihr}, \all{wir} $\rightarrow$ \all{sie}.
\end{enumerate}
\end{methode}

\begin{exoresolu}[Mettre au discours indirect]
\textbf{Énoncé.} Mettez les phrases suivantes au discours indirect, avec ou sans \all{dass} :
\begin{enumerate}
\item \all{Der Zeuge sagt: „Ich war am Montag in Douala.“}
\item \all{Die Schülerin behauptet: „Ich habe die Nachricht nicht verbreitet.“}
\item \all{Der Minister erklärt: „Wir werden die Wahrheit suchen.“}
\item \all{Die Jugendlichen meinen: „Wir glauben den Fake News nicht.“}
\end{enumerate}
\tcblower
\textbf{Solution détaillée.}
\begin{enumerate}
\item Verbe de parole : \all{sagen} ; locuteur : \all{der Zeuge} $\rightarrow$ \all{er}. Fait passé (\all{war}) $\rightarrow$ subjonctif I passé avec \all{sein} : \all{sei gewesen}.\\
\all{Der Zeuge sagt, er \textbf{sei} am Montag in Douala \textbf{gewesen}.} « Le témoin dit qu'il était à Douala lundi. »
\item Fait passé avec \all{haben} : \all{habe verbreitet}.\\
\all{Die Schülerin behauptet, sie \textbf{habe} die Nachricht nicht \textbf{verbreitet}.} « L'élève affirme qu'elle n'a pas diffusé l'information. »
\item Fait futur : \all{werde} + infinitif ; \all{wir} $\rightarrow$ \all{sie}. Subjonctif I \all{werden} (3\textsuperscript{e} pl. \all{sie werden}) = indicatif $\rightarrow$ substitution subjonctif II : \all{würden}.\\
\all{Der Minister erklärt, sie \textbf{würden} die Wahrheit suchen.} « Le ministre déclare qu'ils chercheront la vérité. »
\item \all{wir} $\rightarrow$ \all{sie} ; \all{glauben} au subjonctif I (3\textsuperscript{e} pl. \all{sie glauben}) = indicatif $\rightarrow$ subjonctif II : \all{glaubten}.\\
\all{Die Jugendlichen meinen, sie \textbf{glaubten} den Fake News nicht.} « Les jeunes estiment qu'ils ne croient pas aux fausses informations. »
\end{enumerate}
\end{exoresolu}

\courssub{Comparaison français / allemand et pièges à éviter}

\begin{center}
\begin{tabularx}{\linewidth}{|X|X|X|}
\hline
\rowcolor{popblueL} Français & Deutsch & Remarque \\
\hline
Il dit : « Je suis innocent. » & \all{Er sagt: „Ich bin unschuldig.“} & discours direct : identique dans les deux langues \\
\hline
Il dit qu'il \textbf{est} innocent. & \all{Er sagt, er \textbf{sei} unschuldig.} & français : indicatif ; allemand : subjonctif I \\
\hline
Il \textbf{serait} innocent (selon lui). & \all{Er \textbf{sei} unschuldig, behauptet er.} & le conditionnel journalistique français $\approx$ subjonctif I \\
\hline
Elle a dit qu'elle \textbf{avait} menti. & \all{Sie sagte, sie \textbf{habe} gelogen.} & passé : \all{habe} + participe \\
\hline
Il a déclaré qu'il \textbf{viendrait}. & \all{Er erklärte, er \textbf{werde} kommen.} & futur : \all{werde} + infinitif \\
\hline
\end{tabularx}
\end{center}

\begin{attention}[Erreurs fréquentes des francophones]
\begin{itemize}
\item \textbf{Ne traduisez pas systématiquement le subjonctif I par un conditionnel français.} \all{Er sagt, er sei krank} se traduit le plus souvent par « Il dit qu'il \emph{est} malade » (indicatif), et non « qu'il serait malade ». Le conditionnel français (il \emph{serait} malade) n'est correct que dans le style journalistique, quand on veut marquer le doute.
\item \textbf{N'oubliez pas la place du verbe.} Avec \all{dass}, le verbe conjugué va en \textbf{fin} de phrase : \all{Er sagt, dass er unschuldig \textbf{sei}.} Sans \all{dass}, il reste en 2\textsuperscript{e} position : \all{Er sagt, er \textbf{sei} unschuldig.}
\item \textbf{Attention à \all{sein}.} On écrit \all{er sei} (jamais \all{er seie}). C'est la forme la plus fréquente du subjonctif I dans la presse : apprenez-la par cœur.
\item \textbf{Adaptez les pronoms.} Le français et l'allemand exigent tous deux la transformation \all{ich} $\rightarrow$ \all{er/sie} ; un oubli change complètement le sens de la phrase rapportée.
\item \textbf{Faux ami.} \all{behaupten} ne signifie pas « biaiser » mais « affirmer, prétendre » — souvent avec une nuance de doute, idéal pour parler des fake news : \all{Er behauptet, die Wahrheit zu kennen.} « Il prétend connaître la vérité. »
\end{itemize}
\end{attention}

\begin{aretenir}[L'essentiel du discours indirect]
\begin{itemize}
\item Le discours indirect rapporte une parole \textbf{sans la garantir} ; l'allemand la marque par le \textbf{subjonctif I}, le français par l'indicatif ou le conditionnel journalistique.
\item Formation du présent : infinitif sans \all{-n} + \textbf{-e, -est, -e, -en, -et, -en}. Seul \all{sein} est irrégulier : \all{ich sei, du seist, er sei, wir seien, ihr seiet, sie seien}.
\item Passé : \all{habe}/\all{sei} + participe passé ; futur : \all{werde} + infinitif.
\item Si le subjonctif I est identique à l'indicatif, on le remplace par le \textbf{subjonctif II} (\all{sie hätten}, \all{wir kämen}).
\item Construction avec \all{dass} (verbe en fin) ou sans \all{dass} (verbe en 2\textsuperscript{e} position), après \all{sagen}, \all{behaupten}, \all{meinen}, \all{berichten}, \all{erklären}, \all{fragen}.
\end{itemize}
\end{aretenir}

\begin{exercice}[À vous de jouer]
\begin{enumerate}
\item Conjuguez au subjonctif I présent (en appliquant la règle de substitution si nécessaire) : \all{verbreiten} (er), \all{glauben} (sie, pluriel), \all{haben} (ich), \all{sein} (wir), \all{kommen} (du).
\item Mettez au discours indirect : \all{Der Blogger sagt: „Ich werde den Artikel löschen.“} / \all{Die Zeugin erklärt: „Ich habe den Beweis gefunden.“}
\item Traduisez en allemand : « Le journaliste rapporte que le ministre a menti. » / « Elle affirme qu'elle est innocente. »
\end{enumerate}
\end{exercice}

\begin{corrige}[Exercice « À vous de jouer »]
\begin{enumerate}
\item \all{er verbreite} ; \all{sie glaubten} (subjonctif II de substitution, car \all{glauben} = indicatif) ; \all{ich hätte} (subjonctif II de substitution, car \all{habe} = indicatif) ; \all{wir seien} ; \all{du kommest}.
\item \all{Der Blogger sagt, er \textbf{werde} den Artikel löschen.} « Le blogueur dit qu'il supprimera l'article. » / \all{Die Zeugin erklärt, sie \textbf{habe} den Beweis gefunden.} « La témoin déclare qu'elle a trouvé la preuve. »
\item \all{Der Journalist berichtet, der Minister \textbf{habe} gelogen.} / \all{Sie behauptet, sie \textbf{sei} unschuldig.}
\end{enumerate}
\end{corrige}

Nous savons désormais rapporter les paroles avec la distance du subjonctif I. Reste à maîtriser la construction qui accompagne le plus souvent ce discours rapporté : la subordonnée introduite par \all{dass}, que nous étudierons en détail dans la section suivante.

\coursec{Grammaire 2 : la subordonnée en « dass »}

Dans la section précédente, nous avons appris à rapporter les paroles de quelqu'un au discours indirect grâce au \cle{subjonctif I} : \all{Der Experte sagte, er sei unschuldig.} Mais cette construction sans conjonction n'est pas la seule possibilité. L'allemand, comme le français, dispose d'une conjonction équivalente à « que » : \cle{dass}. C'est elle qui va nous occuper maintenant, car elle entraîne une conséquence capitale pour l'ordre des mots : le verbe conjugué part se placer tout à la fin de la phrase. C'est l'une des structures les plus fréquentes de l'allemand écrit et parlé — et l'une des plus sanctionnées au baccalauréat.

\courssub{Observation : que remarque-t-on ?}

Reprenons une phrase de notre texte support et observons-la attentivement :

\all{Der Experte sagte, dass viele Jugendliche alles glaubten.} « L'expert a dit que beaucoup de jeunes croyaient tout. »

Comparons avec la version sans conjonction étudiée à la section 4 :

\begin{center}
\begin{tabularx}{\linewidth}{|X|X|}
\hline
\rowcolor{popblueL} \textbf{Sans dass (verbe en 2\up{e} position)} & \textbf{Avec dass (verbe en fin)} \\
\hline
\all{Er sagte, er sei unschuldig.} & \all{Er sagte, dass er unschuldig sei.} \\
\hline
\all{Sie glaubt, die Quelle sei zuverlässig.} & \all{Sie glaubt, dass die Quelle zuverlässig sei.} \\
\hline
\end{tabularx}
\end{center}

\begin{exemplebox}[Ce que l'on observe]
\begin{enumerate}
\item Une \cle{virgule} sépare toujours la proposition principale de la proposition introduite par \all{dass}.
\item Après \all{dass}, on trouve d'abord le \cle{sujet} (\all{viele Jugendliche}), puis les compléments (\all{alles}), et le \cle{verbe conjugué} (\all{glaubten}) se trouve \textbf{en toute dernière position}.
\item Le verbe de la subordonnée peut être au subjonctif I (discours indirect soutenu) ou à l'indicatif (style courant, oral) : \all{Er sagt, dass er unschuldig ist.} / \all{Er sagt, dass er unschuldig sei.}
\end{enumerate}
\end{exemplebox}

\courssub{La règle : forme et emploi}

\begin{propriete}[Règle de la subordonnée en « dass »]
\all{dass} est une \cle{conjonction de subordination} qui introduit une \textbf{proposition complétive} (objet du verbe de la principale). Sa structure est rigide :

\begin{center}
\textbf{[Proposition principale] , dass + Sujet + Compléments + Verbe conjugué (position finale)}
\end{center}

\begin{itemize}
\item \textbf{Virgule obligatoire} : en allemand, on met \emph{toujours} une virgule avant \all{dass}, contrairement au français où « que » n'est jamais précédé d'une virgule.
\item \textbf{Verbe en fin} : le verbe conjugué occupe la dernière place de la subordonnée. S'il y a un auxiliaire et un participe ou un infinitif, l'ordre est : participe/infinitif puis auxiliaire conjugué : \all{..., dass er die Nachricht gelesen hat.}
\item \textbf{Verbes introducteurs} : \all{dass} s'emploie après les verbes de parole, d'opinion, de sentiment et de connaissance : \all{sagen} (dire), \all{glauben} (croire), \all{denken} (penser), \all{hoffen} (espérer), \all{wissen} (savoir), \all{behaupten} (affirmer), \all{meinen} (estimer), \all{berichten} (rapporter), \all{befürchten} (craindre), \all{vermuten} (supposer).
\item \textbf{Mode} : à l'oral et dans le style courant, l'indicatif est normal ; au discours indirect écrit et soutenu, on emploie le subjonctif I (voir section 4). Lorsque la forme du subjonctif I est identique à celle de l'indicatif (surtout au pluriel), on utilise la \textbf{forme de remplacement} du subjonctif II : \all{..., dass sie es glaubten} (au lieu de \all{glauben}, identique à l'indicatif).
\end{itemize}
\end{propriete}

\begin{popfigure}
\begin{center}
\begin{tikzpicture}[
  box/.style={draw=popblue, fill=popblueL, rounded corners=2pt, align=center, font=\small, inner sep=4pt},
  vbox/.style={draw=poporange, fill=poporangeL, rounded corners=2pt, align=center, font=\small, inner sep=4pt},
  cbox/.style={draw=popgreen, fill=popgreenL, rounded corners=2pt, align=center, font=\small, inner sep=4pt},
  arr/.style={-{Stealth[length=2.5mm]}, thick, popmuted}
]
\node[box, align=center] (main) at (0,0){\textbf{Proposition principale}\\Der Experte sagte};
\node[cbox] (virg) at (4.1,0) {\textbf{,}};
\node[box] (dass) at (5.3,0) {\textbf{dass}};
\node[box, align=center] (subj) at (7.0,0){\textbf{Sujet}\\viele Jugendliche};
\node[box, align=center] (compl) at (9.3,0){\textbf{Compléments}\\alles};
\node[vbox, align=center] (verb) at (11.6,0){\textbf{Verbe conjugué}\\glaubten};
\draw[arr] (main) -- (virg);
\draw[arr] (virg) -- (dass);
\draw[arr] (dass) -- (subj);
\draw[arr] (subj) -- (compl);
\draw[arr] (compl) -- (verb);
\node[align=center, font=\small\itshape, popdark] at (11.6,-1.1) {position finale !};
\draw[arr, poporange] (11.6,-0.9) -- (verb.south);
\end{tikzpicture}
\end{center}
\legende{Structure de la subordonnée en « dass » : le verbe conjugué ferme la phrase.}
\end{popfigure}

\courssub{Comparaison français / allemand : deux logiques opposées}

C'est ici que le francophone doit faire une véritable « gymnastique mentale ». En français, le verbe suit immédiatement le sujet ; en allemand, dans une subordonnée en \all{dass}, il faut \textbf{retenir le verbe} et ne le lâcher qu'à la toute fin.

\begin{center}
\begin{tabularx}{\linewidth}{|X|X|X|}
\hline
\rowcolor{popblueL} \textbf{Français} & \textbf{Deutsch} & \textbf{Remarque} \\
\hline
Je crois qu'il dit la vérité. & \all{Ich glaube, dass er die Wahrheit sagt.} & verbe \all{sagt} en fin \\
\hline
Elle espère que la source est fiable. & \all{Sie hofft, dass die Quelle zuverlässig ist.} & verbe \all{ist} en fin \\
\hline
Nous savons que beaucoup de jeunes partagent ces messages. & \all{Wir wissen, dass viele Jugendliche diese Nachrichten teilen.} & longue subordonnée : le verbe attend la fin \\
\hline
Il affirme qu'il a lu l'article hier. & \all{Er behauptet, dass er den Artikel gestern gelesen hat.} & auxiliaire \all{hat} après le participe \\
\hline
Je pense que tu dois vérifier l'information. & \all{Ich denke, dass du die Information prüfen musst.} & l'infinitif précède le modal conjugué \\
\hline
\end{tabularx}
\end{center}

\begin{attention}[La virgule : différence capitale]
En français, on écrit : « Je crois qu'il dit la vérité » — \textbf{sans virgule}. En allemand, la virgule avant \all{dass} est \textbf{obligatoire}, sans aucune exception : \all{Ich glaube, dass er die Wahrheit sagt.} Oublier cette virgule est une faute d'orthographe en allemand, pénalisée au baccalauréat. Inversement, ne mettez jamais de virgule après \all{dass}.
\end{attention}

\courssub{Discours indirect : avec ou sans « dass » ?}

Nous avons vu à la section 4 que le discours indirect peut se construire \textbf{sans} conjonction. Il existe donc deux options correctes pour rapporter une parole :

\begin{exemplebox}[Les deux constructions du discours indirect]
\textbf{Option 1 — avec \all{dass}} : le verbe conjugué va en position finale.

\all{Er sagt, dass er unschuldig sei.} « Il dit qu'il est innocent. »

\all{Die Zeitung berichtet, dass die Polizei den Täter gefasst habe.} « Le journal rapporte que la police aurait arrêté le coupable. »

\textbf{Option 2 — sans \all{dass}} : le verbe conjugué prend la place qu'occuperait \all{dass}, c'est-à-dire la deuxième position (construction de type principale).

\all{Er sagt, er sei unschuldig.}

\all{Die Zeitung berichtet, die Polizei habe den Täter gefasst.}
\end{exemplebox}

\begin{methode}[Comment traduire « il dit que... »]
\begin{enumerate}
\item \textbf{Identifier} le verbe introducteur (dire, croire, penser...) et choisir son équivalent allemand : \all{sagen, glauben, denken...}
\item \textbf{Choisir la construction} : avec \all{dass} (plus sûre pour les élèves, structure fixe) ou sans \all{dass} (plus élégante, style journalistique).
\item \textbf{Si l'option « dass » est choisie} : virgule, \all{dass}, sujet, compléments, puis placer le verbe conjugué à la toute fin, au subjonctif I si le contexte est un discours rapporté.
\item \textbf{Vérifier} : la virgule est-elle présente ? Le verbe est-il bien en dernière position ? Le mode (subjonctif I) est-il correct ?
\end{enumerate}
Au baccalauréat, \textbf{les deux options sont acceptées} : \all{Er sagt, dass er unschuldig sei.} et \all{Er sagt, er sei unschuldig.} valent le même point. Choisissez celle que vous maîtrisez le mieux.
\end{methode}

\courssub{Texte d'application : le discours rapporté dans la presse}

\begin{textebox}[Aus einer Schülerzeitung: „Gerüchte an unserer Schule“]
Letzte Woche verbreitete sich an unserer Schule ein Gerücht: Ein Schüler behauptete, dass die Prüfungen abgesagt würden. Viele glaubten sofort, dass die Nachricht stimme, und teilten sie in allen Gruppen. Eine Mitschülerin sagte, sie habe die Information von einer „sicheren Quelle“. Doch die Direktorin erklärte am Montag, dass das Gerücht falsch sei. Sie betonte, dass jede offizielle Mitteilung nur auf der Schulwebsite erscheine. Ein Lehrer meinte, dass solche Falschmeldungen gefährlich seien, weil sie Panik verbreiteten. Er riet uns, dass wir jede Nachricht zuerst prüfen sollten. Die Schülerzeitung berichtete daraufhin, dass die meisten Schüler aus diesem Fehler gelernt hätten. Ein Zehntklässler gab zu, dass er die Meldung weitergeleitet habe, ohne nachzudenken. Er sagte, er schäme sich dafür und werde künftig vorsichtiger sein. Die Direktorin hofft nun, dass die Schüler kritischer mit Informationen im Netz umgehen. Experten warnen, dass Fake News besonders schnell wirken, wenn sie unsere Wünsche bestätigen — zum Beispiel den Wunsch nach prüfungsfreien Tagen!
\end{textebox}

\begin{definition}[Glossaire du texte]
\begin{itemize}
\item \all{das Gerücht, -e} : la rumeur
\item \all{absagen} : annuler (verbe à particule séparable)
\item \all{stimmen} : être exact, être vrai
\item \all{die Mitteilung, -en} : la communication, l'annonce
\item \all{betonen} : souligner, insister
\item \all{raten (rät, riet, hat geraten) + Dat.} : conseiller (à quelqu'un)
\item \all{weiterleiten} : transférer, faire suivre
\item \all{zugeben (gibt zu, gab zu, hat zugegeben)} : avouer
\item \all{sich schämen (für + Akk.)} : avoir honte (de)
\item \all{umgehen mit + Dat.} : se comporter avec, gérer
\item \all{der Wunsch, „-e} : le souhait, le désir
\end{itemize}
\end{definition}

\textbf{Questions de compréhension et de grammaire :}
\begin{enumerate}
\item \all{Was behauptete ein Schüler?} — Relevez la subordonnée en \all{dass} et soulignez le verbe en position finale.
\item \all{Was erklärte die Direktorin am Montag?}
\item Trouvez dans le texte deux exemples de discours indirect \textbf{sans} \all{dass} et transformez-les avec \all{dass}.
\item Pourquoi, selon les experts, les fake news fonctionnent-elles si vite ?
\end{enumerate}

\begin{corrige}[Questions sur le texte]
\begin{enumerate}
\item \all{Ein Schüler behauptete, dass die Prüfungen abgesagt würden.} — Verbe en fin : \all{würden} (forme de remplacement du subjonctif II, car le subjonctif I \all{werden} serait identique à l'indicatif).
\item \all{Die Direktorin erklärte, dass das Gerücht falsch sei.}
\item \all{Sie sagte, sie habe die Information von einer „sicheren Quelle“.} $\rightarrow$ \all{Sie sagte, dass sie die Information von einer „sicheren Quelle“ habe.} ; \all{Er sagte, er schäme sich dafür.} $\rightarrow$ \all{Er sagte, dass er sich dafür schäme.}
\item \all{Fake News wirken besonders schnell, wenn sie unsere Wünsche bestätigen.} — Les fake news fonctionnent vite parce qu'elles confirment nos désirs (par exemple le désir de journées sans examen).
\end{enumerate}
\end{corrige}

\courssub{Exercice résolu}

\begin{exoresolu}[Construire des subordonnées en « dass »]
\textbf{Énoncé.} Transformez les phrases suivantes en introduisant une subordonnée en \all{dass} (discours indirect, subjonctif I).

\begin{enumerate}
\item \all{Der Experte sagt: „Die Meldung ist falsch.“}
\item \all{Die Schülerin glaubt: „Die Quelle ist zuverlässig.“}
\item \all{Wir hoffen: „Die Polizei findet den Täter.“}
\item \all{Er behauptet: „Ich habe den Artikel nicht gelesen.“}
\end{enumerate}

\tcblower
\textbf{Solution détaillée.}

\begin{enumerate}
\item \all{Der Experte sagt, dass die Meldung falsch sei.} — Virgule + \all{dass} + sujet (\all{die Meldung}) + adjectif (\all{falsch}) + verbe \all{sein} au subjonctif I (\all{sei}) en position finale.

\item \all{Die Schülerin glaubt, dass die Quelle zuverlässig sei.} — Même schéma ; \all{sei} reste en fin.

\item \all{Wir hoffen, dass die Polizei den Täter finde.} — \all{finden} au subjonctif I, 3\up{e} personne du singulier : \all{finde} (la Polizei est un nom singulier : \all{sie finde}).

\item \all{Er behauptet, dass er den Artikel nicht gelesen habe.} — Au passé (Perfekt du subjonctif I) : participe \all{gelesen} puis auxiliaire \all{habe} en toute fin. La négation \all{nicht} se place avant le participe.
\end{enumerate}

\textbf{Vérification} : virgule présente, verbe conjugué en dernière position, subjonctif I correct.
\end{exoresolu}

\begin{exercice}[À vous de jouer]
Traduisez en allemand avec une subordonnée en \all{dass} :
\begin{enumerate}
\item Je crois qu'il dit la vérité.
\item Elle espère que la source est fiable.
\item Le professeur dit que les élèves ont bien travaillé.
\item Nous savons que les fake news se répandent vite.
\end{enumerate}
\end{exercice}

\begin{corrige}[Exercice « À vous de jouer »]
\begin{enumerate}
\item \all{Ich glaube, dass er die Wahrheit sagt.} (ou discours indirect : \all{..., dass er die Wahrheit sage.})
\item \all{Sie hofft, dass die Quelle zuverlässig ist.}
\item \all{Der Lehrer sagt, dass die Schüler gut gearbeitet haben.} — auxiliaire \all{haben} après le participe \all{gearbeitet}.
\item \all{Wir wissen, dass sich Fake News schnell verbreiten.} — attention au verbe pronominal : \all{sich} suit immédiatement le sujet ; le verbe \all{verbreiten} reste en fin.
\end{enumerate}
\end{corrige}

\courssub{Piège à éviter : « dass » ou « das » ?}

\begin{attention}[dass (deux s) ou das (un s) ?]
Depuis la réforme orthographique, la conjonction s'écrit \textbf{toujours} avec deux s : \all{dass}. Le mot \all{das} avec un seul s est l'\textbf{article} neutre (\all{das Kind} = l'enfant) ou le \textbf{pronom démonstratif/relatif} (\all{Das ist mein Bruder.} = Voici/C'est mon frère).

\textbf{Le test infaillible} : si l'on peut remplacer le mot par \all{dieses} ou \all{welches}, on écrit \all{das} (un seul s) ; sinon, on écrit \all{dass}.

\begin{itemize}
\item \all{Ich sehe das (= dieses) Haus.} $\rightarrow$ article : un seul s.
\item \all{Das (= dieses) ist gefährlich.} $\rightarrow$ pronom : un seul s.
\item \all{Ich glaube, dass er lügt.} $\rightarrow$ impossible de dire \all{dieses/welches} : conjonction, deux s.
\end{itemize}

\textbf{Autre piège de francophone} : ne jamais traduire « que » par \all{das} ni oublier la virgule : « Je sais que c'est vrai » se traduit \all{Ich weiß, dass das wahr ist.} — notez les deux mots côte à côte : \all{dass} (conjonction) + \all{das} (pronom) !
\end{attention}

\begin{aretenir}[L'essentiel de la subordonnée en « dass »]
\begin{itemize}
\item \all{dass} = « que » ; il introduit une complétive après un verbe de parole, d'opinion ou de connaissance.
\item Schéma : \textbf{Principale , dass + sujet + compléments + verbe conjugué en fin}.
\item Virgule \textbf{obligatoire} avant \all{dass} (contrairement au français).
\item Discours indirect : deux options acceptées au Bac — \all{Er sagt, dass er unschuldig sei.} ou \all{Er sagt, er sei unschuldig.}
\item Si le subjonctif I est identique à l'indicatif (pluriel), on utilise la forme de remplacement du subjonctif II : \all{..., dass sie es glaubten.}
\item Orthographe : conjonction \all{dass} (deux s) ; article/pronom \all{das} (un s). Test : remplacement par \all{dieses/welches}.
\end{itemize}
\end{aretenir}

\coursec{Méthode du commentaire et commentaire modèle}

Nous savons désormais rapporter les propos d'autrui au discours indirect (Konjunktiv I) et construire des subordonnées en \all{dass}. Ces deux outils grammaticaux ne sont pas des exercices isolés : ils sont précisément les instruments du \cle{commentaire de texte}, épreuve reine de l'expression écrite au Baccalauréat. Cette dernière section montre comment les mobiliser pour analyser et résumer le texte étudié, « Fake News – die neue Gefahr im Internet ».

\courssub{Qu'est-ce qu'un commentaire de texte au Bac ?}

\begin{definition}[Le commentaire de texte]
Le \cle{commentaire de texte} (allemand : \all{die Textanalyse} ou \all{der Kommentar}) est une production écrite organisée dans laquelle l'élève présente un texte, en dégage le thème et les arguments, analyse les procédés de l'auteur, puis donne son avis personnel argumenté. Il ne s'agit ni d'une traduction, ni d'une simple paraphrase : il faut \emph{comprendre, expliquer et juger}.
\end{definition}

Au Baccalauréat, le correcteur attend trois qualités : une \cle{structure claire} (introduction, analyse, conclusion), une \cle{langue correcte} (place du verbe, discours indirect pour rapporter les propos de l'auteur) et un \cle{avis personnel} exprimé avec des marqueurs d'opinion.

\courssub{Les quatre étapes du commentaire}

\begin{methode}[Rédiger un commentaire de texte en quatre temps]
\begin{enumerate}
\item \textbf{Einleitung (introduction)} : présenter le texte en deux ou trois phrases — type de texte (\all{Zeitungsartikel}, \all{Interview}, \all{Bericht}), source, thème, et problématique. Formule utile : \all{Der Text ist ein Zeitungsartikel über ... Er behandelt das Problem ...} « Le texte est un article de journal sur... Il traite le problème... »
\item \textbf{Analyse} : dégager les idées principales dans l'ordre du texte, en rapportant les propos de l'auteur au \cle{subjonctif I} et en reliant les étapes par des connecteurs : \all{zuerst} (d'abord), \all{dann} (ensuite), \all{außerdem} (de plus), \all{schließlich} (enfin). Relever les procédés : exemples, citations d'experts, chiffres, questions rhétoriques.
\item \textbf{Meinung (avis personnel)} : donner son point de vue avec \all{meiner Meinung nach} (à mon avis), \all{ich finde, dass ...}, \all{ich bin der Meinung, dass ...}. Attention : après \all{dass}, le verbe conjugué va à la fin !
\item \textbf{Schluss (conclusion)} : répondre brièvement à la problématique et ouvrir sur un contexte plus large, par exemple la situation au Cameroun : \all{Auch in Kamerun ist dieses Problem aktuell.} « Au Cameroun aussi, ce problème est d'actualité. »
\end{enumerate}
\end{methode}

\begin{popfigure}
\begin{center}
\begin{tikzpicture}[node distance=0.55cm,
box/.style={rectangle, rounded corners=4pt, draw=popblue, fill=popblueL, text width=4.6cm, align=center, minimum height=1.0cm, font=\small},
arr/.style={-{Stealth[length=3mm]}, thick, popdark}]
\node[box, align=center] (e){\textbf{Einleitung}\\ Type de texte, source, thème, problématique};
\node[box, below=of e, align=center] (a){\textbf{Analyse}\\ Idées principales + Konjunktiv I\\ (zuerst, dann, außerdem, schließlich)};
\node[box, below=of a, draw=poporange, fill=poporangeL, align=center] (m){\textbf{Meinung}\\ Avis personnel\\ (meiner Meinung nach, ich finde, dass ...)};
\node[box, below=of m, draw=popgreen, fill=popgreenL, align=center] (s){\textbf{Schluss}\\ Réponse à la problématique + ouverture (Kamerun)};
\draw[arr] (e) -- (a);
\draw[arr] (a) -- (m);
\draw[arr] (m) -- (s);
\end{tikzpicture}
\end{center}
\legende{Organigramme du commentaire de texte : Einleitung → Analyse → Meinung → Schluss}
\end{popfigure}

\courssub{Les outils linguistiques du commentaire}

Le commentaire exige un vocabulaire précis pour présenter, rapporter et juger. Voici le tableau récapitulatif des expressions indispensables.

\begin{center}
\begin{tabularx}{\linewidth}{|X|X|X|}
\hline
\rowcolor{popblueL} Fonction & Deutsch & Français \\
\hline
Présenter le texte & \all{Der Text handelt von ...} & Le texte traite de... \\
\hline
Présenter le texte & \all{Das Thema ist ...} & Le thème est... \\
\hline
Rapporter (indirect) & \all{Der Autor sagt, Fake News seien gefährlich.} & L'auteur dit que les fake news sont dangereuses. \\
\hline
Rapporter (indirect) & \all{Er behauptet, sie verbreiteten sich schnell.} & Il affirme qu'elles se répandent vite. \\
\hline
Ordonner l'analyse & \all{zuerst – dann – außerdem – schließlich} & d'abord – ensuite – de plus – enfin \\
\hline
Donner son avis & \all{Meiner Meinung nach ...} & À mon avis... \\
\hline
Donner son avis & \all{Ich finde, dass das Problem wichtig ist.} & Je trouve que le problème est important. \\
\hline
Conclure & \all{Zusammenfassend kann man sagen, ...} & En résumé, on peut dire... \\
\hline
Ouvrir & \all{Auch in Kamerun ...} & Au Cameroun aussi... \\
\hline
\end{tabularx}
\end{center}

\begin{propriete}[Rapporter les propos de l'auteur : la règle d'or]
Dans l'analyse, on ne rapporte jamais les propos de l'auteur à l'indicatif : on utilise le \cle{Konjunktiv I}, seul ou introduit par \all{dass}.
\begin{itemize}
\item Sans \all{dass} (verbe en position 2) : \all{Der Autor sagt, viele Menschen informierten sich nur noch über soziale Medien.} « L'auteur dit que beaucoup de gens ne s'informent plus que par les réseaux sociaux. »
\item Avec \all{dass} (verbe à la fin) : \all{Der Autor erklärt, dass sich Falschmeldungen schnell verbreiteten.} « L'auteur explique que les fausses informations se répandent rapidement. »
\end{itemize}
\textbf{Forme de remplacement (Ersatzform)} : quand la forme du Konjunktiv I est identique à l'indicatif présent, on la remplace par le Konjunktiv II pour que le discours indirect reste reconnaissable. C'est le cas le plus fréquent au pluriel : \all{sie informieren} (indicatif = KI) → \all{sie informierten} ; \all{sie verbreiten} → \all{sie verbreiteten} ; \all{sie glauben} → \all{... dass viele Leser die Nachrichten glauben würden.} « ... que beaucoup de lecteurs croiraient les informations. » En revanche, à la 3\up{e} personne du singulier, le Konjunktiv I est toujours distinct (\all{er sagt} → \all{er sage}, \all{sie ist} → \all{sie sei}, \all{er muss} → \all{er müsse}) : on l'utilise donc tel quel.
\end{propriete}

\courssub{Commentaire modèle rédigé}

Lisons maintenant un commentaire complet (10 à 12 lignes) sur notre texte support. Observons comment les quatre étapes s'enchaînent et comment le subjonctif I rapporte fidèlement les propos de l'auteur.

\begin{textebox}[Kommentar : « Fake News – die neue Gefahr im Internet »]
Der Text ist ein Zeitungsartikel über Fake News im Internet. Er behandelt die Frage, warum sich Falschmeldungen so schnell verbreiten und wie man sie erkennen kann. Zuerst erklärt der Autor, dass sich viele Menschen nur noch über soziale Medien informierten. Dann zitiert er einen Experten, der sage, Falschmeldungen verbreiteten sich schneller als die Wahrheit. Außerdem behauptet ein Journalist, manche Fake News seien bezahlt worden. Schließlich rät der Autor, jede Nachricht zu überprüfen, bevor man sie weiterschickt. Meiner Meinung nach ist das ein sehr wichtiges Problem, auch in Kamerun, wo viele falsche Nachrichten über WhatsApp verbreitet werden. Ich finde, dass die Schulen die Jugendlichen besser informieren müssen. Zusammenfassend kann man sagen: Nur kritische Leser können die Wahrheit finden.
\end{textebox}

\textbf{Glossaire} : \all{sich verbreiten} (se répandre) ; \all{erkennen, erkannte, hat erkannt} (reconnaître) ; \all{weiterschicken} (transférer, faire suivre) ; \all{überprüfen} (vérifier) ; \all{bevor} (avant de, avant que — en allemand suivi de l'indicatif) ; \all{kritisch} (critique, qui a l'esprit critique).

\textbf{Analyse du modèle} : l'introduction (phrases 1 et 2) présente le type de texte, le thème et la problématique. L'analyse (connecteurs \all{zuerst, dann, außerdem, schließlich}) rapporte chaque argument au Konjunktiv I : \all{informierten} et \all{verbreiteten} (formes de remplacement Konjunktiv II, car le KI pluriel serait identique à l'indicatif), \all{sage} (KI singulier de \all{sagen}, distinct de l'indicatif \all{sagt}), \all{seien bezahlt worden} (Konjunktiv I du passif au passé : « auraient été payées »). L'avis personnel est marqué par \all{meiner Meinung nach} et \all{ich finde, dass ...}. La conclusion répond à la problématique et ouvre sur le Cameroun.

\begin{attention}[Les erreurs à éviter dans le résumé et le commentaire]
\begin{itemize}
\item \textbf{Ne pas recopier le texte.} Un résumé n'est pas un collage de phrases du texte : il faut \emph{reformuler} avec ses propres mots et rapporter les propos au discours indirect. Recopier trois phrases du texte est sanctionné.
\item \textbf{Place du verbe.} Après \all{meiner Meinung nach}, le verbe reste en position 2 : \all{Meiner Meinung nach \textbf{ist} das Problem wichtig.} et non \emph{Meiner Meinung nach das Problem ist wichtig}.
\item \textbf{Subordonnée en dass.} Le verbe conjugué va à la fin : \all{Ich finde, dass Fake News gefährlich \textbf{sind}.}
\item \textbf{Discours indirect.} Ne pas mélanger indicatif et Konjunktiv I dans le même résumé : une fois le discours indirect choisi, on s'y tient.
\item \textbf{Faux ami.} \all{kritisch} signifie « qui a l'esprit critique » ; « une critique » au sens de « texte d'analyse » se dit \all{die Kritik} ou \all{der Kommentar}. De même, \all{die Lüge} = « le mensonge », mais \all{lügen} (verbe fort : \all{lügt, log, hat gelogen}) = « mentir ».
\item \textbf{Majuscules.} Tous les noms prennent la majuscule : \all{die Wahrheit, der Beweis, die Falschmeldung}.
\end{itemize}
\end{attention}

\courssub{Exercice résolu et production guidée}

\begin{exoresolu}[Transformer un résumé au discours indirect]
\textbf{Énoncé} — Voici trois phrases du texte, au discours direct. Réécrivez-les comme dans un résumé, en les introduisant par le verbe entre parenthèses.
\begin{enumerate}
\item \all{Viele Menschen informieren sich nur noch über soziale Medien.} (Der Autor erklärt, dass ...)
\item \all{Falschmeldungen verbreiten sich schneller als die Wahrheit.} (Ein Experte sagt, ...)
\item \all{Man muss jede Nachricht überprüfen.} (Der Autor meint, ...)
\end{enumerate}
\tcblower
\textbf{Solution détaillée} —
\begin{enumerate}
\item Avec \all{dass}, le verbe conjugué va à la fin. Le sujet \all{viele Menschen} est au pluriel : le Konjunktiv I \all{informieren} serait identique à l'indicatif, on utilise donc la forme de remplacement Konjunktiv II \all{informierten} : \all{Der Autor erklärt, dass sich viele Menschen nur noch über soziale Medien informierten.} « L'auteur explique que beaucoup de gens ne s'informent plus que par les réseaux sociaux. »
\item Sans \all{dass}, le verbe reste en position 2 ; même remplacement au pluriel : \all{verbreiten} → \all{verbreiteten} : \all{Ein Experte sagt, Falschmeldungen verbreiteten sich schneller als die Wahrheit.} « Un expert dit que les fausses informations se répandraient plus vite que la vérité. »
\item \all{muss} (3\up{e} pers. sing.) → Konjunktiv I \all{müsse}, forme distincte de l'indicatif : \all{Der Autor meint, man müsse jede Nachricht überprüfen.} « L'auteur estime qu'il faudrait vérifier chaque information. »
\end{enumerate}
\end{exoresolu}

\begin{exercice}[Production guidée : résumé et opinion]
\textbf{A. Résumé en 5 phrases.} Résumez le texte « Fake News – die neue Gefahr im Internet » en cinq phrases, sans recopier le texte. Consignes : commencez par \all{Der Text handelt von ...}, utilisez au moins trois verbes de parole (\all{sagen, erklären, behaupten, raten, meinen}) et mettez les propos rapportés au Konjunktiv I.

\textbf{B. Court paragraphe d'opinion (5-6 lignes).} Répondez à la question : \all{Sind Fake News auch in Kamerun ein Problem?} Utilisez \all{meiner Meinung nach}, une subordonnée en \all{dass}, et un exemple concret (WhatsApp, Facebook, les bendskins, le football).
\end{exercice}

\begin{corrige}[Exercice : éléments de réponse]
\textbf{A. Résumé possible} : \all{Der Text handelt von Fake News im Internet. Der Autor erklärt, dass sich viele Menschen nur noch über soziale Medien informierten. Ein Experte sagt, Falschmeldungen verbreiteten sich schneller als die Wahrheit. Ein Journalist behauptet, manche Fake News seien bezahlt worden. Schließlich rät der Autor, jede Nachricht zu überprüfen.}

\textbf{B. Opinion possible} : \all{Meiner Meinung nach sind Fake News auch in Kamerun ein großes Problem. Ich finde, dass viele Menschen falsche Nachrichten über WhatsApp bekommen und sie sofort weiterschicken. Zum Beispiel gab es falsche Meldungen über Fußballspieler oder über die Schulen. Ich bin der Meinung, dass man jede Nachricht zuerst überprüfen muss. Die Schulen sollten die Jugendlichen über diese Gefahr informieren.}
\end{corrige}

\begin{experience}[Débat en classe : « Wahrheit oder Lüge ? »]
Par groupes de quatre : deux élèves jouent des journalistes sérieux, deux autres défendent la liberté totale de publier sur les réseaux sociaux. Chaque groupe prépare trois arguments avec des phrases au discours indirect (\all{Unsere Gruppe meint, man müsse ... ; Der Journalist hat gesagt, ...}). La classe vote ensuite pour les arguments les plus convaincants. Durée : 15 minutes.
\end{experience}

\begin{aretenir}[L'essentiel de la section]
\begin{itemize}
\item Un commentaire de texte suit toujours le plan : \textbf{Einleitung → Analyse → Meinung → Schluss}.
\item Dans l'analyse, les propos de l'auteur se rapportent au \textbf{Konjunktiv I}, avec ou sans \all{dass} (verbe en position 2 ou à la fin) ; si la forme du KI est identique à l'indicatif, on la remplace par le Konjunktiv II.
\item Les connecteurs \all{zuerst, dann, außerdem, schließlich} structurent l'analyse ; \all{meiner Meinung nach} introduit l'avis personnel (verbe en position 2).
\item Un résumé reformule, ne recopie jamais le texte.
\item La conclusion répond à la problématique et peut s'ouvrir sur le contexte camerounais.
\end{itemize}
\end{aretenir}

\newpage

\coursec{Activité d'intégration}

\courssub{Objectifs de l'activité}

Cette activité d'intégration clôture la leçon sur les fake news en mobilisant simultanément les trois compétences travaillées : le vocabulaire thématique de la vérité et du mensonge, le discours indirect au subjonctif I et la subordonnée introduite par \all{dass}. Elle reproduit une situation authentique du métier de journaliste : vérifier une information avant de la publier. À l'issue de l'activité, l'élève sera capable de :

\begin{itemize}
\item identifier une fausse information en analysant sa source, son expéditeur et ses preuves ;
\item employer correctement le lexique de la leçon : \all{die Falschmeldung}, \all{die Lüge}, \all{die Wahrheit}, \all{verbreiten}, \all{glauben}, \all{der Beweis}, \all{die Quelle}, \all{der Absender} ;
\item rapporter des propos au discours indirect, avec et sans \all{dass}, en conjuguant correctement le subjonctif I ;
\item rédiger et présenter oralement un court article de presse en allemand, compétence directement réinvestissable dans l'épreuve de rédaction du Baccalauréat.
\end{itemize}

\begin{experience}[Fact-Checker für einen Tag]
Travail en groupes de 4. Durée : 45 minutes (20 min de préparation écrite, 20 min de présentations orales, 5 min de bilan). Chaque groupe reçoit trois messages WhatsApp et doit jouer le rôle d'une équipe de vérificateurs d'informations (\all{die Faktenchecker}) pour un journal lycéen de Yaoundé.
\end{experience}

\courssub{Situation}

Vous êtes membres de la rédaction du journal de votre lycée à Yaoundé, \all{„Die junge Stimme“}. En période d'examens, de nombreux messages circulent sur les groupes WhatsApp des élèves. Votre rédacteur en chef vous confie une mission : vérifier trois messages reçus ce matin et publier un article d'avertissement intitulé \all{„Achtung, Fake News!“}.

\begin{textebox}[Drei WhatsApp-Nachrichten aus der Gruppe „Tle A Yaoundé“]
\textbf{Nachricht 1} (Absender: „InfoRapide237“, unbekannte Nummer, 6 Uhr 12) : „SCHOCK!!! Das BAC-Examen wurde abgesagt! Alle Schüler bekommen automatisch 20/20! Ein Freund, der im Ministerium arbeitet, hat das bestätigt. Teilt diese Nachricht schnell, bevor sie gelöscht wird!!!“

\textbf{Nachricht 2} (Absender: offizielle Seite des Gymnasiums, 7 Uhr 30) : „Liebe Schülerinnen und Schüler, die Deutschprüfung findet wie geplant am Montag um 8 Uhr statt. Bitte bringen Sie Ihre Ausweise mit. Die Direktion.“

\textbf{Nachricht 3} (Absender: ein Mitschüler, 8 Uhr 05) : „Ich habe gehört, dass der Direktor die Kantine schließen will. Weiß jemand, ob das stimmt?“
\end{textebox}

Glossaire : \all{absagen} (annuler), \all{automatisch} (automatiquement), \all{bestätigen} (confirmer), \all{teilen} (partager), \all{löschen} (supprimer), \all{der Ausweis, -e} (la pièce d'identité), \all{die Kantine, -n} (la cantine), \all{stimmen} (être exact), \all{unbekannt} (inconnu), \all{der Mitschüler, -} (le camarade de classe).

\courssub{Consignes}

\begin{enumerate}
\item \textbf{Analyse des messages.} Lisez les trois messages. Pour chacun, notez dans un tableau : \all{Absender} (expéditeur), \all{Quelle} (source), \all{Beweis} (preuve), \all{Sprache} (style : majuscules, points d'exclamation, fautes). Identifiez la vraie information, la fake news et la simple rumeur (\all{das Gerücht}).

\item \textbf{Justification orale.} Dans votre groupe, justifiez votre choix avec le vocabulaire de la leçon. Exemple de départ : \all{„Wir glauben, dass Nachricht 1 eine Falschmeldung ist, weil ...“}

\item \textbf{Rédaction de l'article.} Rédigez un article de 8 à 10 lignes intitulé \all{„Achtung, Fake News!“}. Vous devez rapporter les propos des internautes et des élèves au \cle{discours indirect} (subjonctif I) : au moins deux phrases avec \all{dass} et deux phrases sans \all{dass}. Exemple : \all{Ein Nutzer sagte, er habe die Nachricht geglaubt.}

\item \textbf{Relecture collective.} Vérifiez : la conjugaison du subjonctif I, la place du verbe conjugué à la fin dans la subordonnée en \all{dass}, la majuscule à tous les noms, le vocabulaire du thème (au moins cinq mots de la leçon).

\item \textbf{Présentation orale.} Un porte-parole du groupe lit l'article devant la classe. Les autres groupes écoutent et posent une question, par exemple : \all{„Woher wisst ihr, dass die Quelle nicht seriös ist?“}
\end{enumerate}

\begin{methode}[Comment vérifier une information en quatre étapes]
\begin{enumerate}
\item \textbf{Wer?} Qui parle ? Identifier l'expéditeur (\all{der Absender}) : numéro inconnue, page officielle, ami ? Un expéditeur anonyme est un premier signal d'alarme.
\item \textbf{Woher?} D'où vient l'information ? Chercher la source (\all{die Quelle}) : « un ami qui travaille au ministère » n'est pas une source vérifiable.
\item \textbf{Welche Beweise?} Quelles preuves (\all{die Beweise}) sont données ? Une vraie information cite des faits, des dates, des documents.
\item \textbf{Wie?} Comment le message est-il écrit ? Majuscules partout, points d'exclamation multiples, urgence artificielle (« partagez vite avant suppression ! ») : le style trahit souvent la fake news.
\end{enumerate}
\end{methode}

\courssub{Ressources à mobiliser}

\begin{aretenir}[Boîte à outils de l'élève]
\begin{itemize}
\item \textbf{Lexique :} \all{die Falschmeldung, -en} (la fausse information), \all{die Lüge, -n} (le mensonge), \all{die Wahrheit, -en} (la vérité), \all{das Gerücht, -e} (la rumeur), \all{die Quelle, -n} (la source), \all{der Absender, -} (l'expéditeur), \all{der Beweis, -e} (la preuve), \all{verbreiten} (diffuser), \all{glauben} (croire), \all{prüfen} (vérifier), \all{seriös} (sérieux, fiable).
\item \textbf{Subjonctif I :} formé sur le radical de l'infinitif + \all{-e, -est, -e, -en, -et, -en}. Formes clés : \all{er habe}, \all{er sei}, \all{er werde}, \all{er könne}. Au passé : \all{er habe geglaubt}, \all{er sei gekommen}.
\item \textbf{Avec dass :} verbe conjugué à la fin. \all{Ein Schüler sagte, dass er die Nachricht geglaubt habe.}
\item \textbf{Sans dass :} verbe conjugué en deuxième position. \all{Ein Schüler sagte, er habe die Nachricht geglaubt.}
\end{itemize}
\end{aretenir}

\begin{popfigure}
\begin{tikzpicture}[node distance=0.6cm and 1.6cm]
\node[draw=popblue, fill=popblueL, rounded corners, align=center, text width=4.6cm] (direct) {\textbf{Discours direct}\\ \all{„Ich habe die}\\ \all{Nachricht geglaubt.“}};
\node[draw=poporange, fill=poporangeL, rounded corners, align=center, text width=5.2cm, right=of direct] (indirect) {\textbf{Discours indirect (Konj. I)}\\ \all{Er sagte, er habe die}\\ \all{Nachricht geglaubt.}};
\draw[-{Stealth}, very thick, popdark] (direct) -- (indirect) node[midway, above, align=center, font=\small] {indicatif $\rightarrow$ Konj. I\\ 1\iere{} pers. $\rightarrow$ 3\ieme{} pers.};
\node[draw=popgreen, fill=popgreenL, rounded corners, align=center, text width=5.2cm, below=of indirect] (dass) {\textbf{Variante avec \all{dass}}\\ \all{Er sagte, dass er die}\\ \all{Nachricht geglaubt habe.}};
\draw[-{Stealth}, very thick, popdark] (indirect) -- (dass) node[midway, right, align=center, font=\small] {verbe conjugué\\ à la fin};
\end{tikzpicture}
\legende{Du discours direct au discours indirect : changement de personne, subjonctif I et place du verbe.}
\end{popfigure}

\begin{attention}[Pièges à éviter]
\begin{itemize}
\item Ne dites pas \all{„er haben“} au subjonctif I : la troisième personne du singulier est \all{er habe}. C'est la forme qui distingue le discours indirect de l'indicatif.
\item Après \all{dass}, le verbe conjugué va \textbf{à la fin} : pas \all{„dass er habe geglaubt die Nachricht“}, mais \all{dass er die Nachricht geglaubt habe}.
\item Quand le subjonctif I est identique à l'indicatif (par exemple \all{sie haben}), on le remplace par le subjonctif II : \all{sie hätten}. C'est la règle de substitution, obligatoire pour marquer le discours rapporté.
\item Faux ami : \all{die Lüge} signifie « le mensonge », pas « le lieu ». Et \all{glauben} signifie « croire », pas « regarder ».
\item N'oubliez pas la majuscule à tous les noms : \all{die Nachricht}, \all{die Quelle}, \all{der Beweis}.
\end{itemize}
\end{attention}

\courssub{Proposition de production et corrigé commenté}

\begin{exoresolu}[Article modèle : „Achtung, Fake News!“]
Rédigez un article de 8 à 10 lignes rapportant les propos recueillis, avec au moins deux subordonnées en \all{dass} et deux discours indirects sans \all{dass}.
\tcblower
\textbf{Production modèle :}

\all{„Achtung, Fake News! Seit gestern verbreitet sich eine Falschmeldung auf WhatsApp: Das BAC-Examen sei abgesagt worden. Ein unbekannter Absender behauptete, dass ein Freund im Ministerium die Nachricht bestätigt habe. Doch diese Quelle ist nicht seriös: Es gibt keinen Beweis. Viele Schüler sagten, sie hätten die Nachricht zuerst geglaubt. Eine Schülerin erklärte, dass sie die Lüge sofort erkannt habe, weil der Text voller Fehler sei. Der Direktor versicherte, dass die Prüfung wie geplant stattfinden werde. Unser Rat: Prüft immer die Quelle, bevor ihr eine Nachricht teilt! Verbreitet die Wahrheit, keine Lügen!“}

\textbf{Commentaire des choix de langue :}
\begin{itemize}
\item \all{sei abgesagt worden} : subjonctif I du passif au passé, sans \all{dass}, verbe en deuxième position — forme typique du style journalistique.
\item \all{behauptete, dass ein Freund ... bestätigt habe} : subordonnée en \all{dass}, verbe conjugué \all{habe} rejeté à la fin, auxiliaire du subjonctif I pour le discours rapporté.
\item \all{sie hätten die Nachricht zuerst geglaubt} : discours indirect sans \all{dass} ; le subjonctif II \all{hätten} remplace ici le subjonctif I \all{haben}, identique à l'indicatif au pluriel — c'est la règle de substitution.
\item \all{dass sie die Lüge sofort erkannt habe, weil der Text voller Fehler sei} : double discours indirect, avec verbe conjugué final dans les deux subordonnées.
\item \all{stattfinden werde} : subjonctif I de \all{werden} + infinitif, pour rapporter une promesse portant sur l'avenir.
\item Lexique de la leçon mobilisé : \all{Falschmeldung}, \all{Absender}, \all{Quelle}, \all{Beweis}, \all{Lüge}, \all{Wahrheit}, \all{verbreiten}, \all{glauben}, \all{seriös}.
\end{itemize}
\end{exoresolu}

\begin{corrige}[Grille d'évaluation]
\begin{center}
\begin{tabularx}{\linewidth}{|X|c|}
\hline
\rowcolor{popblueL} \textbf{Critère} & \textbf{Points} \\
\hline
Vocabulaire du thème (au moins 5 mots correctement employés) & 5 \\
\hline
Subjonctif I (formes correctes, y compris la substitution) & 5 \\
\hline
Subordonnées en \all{dass} (verbe à la fin) et tournures sans \all{dass} & 5 \\
\hline
Clarté, structure de l'article, présentation orale & 5 \\
\hline
\textbf{Total} & \textbf{20} \\
\hline
\end{tabularx}
\end{center}
\end{corrige}

\courssub{Bilan de l'activité}

Cette activité a permis de réinvestir l'ensemble des acquis de la leçon dans une situation de communication authentique et proche de la vie des élèves : la circulation de rumeurs sur WhatsApp en période d'examens. Sur le plan lexical, les élèves ont manipulé le vocabulaire de la vérité et du mensonge dans un contexte réel de vérification. Sur le plan grammatical, le discours indirect au subjonctif I et la subordonnée en \all{dass} ont été employés non comme des exercices mécaniques, mais comme des outils du journaliste qui rapporte des propos avec distance et prudence — exactement la compétence attendue dans les épreuves de rédaction et de commentaire du Baccalauréat. Sur le plan citoyen enfin, l'activité développe l'esprit critique face aux informations numériques : vérifier la source, chercher des preuves, ne pas partager trop vite. Cette compétence, essentielle à l'ère des réseaux sociaux, pourra être réinvestie dans la leçon suivante consacrée aux autres formes de cybercriminalité.

\newpage

\coursec{Méthodes et exercices résolus}

\courssub{Méthodes et exercices résolus}

\begin{methode}[Comprendre un texte d’actualité : lecture globale et détaillée]
\textbf{Objectif :} Extraire l’information essentielle, identifier la structure argumentative et le lexique thématique d’un texte journalistique sur les \emph{Fake News}.
\begin{enumerate}
\item \textbf{Anticiper :} Lire le titre, le sous-titre, la source, la date. Formuler des hypothèses sur le contenu (mots-clés attendus : \all{die Falschmeldung}, \all{verbreiten}, \all{die soziale Medien}, \all{der Wahrheitsgehalt}).
\item \textbf{Lecture globale (Skimming) :} Lire vite sans s’arrêter au vocabulaire inconnu. Répondre : Quel est le \textbf{sujet} ? Quel est le \textbf{thème principal} (danger, solution, exemple) ? Quel est le \textbf{ton} (alarmiste, analytique, neutre) ?
\item \textbf{Lecture détaillée (Scanning) :} Relire par paragraphes. Souligner : les \textbf{connecteurs logiques} (\all{denn}, \all{weil}, \all{jedoch}, \all{außerdem}), les \textbf{chiffres/dates}, les \textbf{verbes d’opinion} (\all{behaupten}, \all{vermuten}, \all{bestätigen} — souvent au \textbf{Konjunktiv I}).
\item \textbf{Lexique en contexte :} Deviner le sens des mots inconnus par la structure (préfixes \all{ver-}, \all{ent-}, \all{miss-}) et le co-texte. Noter les \textbf{familles de mots} (\all{die Lüge} $\to$ \all{lügen}, \all{der Lügner}, \all{belügen}).
\item \textbf{Synthèse :} Rédiger en 2 phrases (allemand) : l’idée centrale + l’argument clé de l’auteur.
\end{enumerate}
\end{methode}

\begin{exoresolu}[Compréhension détaillée : « Fake News im Netz »]
\textbf{Texte support (extrait) :}
\begin{textebox}[Extrait d’un article de la \emph{Deutsche Welle}, adapté]
„Laut einer Studie der Universität Oxford verbreiten sich Falschmeldungen im Netz sechsmal schneller als wahre Nachrichten. Besonders in sozialen Netzwerken wie X oder TikTok erreichen sie Millionen Nutzer, bevor Faktenchecker sie widerlegen können. Der Algorithmus belohnt Emotionen, nicht die Wahrheit. Experten warnen: „Wer alles glaubt, was er liest, verliert den Bezug zur Realität.“ Die Bundesregierung plant deshalb ein neues Gesetz gegen Desinformation.“
\end{textebox}

\textbf{Exercices :}
\begin{enumerate}
\item Répondez en allemand, par des phrases complètes :
\begin{enumerate}
\item Wie schnell verbreiten sich Falschmeldungen im Vergleich zu wahren Nachrichten?
\item Welche Rolle spielt der Algorithmus laut dem Text?
\item Was plant die Bundesregierung?
\end{enumerate}
\item Trouvez dans le texte les mots correspondant à ces définitions :
\begin{enumerate}
\item Person, die Fakten überprüft (Substantiv, maskulin, Plural: \ldots).
\item Etwas als falsch beweisen (Verb, Infinitiv).
\item Verlust der Verbindung zur Realität (Ausdruck mit \all{Verlust} / \all{Verbindung}).
\end{enumerate}
\item Grammaire en contexte : Relevez \textbf{deux} verbes au \textbf{Konjunktiv I} dans la citation directe des experts. Donnez leur infinitif et leur personne.
\end{enumerate}

\tcblower
\textbf{Solution détaillée :}
\textbf{1. Questions de compréhension :}
\begin{itemize}
\item[a)] \all{Sich Falschmeldungen im Netz sechsmal schneller als wahre Nachrichten.} (Reprise du \all{sechsmal schneller} + comparatif \all{als}). \textit{Attention : ordre des mots en phrase complète : Subjekt + Verb + Adverbialbestimmung.}
\item[b)] \all{Der Algorithmus belohnt Emotionen und nicht die Wahrheit.} (Présent indicatif, verbe à la 2e position).
\item[c)] \all{Die Bundesregierung plant ein neues Gesetz gegen Desinformation.} (Vocabulaire : \all{die Desinformation} = la désinformation).
\end{itemize}
\textbf{2. Recherche lexicale :}
\begin{itemize}
\item[a)] \all{der Faktenchecker, die Faktenchecker} (Mot composé : \all{Fakt} + \all{en} + \all{checker}).
\item[b)] \all{widerlegen} (Préfixe \all{wider-} = contre ; \all{legen} = poser $\to$ réfuter). \textit{Ne pas confondre avec \all{wiederlegen} (orthographe erronée fréquente).}
\item[c)] \all{den Bezug zur Realität verlieren} (Expression figée : \all{Verbindung} possible mais \all{Bezug} est dans le texte \all{Bezug zur Realität}).
\end{itemize}
\textbf{3. Konjunktiv I repéré :}
\begin{itemize}
\item \all{glaubt} (3e pers. sg. de \all{glauben} $\to$ Konjunktiv I : \all{er glaube} — \textbf{Attention} : dans le texte \all{„Wer alles glaubt\ldots“} c’est de l’indicatif ! C’est une proposition relative introduite par \all{wer}. Le Konjunktiv I apparaît dans le discours indirect \emph{rapporté} par le journaliste : \all{Experten warnen: „Wer alles glaubt\ldots“} $\to$ Discours indirect : \all{Experten warnen, wer alles \textbf{glaube}\ldots}).
\item \all{verliert} (Indicatif dans la citation). Le Konjunktiv I serait \all{verliere}.
\end{itemize}
\textit{Conclusion :} Le texte original utilise le discours direct (guillemets „…“). L’exercice vise à repérer les formes qui \textbf{deviendront} du Konjunktiv I lors de la transformation en discours indirect (voir Méthode 3).
\end{exoresolu}

\begin{methode}[Maîtriser le lexique thématique : familles de mots, collocations et genres]
\textbf{Objectif :} Construire un réseau lexical actif autour de \all{Wahrheit} et \all{Lüge} pour l’expression écrite/orale.
\begin{enumerate}
\item \textbf{Tableau des familles (der/die/das + Plural) :} Apprendre par \textbf{bloc genre+pluriel}.
\begin{center}
\begin{tabularx}{\linewidth}{|X|X|X|}\hline
\rowcolor{popblueL} \textbf{Nom (Art. Pl.)} & \textbf{Verbe} & \textbf{Adjectif / Participe} \\ \hline
\all{die Lüge, die Lügen} & \all{lügen} (regelmäßig) & \all{lügenhaft, gelogen} \\ \hline
\all{die Wahrheit, die Wahrheiten} & \all{wahrheitsgemäß handeln} & \all{wahr, wahrhaftig} \\ \hline
\all{die Falschmeldung, die -en} & \all{falsch melden} (trennbar) & \all{falsch, irreführend} \\ \hline
\all{der Beweis, die Beweise} & \all{beweisen} & \all{beweisbar, bewiesen} \\ \hline
\all{die Verbreitung, -en} & \all{verbreiten} (Akk.) / \all{sich verbreiten} & \all{verbreitet, viral} \\ \hline
\all{der Faktencheck, -s} & \all{faktenchecken} & \all{gecheckt, verifiziert} \\ \hline
\end{tabularx}
\end{center}
\item \textbf{Collocations fixes (à apprendre par cœur) :}
\all{eine Lüge verbreiten} | \all{die Wahrheit ans Licht bringen} | \all{Beweise liefern / erbringen} | \all{etw. widerlegen} | \all{Quellen prüfen} | \all{Desinformation bekämpfen}.
\item \textbf{Préfixes productifs :} \all{ent-} (\all{entlarven} = démasquer), \all{ver-} (\all{verbreiten}, \all{verheimlichen}), \all{miss-} (\all{missbrauchen}).
\item \textbf{Faux amis / Pièges :} \all{die Lüge} (mensonge) $\neq$ \all{die Lage} (situation). \all{beweisen} (prouver) $\neq$ \all{bevorzugen} (préférer). \all{verbreiten} (répandre) $\neq$ \all{verbreiten} (s’étendre — réfléchi \all{sich verbreiten}).
\end{enumerate}
\end{methode}

\begin{exoresolu}[Lexique : transformation et emploi en contexte]
\textbf{Exercices :}
\begin{enumerate}
\item \textbf{Complétez} avec le mot de la même famille (nom, verbe, adjectif) en respectant la casse et le genre :
\begin{enumerate}
\item \all{Die \_\_\_\_\_\_\_\_ (Lüge) im Netz ist ein großes Problem.}
\item \all{Man muss Informationen immer \_\_\_\_\_\_\_\_ (Beweis) überprüfen.}
\item \all{Das Video war \_\_\_\_\_\_\_\_ (irreführen), aber nicht illegal.}
\item \all{Die Polizei konnte den Täter \_\_\_\_\_\_\_\_ (entlarven).}
\end{enumerate}
\item \textbf{Traduisez en allemand} (thème lexical) :
\begin{enumerate}
\item Les fausses informations se propagent viralement.
\item Il faut apporter des preuves.
\item Le gouvernement veut lutter contre la désinformation.
\end{enumerate}
\item \textbf{Choisissez le bon mot} (\all{verbreiten} vs \all{sich verbreiten} ; \all{widerlegen} vs \all{widersprechen}) :
\begin{enumerate}
\item \all{Das Gerücht \_\_\_\_\_\_\_\_ sich schnell in der Schule.}
\item \all{Er \_\_\_\_\_\_\_\_ die These seines Kollegen mit Fakten.}
\item \all{Man sollte keine Falschmeldungen \_\_\_\_\_\_\_\_.}
\end{enumerate}
\end{enumerate}

\tcblower
\textbf{Solution détaillée :}
\textbf{1. Familles de mots (Orthographe \& Genre) :}
\begin{enumerate}
\item \all{die Lüge} (Nom féminin, même forme au singulier). \textit{Piège : ne pas écrire *\all{die Lüg} ou *\all{das Lüge}.}
\item \all{beweisen} (Verbe) ou \all{beweisbar} (Adjectif). Contexte \all{muss \ldots \_\_\_\_\_\_\_\_ überprüfen} $\to$ infinitif \all{beweisen} (séparable ? non). \textit{Mieux :} \all{belegen} (étayer) ou \all{verifizieren}. Attendu : \all{beweisen}.
\item \all{irreführend} (Participe I adjectivé de \all{irreführen}, séparable). \textit{Attention à l’Umlaut : \all{führen} $\to$ \all{irreführend}.}
\item \all{entlarven} (Verbe, préfixe \all{ent-} inaccentué $\to$ pas de \all{ge-} au Participe II : \all{entlarvt}). Ici infinitif après modale \all{könnte}.
\end{enumerate}
\textbf{2. Thème (Traduction vers l’allemand) :}
\begin{enumerate}
\item \all{Falschmeldungen verbreiten sich viral.} (\all{sich verbreiten} réfléchi + adverbe \all{viral}). \textit{Ordre : Sujet + Verbe + Complément + Adverbe.}
\item \all{Man muss Beweise liefern / erbringen.} (\all{Beweise} Accusatif pluriel). \textit{Collocation : \all{Beweise liefern}.}
\item \all{Die Regierung will Desinformation bekämpfen.} (\all{Die Bundesregierung} ou \all{Die Regierung}, \all{bekämpfen} verbe fort \all{kämpfen}).
\end{enumerate}
\textbf{3. Choix verbal (Transitivité / Réflexif) :}
\begin{enumerate}
\item \all{verbreitet} (\all{sich verbreiten} = se propager, intransitif/réfléchi). Sujet = \all{Das Gerücht}.
\item \all{widerlegt} (\all{etw. (Akk.) widerlegen} = réfuter qqch). Objet = \all{die These}.
\item \all{verbreiten} (\all{etw. (Akk.) verbreiten} = répandre qqch, transitif). Objet = \all{Falschmeldungen}.
\end{enumerate}
\textit{Règle :} \all{verbreiten} + Akk = action sur objet ; \all{sich verbreiten} = processus autonome.
\end{exoresolu}

\begin{methode}[Appliquer le discours indirect (Konjunktiv I) : règles de transformation]
\textbf{Objectif :} Reporter des propos (journalistiques, témoignages) sans les citer directement, marque de distance journalistique.
\begin{enumerate}
\item \textbf{Identifier le verbe introducteur :} \all{sagen}, \all{behaupten}, \all{berichten}, \all{melden}, \all{warnen}, \all{bestätigen}, \all{hinzufügen} (souvent au Präteritum : \all{behauptete}, \all{meldete}).
\item \textbf{Choisir le mode :} \textbf{Konjunktiv I} (forme standard). Si forme KI = Indicatif (ex. \all{ich bin} $\to$ \all{ich sei} ; \all{wir haben} $\to$ \all{wir haben} — identique !), utiliser \textbf{Konjunktiv II} (\all{ich wäre}, \all{wir hätten}) ou \all{würde}-Forme.
\item \textbf{Transformer les pronoms :} \all{ich} $\to$ \all{er/sie} ; \all{wir} $\to$ \all{sie} ; \all{mein} $\to$ \all{sein/ihr} ; \all{heute} $\to$ \all{an jenem Tag} ; \all{hier} $\to$ \all{dort}.
\item \textbf{Structure de la phrase :} Verbe introducteur (Pos. 2) + \textbf{\all{dass}} + \ldots + \textbf{Verbe au KI à la fin}. \textit{Exception : style journalistique sans \all{dass} (verbe en position 2), niveau supérieur. Au Bac : exiger \all{dass}.}
\item \textbf{Tableau KI essentiel (verbes forts/irréguliers) :}
\begin{center}
\begin{tabular}{|l|l|l|l|}\hline
\rowcolor{popblueL} \textbf{Infinitif} & \textbf{ich} & \textbf{er/sie/es} & \textbf{wir/sie/Sie} \\ \hline
\all{sein} & \all{sei} & \all{sei} & \all{seien} \\ \hline
\all{haben} & \all{habe} & \all{habe} & \all{haben} (\textsf{= Indik.!}) $\to$ \all{häten} (KII) \\ \hline
\all{werden} & \all{werde} & \all{werde} & \all{werden} (\textsf{= Indik.!}) $\to$ \all{würden} (KII) \\ \hline
\all{können/müssen/sollen/wollen} & \all{könne/müsse/sole/wolle} & \all{könne/müsse/sole/wolle} & \all{können/müssen/sollen/wollen} (\textsf{= Indik.!}) $\to$ KII \\ \hline
\all{verbreiten} & \all{verbreite} & \all{verbreite} & \all{verbreiten} \\ \hline
\end{tabular}
\end{center}
\end{enumerate}
\end{methode}

\begin{exoresolu}[Transformation : Discours direct $\to$ Indirect (Konjunktiv I)]
\textbf{Énoncé :} Transformez les phrases suivantes en \textbf{discours indirect} avec \all{dass} + \textbf{Konjunktiv I}. Le verbe introducteur est au \textbf{Präteritum}.
\begin{enumerate}
\item \all{Der Journalist: „Die Falschmeldung verbreitet sich rasend schnell.“}
\item \all{Der Experte: „Wir haben keine Beweise für diese Theorie.“}
\item \all{Die Ministerin: „Ich werde ein neues Gesetz vorschlagen.“}
\item \all{Ein Nutzer: „Ich kann die Quelle nicht überprüfen.“}
\item \all{Die Polizei: „Der Täter ist gefasst.“}
\end{enumerate}
\textbf{Bonus (Version) :} Traduisez en français la phrase 2 transformée.

\tcblower
\textbf{Solution détaillée :}
\textbf{Rappel :} Verbe introducteur (Prät.) + \all{dass} + Sujet + Compléments + \textbf{Verbe KI à la fin}.
\begin{enumerate}
\item \all{Der Journalist meldete, dass sich die Falschmeldung rasend schnell \textbf{verbreite}.}
(\all{sich verbreiten} réfléchi $\to$ \all{sich} reste devant le verbe final. KI de \all{verbreiten} : \all{er verbreite}).
\item \all{Der Experte berichtete, dass sie \textbf{keine Beweise} für diese Theorie \textbf{hätten}.}
(\all{haben} : KI \all{haben} = Indikatif $\to$ \textbf{remplacement par KII \all{häten}}). \textit{Traduction :} L'expert a rapporté qu'ils n'avaient aucune preuve pour cette théorie.
\item \all{Die Ministerin ankündigte, dass sie ein neues Gesetz \textbf{vorschlagen werde}.}
(\all{werden} + Infinitif : KI de \all{werden} = \all{werde}. \all{vorschlagen} séparable $\to$ \all{vorschlagen} à la fin devant \all{werde}). \textit{Ordre : Particule + Infinitif + Auxiliaire KI.}
\item \all{Ein Nutzer gab zu, dass er die Quelle nicht \textbf{überprüfen könne}.}
(\all{können} : KI \all{könne}. \all{überprüfen} infinitif à l’avant-dernière place).
\item \all{Die Polizei bestätigte, dass der Täter \textbf{gefasst sei}.}
(\all{sein} : KI \all{sei}. Participe II \all{gefasst} devant \all{sei}).
\end{enumerate}
\textbf{Points de vigilance :}
\begin{itemize}
\item N°2 : \all{haben} $\to$ KII obligatoire (\all{häten}) car KI identique à l’indicatif.
\item N°3 : Futur (devenu modal \all{werde} + Inf.) $\to$ KI de \all{werden} (\all{werde}).
\item N°5 : Passif (\all{ist gefasst}) $\to$ Auxiliaire \all{sein} au KI (\all{sei}).
\end{itemize}
\end{exoresolu}

\begin{methode}[Construire la subordonnée en « dass » : syntaxe, temps et connecteurs]
\textbf{Objectif :} Produire des phrases complexes correctes pour exprimer l’opinion, la cause, le contenu d’une information.
\begin{enumerate}
\item \textbf{Structure immuable :} \all{Hauptsatz (V2) , \textbf{dass} + Subjekt + Compléments + \textbf{Verbe conjugué à la fin}.}
\item \textbf{Verbes introducteurs fréquents (gouvernant \all{dass}) :}
\begin{itemize}
\item Opinion/Savoir : \all{denken, glauben, meinen, wissen, bezweifeln}.
\item Parole : \all{sagen, berichten, behaupten, versprechen, bestreiten}.
\item Sentiment : \all{bedauern, freuen, ärgern, wundern} (souvent + \all{sich}).
\end{itemize}
\item \textbf{Concordance des temps (Consecutio temporum) :}
\begin{center}
\begin{tabularx}{\linewidth}{|X|X|X|}\hline
\rowcolor{popblueL} \textbf{Principal (V. intro)} & \textbf{Subordonnée (contenu)} & \textbf{Exemple} \\ \hline
Présent (\all{glaube}) & Présent / Futur / Perfekt & \all{Ich glaube, dass er \textbf{kommt} / \textbf{kommen wird} / \textbf{gekommen ist}.} \\ \hline
Präteritum (\all{glaubte}) & Präteritum / Plusquamperfekt & \all{Ich glaubte, dass er \textbf{kam} / \textbf{gekommen war}.} \\ \hline
Perfekt (\all{hat geglaubt}) & Perfekt / Plusquamperfekt & \all{Er hat geglaubt, dass er \textbf{gekommen ist} / \textbf{war}.} \\ \hline
\end{tabularx}
\end{center}
\item \textbf{Négation :} \all{nicht} / \all{kein} se place \textbf{avant} le verbe final (ou avant l’élément nié). Ex: \all{Ich glaube, dass er \textbf{nicht} kommt.}
\item \textbf{Verbes à particule séparable :} La particule rejoint le verbe à la fin. \all{Ich denke, dass er die Falschmeldung \textbf{weiterverbreitet}.} (mais \all{verbreitet \ldots weiter} au principal).
\end{enumerate}
\end{methode}

\begin{exoresolu}[Syntaxe : « dass »-Sätze, ordre des mots et concordance]
\textbf{Exercices :}
\begin{enumerate}
\item \textbf{Reliez} les deux phrases par une subordonnée \all{dass} (verbe introducteur entre parenthèses). Respectez la concordance des temps.
\begin{enumerate}
\item \all{Die Fake News sind gefährlich.} (\all{der Experte warnen} -- Präteritum)
\item \all{Der Algorithmus belohnt Emotionen.} (\all{die Studie zeigen} -- Präsens)
\item \all{Die Regierung hat ein Gesetz geplant.} (\all{die Zeitung melden} -- Perfekt)
\item \all{Man kann die Wahrheit nicht finden.} (\all{der Nutzer bezweifeln} -- Präteritum)
\end{enumerate}
\item \textbf{Corrigez les erreurs} (ordre des mots, conjugaison, particule) dans ces phrases d’élèves :
\begin{enumerate}
\item \all{Ich glaube, dass er kommt nicht.}
\item \all{Er sagte, dass er die Quelle überprüft hat.} (Contexte : discours indirect $\to$ KI attendu).
\item \all{Wir denken, dass die Lüge sich verbreitet schnell.}
\item \all{Sie behauptete, dass sie die Wahrheit kennt.} (Contexte : discours indirect).
\end{enumerate}
\item \textbf{Traduisez en allemand} (Thème) :
\begin{enumerate}
\item Je pense que les fake news détruisent la confiance.
\item Nous avons douté qu'il ait dit la vérité. (Indice : \all{zweifeln an + Dat.} ou \all{bezweifeln + Akk.} $\to$ ici \all{dass}-Satz avec \all{bezweifeln}).
\end{enumerate}
\end{enumerate}

\tcblower
\textbf{Solution détaillée :}
\textbf{1. Relier (Conjonction \all{dass} + verbe à la fin) :}
\begin{enumerate}
\item \all{Der Experte warnte, dass die Fake News gefährlich \textbf{waren}.} (Principal Prät. $\to$ Subordonnée Prät. \all{waren}). \textit{Note : \all{warnen} prend \all{dass} ou \all{vor + Dat.} Ici \all{dass}.}
\item \all{Die Studie zeigt, dass der Algorithmus Emotionen \textbf{belohnt}.} (Principal Présent $\to$ Subordonnée Présent).
\item \all{Die Zeitung hat gemeldet, dass die Regierung ein Gesetz \textbf{geplant hat}.} (Principal Perfekt $\to$ Subordonnée Perfekt. \all{planen} faible $\to$ \all{geplant}). \textit{Ordre : Participe + Auxiliaire.}
\item \all{Der Nutzer bezweifelte, dass man die Wahrheit \textbf{finden könne}.} (Principal Prät. $\to$ Subordonnée : \textbf{Konjunktiv I} car \all{bezweifeln} introduit un discours indirect/opinion ! \all{können} $\to$ \all{könne}). \textit{Piège majeur : verbes de doute/parole $\to$ KI.}
\end{enumerate}
\textbf{2. Correction d’erreurs types :}
\begin{enumerate}
\item \all{Ich glaube, dass er \textbf{nicht kommt}.} (\all{nicht} devant le verbe final \all{kommt}).
\item \all{Er sagte, dass er die Quelle \textbf{nicht überprüfen könne}.} (\textbf{Konjunktiv I} \all{könne} + infinitif \all{überprüfen} devant). \textit{Erreur élève : Indicatif Perfekt \all{hat überprüft} au lieu de KI.}
\item \all{Wir denken, dass sich die Lüge \textbf{schnell verbreitet}.} (\all{sich verbreiten} réfléchi : \all{sich} devant le verbe final \all{verbreitet}. Adverbe \all{schnell} avant \all{sich} ou avant \all{verbreitet}).
\item \all{Sie behauptete, dass sie die Wahrheit \textbf{kenne}.} (\textbf{Konjunktiv I} \all{kenne} pour \all{sie} (3e pl/sg formel) ou \all{kennen} (1e pl) — contexte \all{sie} = elle $\to$ \all{kenne}). \textit{Erreur : Indicatif \all{kennt}.}
\end{enumerate}
\textbf{3. Thème :}
\begin{enumerate}
\item \all{Ich denke, dass Fake News das Vertrauen \textbf{zerstören}.} (\all{zerstören} fort, Présent pluriel).
\item \all{Wir bezweifelten, dass er die Wahrheit \textbf{gesagt habe}.} (Principal Prät. \all{bezweifelten} $\to$ Subordonnée KI Perfekt : Auxiliaire \all{haben} au KI \all{habe} + Participe \all{gesagt}). \textit{Attention : \all{sagen} prend \all{haben}.}
\end{enumerate}
\end{exoresolu}

\begin{methode}[Rédiger un paragraphe argumenté court (80–100 mots) : structure et connecteurs]
\textbf{Objectif :} Produire une expression écrite cohérente (type Bac : \all{Stellungnahme} ou \all{Kommentar}) sur les dangers des Fake News.
\begin{enumerate}
\item \textbf{Analyser le sujet :} Identifier le thème (ex. \all{Sind soziale Medien schuld an Fake News?}) et la consigne (avantages/inconvénients, opinion personnelle, solutions).
\item \textbf{Plan en 3 étapes (Sandwich) :}
\begin{enumerate}
\item \textbf{These (These-Satz) :} Phrase d’accroche + position claire. \all{Meiner Meinung nach sind Fake News eine ernste Bedrohung für die Demokratie.}
\item \textbf{Arguments + Exemples (Argumente + Beispiele) :} 2 arguments distincts, chacun introduit par un connecteur. Structure : \all{Erstens\ldots / Zweitens\ldots} + \all{weil / da / denn} + \all{dass}-Satz. Exemple concret (Cameroun/Allemagne).
\item \textbf{Conclusion (Fazit / Lösungsvorschlag) :} \all{Zusammenfassend / Abschließend} + appel à l’action (\all{Man muss\ldots}, \all{Es ist wichtig, dass\ldots + KI}).
\end{enumerate}
\item \textbf{Connecteurs logiques (Niveau B1/B2) :}
\all{Einerseits \ldots andererseits} | \all{Nicht nur \ldots sondern auch} | \all{Auf der einen Seite \ldots auf der anderen Seite} | \all{Daher / Deswegen / Folglich}.
\item \textbf{Relecture (Check-liste) :} Majuscules (Noms) ? Verbe position 2 (principal) / fin (subordonnée) ? Cas (Akk/Dat après prépositions/verbes) ? KI après \all{behaupten/denken} au passé ? Umlauts/ß ?
\end{enumerate}
\end{methode}

\begin{exoresolu}[Expression écrite : « Gefahr durch Fake News »]
\textbf{Sujet :} \all{„Fake News sind eine Gefahr für die Gesellschaft. Wie kann man sich schützen?“} (Environ 90 mots).
\textbf{Consigne :} Rédigez un paragraphe structuré en allemand. Utilisez au moins : 1 subordonnée \all{dass} avec KI, 1 subordonnée \all{weil}, 1 verbe à particule séparable, le vocabulaire de la leçon.

\tcblower
\textbf{Modèle de rédaction (corrigé type) :}
\begin{textebox}[Production élève attendue]
Meiner Meinung nach sind Falschmeldungen eine große Gefahr für unsere Demokratie. Erstens verbreiten sie sich in sozialen Netzwerken extrem schnell, weil die Algorithmen Emotionen belohnen und nicht die Wahrheit. Zweitens glauben viele Nutzer alles, was sie lesen, ohne die Quellen zu prüfen. Die Bundesregierung behauptete kürzlich, dass sie ein Gesetz gegen Desinformation \textbf{plane} (KI). Meiner Ansicht nach muss man aber vor allem die Medienkompetenz in den Schulen \textbf{stärken}. Nur so können wir die Verbreitung von Lügen \textbf{eindämmen}.
\end{textebox}

\textbf{Analyse grammaticale du modèle :}
\begin{itemize}
\item \all{Meiner Meinung nach} (Prépositionnel + Génitif) $\to$ Position 1, Verbe \all{sind} Position 2.
\item \all{weil \ldots belohnen} (Subordonnée causale, verbe à la fin).
\item \all{ohne \ldots zu prüfen} (Construction infinitive, \all{zu} + infinitif à la fin).
\item \all{dass sie ein Gesetz \ldots plane} (\textbf{Konjunktiv I} \all{plane} après \all{behauptete} (Prät.) + \all{dass}). \textit{Verbe séparable \all{planen} (non séparable ici) mais \all{vorschlagen} serait \all{vorschlage}.}
\item \all{dass \ldots stärken} (Subordonnée finale implicite après \all{wichtig ist}, mais ici \all{dass} + Indicatif présent car certitude du scripteur). \textit{Variante KI :} \all{dass man die Medienkompetenz stärke}.
\item \all{Nur so können wir \ldots eindämmen} (Inversion \all{Nur so} $\to$ Verbe \all{können} Pos 2, Infinitif \all{eindämmen} à la fin).
\item \textbf{Vocabulaire ciblé :} \all{Falschmeldungen, verbreiten, Algorithmus, Quellen prüfen, Desinformation, Medienkompetenz, eindämmen}.
\end{itemize}
\textbf{Note Bac :} 14-16/20 si structure respectée, vocabulaire précis, KI correct, peu de fautes de cas/maj.
\end{exoresolu}

\begin{attention}[Les 5 erreurs qui coûtent des points au Bac]
\begin{enumerate}
\item \textbf{Majuscules oubliées sur les noms :} \all{*die lüge, *das internet, *die wahrheit} $\to$ \textbf{\all{die Lüge, das Internet, die Wahrheit}}. \textit{Règle d'or : Tout nom = Majuscule.}
\item \textbf{Ordre des mots dans la subordonnée \all{dass} :} \all{*Ich glaube, dass er kommt nicht.} $\to$ \textbf{\all{Ich glaube, dass er nicht kommt.}} (Verbe \textbf{à la fin}, \all{nicht} devant le verbe).
\item \textbf{Konjunktiv I absent ou mal formé (Discours indirect) :} \all{*Er sagte, dass er die Quelle überprüft hat.} $\to$ \textbf{\all{Er sagte, dass er die Quelle überprüfen könne.}} (KI de \all{können} + Infinitif). \textit{Piège \all{haben/sein/werden/modaux} $\to$ KII (\all{häten, wäre, würde}).}
\item \textbf{Faux amis / Calques français :}
\begin{itemize}
\item \all{*aktuell} = actuel (en temps) $\neq$ \all{actuellement} (\all{derzeit, momentan}).
\item \all{*kontrollieren} = vérifier/contrôler $\neq$ \all{contrôler} au sens \all{überwachen}.
\item \all{*realisieren} = réaliser (prendre conscience) $\neq$ \all{réaliser} un projet (\all{verwirklichen, umsetzen}).
\item \all{*desinformation} (minuscule, pas de \all{Des-} comme préfixe nominal standard) $\to$ \textbf{\all{Desinformation}} (Majuscule, nom féminin).
\end{itemize}
\item \textbf{Cas accusatif/datif après verbes/prépositions :} \all{*glauben dem Mann} (Akk !) $\to$ \all{glauben \textbf{dem} Mann} (Dativ). \all{*warten für den Bus} $\to$ \all{warten \textbf{auf} \textbf{den} Bus} (Akk). \all{*sich interessieren für die Thema} $\to$ \all{sich interessieren für \textbf{das} Thema} (Neutre !).
\end{enumerate}
\end{attention}

\newpage

\coursec{Exercices}

\courssub{Je vérifie mes connaissances}

\begin{exercice}[QCM : le subjonctif I]
Entoure la bonne réponse.

1. Le discours indirect en allemand se rend généralement par :
\begin{itemize}
\item a) le conditionnel
\item b) le subjonctif I (Konjunktiv I)
\item c) le futur II
\end{itemize}

2. \all{Er sagte, er \ldots\ krank.} (discours indirect sans \all{dass})
\begin{itemize}
\item a) ist
\item b) sei
\item c) wäre
\end{itemize}

3. Le subjonctif I de \all{haben} à la 3\up{e} personne du singulier est :
\begin{itemize}
\item a) hat
\item b) habe
\item c) hätte
\end{itemize}

4. Dans une subordonnée en \all{dass}, le verbe conjugué se place :
\begin{itemize}
\item a) en deuxième position
\item b) juste après \all{dass}
\item c) à la fin de la subordonnée
\end{itemize}

5. \all{Die Journalistin behauptet, sie habe nicht gelogen.} signifie :
\begin{itemize}
\item a) La journaliste affirme qu'elle a menti.
\item b) La journaliste affirme qu'elle n'a pas menti.
\item c) La journaliste nie avoir dit la vérité.
\end{itemize}
\end{exercice}

\begin{exercice}[Vrai ou faux ? Justifie ta réponse]
Réponds par \all{richtig} ou \all{falsch} et justifie en français ou en allemand.

1. \all{Die Falschmeldung} est un nom féminin.
2. Le pluriel de \all{der Beweis} est \all{die Beweise}.
3. \all{verbreiten} est un verbe faible qui signifie « vérifier ».
4. Au subjonctif I présent, \all{sein} se conjugue : \all{ich sei, du seist, er sei}.
5. On peut construire le discours indirect sans \all{dass} ; dans ce cas, le verbe conjugué reste en deuxième position, comme dans une phrase principale.
6. \all{glauben} et \all{beweisen} sont des synonymes.
\end{exercice}

\begin{exercice}[Vocabulaire : trouve le mot]
Complète chaque phrase avec le mot qui convient : \all{die Lüge, die Wahrheit, der Beweis, die Quelle, verbreiten, überprüfen}.

1. Bevor man eine Nachricht teilt, muss man die \ldots\ \ldots\ . (vérifier la source)
2. Er hat die \ldots\ gesagt : er war wirklich nicht dort. (la vérité)
3. Fake News \ldots\ sich sehr schnell im Internet. (se répandre)
4. Die Polizei hat einen \ldots\ gefunden : das Foto war manipuliert. (une preuve)
5. Eine \ldots\ ist das Gegenteil von der Wahrheit. (un mensonge)
\end{exercice}

\begin{exercice}[Conjugaison : tableau du subjonctif I]
Complète le tableau du subjonctif I présent des verbes suivants.

\begin{center}
\begin{tabular}{|l|l|l|l|l|l|}
\hline
\rowcolor{popblueL} Personne & sein & haben & sagen & verbreiten & werden \\
\hline
ich & \ldots & \ldots & \ldots & \ldots & \ldots \\
\hline
du & \ldots & \ldots & \ldots & \ldots & \ldots \\
\hline
er/sie/es & \ldots & \ldots & \ldots & \ldots & \ldots \\
\hline
wir & \ldots & \ldots & \ldots & \ldots & \ldots \\
\hline
ihr & \ldots & \ldots & \ldots & \ldots & \ldots \\
\hline
sie/Sie & \ldots & \ldots & \ldots & \ldots & \ldots \\
\hline
\end{tabular}
\end{center}
\end{exercice}

\courssub{Je m'entraîne}

\begin{exercice}[Du discours direct au discours indirect]
Mets les phrases suivantes au discours indirect, une fois \textbf{avec} \all{dass}, une fois \textbf{sans} \all{dass}. Commence par : \all{Der Experte sagte, \ldots}

1. \all{„Die Falschmeldungen verbreiten sich schneller als die Wahrheit.“}
2. \all{„Viele Jugendliche glauben alles, was sie im Netz lesen.“}
3. \all{„Die Quelle dieser Nachricht ist nicht zuverlässig.“}
4. \all{„Wir haben den Beweis gefunden.“}
5. \all{„Die Regierung wird gegen Fake News kämpfen.“}
6. \all{„Ich überprüfe jede Information vor dem Teilen.“}
\end{exercice}

\begin{exercice}[Texte à trous : Fake News]
Complète le texte avec les mots suivants, conjugués au subjonctif I quand c'est nécessaire : \all{sein (2 fois), haben, verbreiten, glauben, die Lüge, die Quelle, überprüfen}.

\all{Ein Medienexperte erklärte in einem Interview, Fake News \ldots\ heute ein großes Problem. Viele Menschen \ldots\ nicht mehr die Zeit, jede Nachricht zu \ldots . Deshalb \ldots\ sich \ldots\ oft schneller als die Wahrheit. Er sagte auch, die \ldots\ \ldots\ immer das Wichtigste. Wer sie nicht prüfe, \ldots\ oft alles, was er im Internet lese.}
\end{exercice}

\begin{exercice}[Appariement : mot et définition]
Relie chaque mot allemand à sa définition française.

\begin{center}
\begin{tabularx}{\linewidth}{|l|X|}
\hline
\rowcolor{popblueL} Mot allemand & Définition française \\
\hline
1. \all{die Falschmeldung, -en} & a. dire quelque chose qui n'est pas vrai \\
\hline
2. \all{der Beweis, -e} & b. information inventée pour tromper le public \\
\hline
3. \all{die Quelle, -n} & c. élément qui montre qu'un fait est vrai \\
\hline
4. \all{lügen} & d. origine d'une information (journal, site, témoin) \\
\hline
5. \all{überprüfen} & e. contrôler si une information est exacte \\
\hline
\end{tabularx}
\end{center}
\end{exercice}

\begin{exercice}[Traduction : thème guidé]
Traduis les phrases suivantes en allemand en utilisant le subjonctif I et, quand c'est possible, une subordonnée en \all{dass}.

1. Le journaliste a dit que les fausses informations se répandaient plus vite que la vérité.
2. Elle affirme qu'elle n'a pas menti.
3. Nous croyons que la source n'est pas fiable.
4. Le ministre a déclaré que le gouvernement lutterait contre les fake news.
\end{exercice}

\courssub{Je me prépare au Bac}

\begin{textebox}[Fake News — eine Gefahr für Kamerun?]
In Kamerun verbreiten sich Falschmeldungen besonders schnell über WhatsApp und Facebook. Ein Kommunikationsexperte aus Yaoundé erklärte in einem Interview, viele Bürger seien nicht vorsichtig genug. Sie teilten Nachrichten, ohne die Quelle zu überprüfen. Er sagte, manche Falschmeldungen hätten sogar Konflikte in Dörfern verursacht. Besonders vor Wahlen sei das Problem groß: Politiker oder ihre Anhänger erfänden Lügen über ihre Gegner.

Eine Journalistin aus Douala berichtete, ihre Redaktion habe ein Team gebildet, das jede verdächtige Nachricht kontrolliere. Sie betonte, die Wahrheit brauche manchmal Tage, aber eine Lüge verbreite sich in wenigen Minuten. Sie forderte, die Schulen müssten den Schülern beibringen, wie man Informationen überprüft. Nur so könne man gegen Fake News kämpfen. (environ 130 mots)
\end{textebox}

\textbf{Glossaire :} \all{der Bürger, -} (le citoyen) ; \all{vorsichtig} (prudent) ; \all{verursachen} (causer, provoquer) ; \all{die Wahl, -en} (l'élection) ; \all{der Anhänger, -} (le partisan) ; \all{der Gegner, -} (l'adversaire) ; \all{verdächtig} (suspect) ; \all{fordern} (exiger, réclamer) ; \all{beibringen} (enseigner, apprendre à).

\begin{exercice}[Compréhension écrite et questions de langue]
Lis le texte \all{„Fake News — eine Gefahr für Kamerun?“} ci-dessus, puis réponds aux questions \textbf{en allemand} par des phrases complètes.

\textbf{A. Compréhension (réponds en allemand)}
\begin{enumerate}
\item Wo verbreiten sich Falschmeldungen in Kamerun besonders schnell ?
\item Was tun viele Bürger nicht, bevor sie Nachrichten teilen ?
\item Wann ist das Problem der Fake News besonders groß ? Warum ?
\item Was hat die Redaktion der Journalistin aus Douala gemacht ?
\item Was fordert die Journalistin für die Schulen ?
\end{enumerate}

\textbf{B. Questions de langue}
\begin{enumerate}
\item Releve dans le texte trois verbes conjugués au subjonctif I et explique pourquoi ce mode est utilisé.
\item Mets au discours indirect (avec \all{dass}) : \all{„Wir kontrollieren jede verdächtige Nachricht.“} (Die Journalistin sagte, \ldots)
\item Trouve dans le texte le contraire de : \all{die Wahrheit} et \all{vorsichtig}.
\end{enumerate}
\end{exercice}

\begin{exercice}[Thème et version]
\textbf{A. Thème.} Traduis en allemand (attention au subjonctif I et à la place du verbe) :

\begin{enumerate}
\item Le journaliste a dit que les fausses informations se répandaient plus vite que la vérité.
\item Elle affirme qu'elle n'a pas menti et qu'elle a toujours vérifié ses sources.
\item Nous croyons que la source n'est pas fiable.
\end{enumerate}

\textbf{B. Version.} Traduis en français l'extrait suivant (attention au choix entre indicatif et conditionnel en français pour rendre le subjonctif I) :
\end{exercice}

\begin{textebox}[Aus einer Pressemitteilung]
Der Pressesprecher erklärte, die Regierung sei besorgt über die Verbreitung von Falschmeldungen. Man habe Beweise dafür, dass ausländische Gruppen Lügen verbreiteten. Er sagte, die Bürger sollten jede Nachricht überprüfen, und fügte hinzu, die Wahrheit werde immer siegen.
\end{textebox}

\begin{exercice}[Expression écrite type Bac]
\textbf{Sujet :} \all{Fake News in Kamerun : ein Problem?}

Tu es journaliste pour le journal de ton lycée. Rédige un paragraphe d'environ \textbf{8 lignes (80 à 100 mots)} dans lequel tu :
\begin{itemize}
\item expliques ce que sont les fake news et comment elles se répandent au Cameroun ;
\item rapportes les paroles d'au moins deux personnes (un expert, un élève, un enseignant\ldots) \textbf{au discours indirect} ;
\item donnes ton opinion et proposes une solution.
\end{itemize}

\textbf{Contraintes obligatoires :}
\begin{itemize}
\item utilise au moins \textbf{3 verbes au subjonctif I} ;
\item utilise au moins \textbf{2 subordonnées en} \all{dass} ;
\item emploie au moins 4 mots du vocabulaire de la leçon (\all{die Falschmeldung, die Lüge, die Wahrheit, der Beweis, die Quelle, verbreiten, überprüfen, glauben}).
\end{itemize}
\end{exercice}

\begin{methode}[Réussir ta production écrite]
\begin{enumerate}
\item \textbf{Planifie} : 2 phrases de présentation du problème, 2 phrases de discours rapporté, 2 phrases d'opinion et de solution.
\item \textbf{Discours indirect} : après \all{sagen, erklären, behaupten, meinen}, utilise le subjonctif I ; vérifie la place du verbe (fin avec \all{dass}, deuxième position sans \all{dass}).
\item \textbf{Vocabulaire} : choisis 4 mots de la liste et souligne-les dans ton brouillon pour vérifier la contrainte.
\item \textbf{Relis} : majuscules des noms, Umlaute, accord du verbe, ponctuation.
\end{enumerate}
\end{methode}

\begin{exemplebox}[Production modèle]
\all{Fake News sind Falschmeldungen, die sich besonders über WhatsApp und Facebook schnell verbreiten. Ein Experte aus Yaoundé erklärte, dass viele Bürger die Quelle einer Nachricht nicht überprüften. Er sagte auch, manche Lügen hätten schon Konflikte verursacht. Eine Schülerin meinte, sie glaube nicht mehr alles, was sie im Netz lese. Ich finde, dass jeder die Wahrheit suchen muss: Man soll jede Information prüfen, bevor man sie teilt. Die Schulen müssen den Schülern beibringen, wie man Fake News erkennt. Nur so kann man gegen dieses Problem kämpfen.} (environ 95 mots)

\emph{Traduction :} « Les fake news sont de fausses informations qui se répandent très vite, surtout via WhatsApp et Facebook. Un expert de Yaoundé a expliqué que beaucoup de citoyens ne vérifiaient pas la source d'une nouvelle. Il a dit aussi que certains mensonges avaient déjà provoqué des conflits. Une élève a dit qu'elle ne croyait plus tout ce qu'elle lisait sur le Net. Je pense que chacun doit chercher la vérité : on devrait vérifier toute information avant de la partager. Les écoles doivent apprendre aux élèves à reconnaître les fake news. C'est la seule façon de lutter contre ce problème. »
\end{exemplebox}

\begin{attention}[Pièges fréquents au Bac]
\begin{itemize}
\item Ne confonds pas subjonctif I (\all{er sei, er habe}) et subjonctif II (\all{er wäre, er hätte}) : le subjonctif II exprime l'irréel, pas le discours rapporté.
\item Sans \all{dass}, le verbe conjugué reste en \textbf{deuxième position} : \all{Sie sagte, sie habe nicht gelogen.} (et non \all{sie nicht gelogen habe}).
\item Au passé du discours indirect : \all{habe/sei} + participe passé : \all{Er sagte, er habe den Beweis gefunden.}
\item \all{die Lüge} (le mensonge) n'est pas \all{die Luge} ; attention à l'Umlaut. \all{beweisen} = « prouver », pas « prévenir ».
\end{itemize}
\end{attention}

\newpage

\coursec{Corrigés des exercices}

\courssub{Je vérifie mes connaissances}

\begin{corrige}[Exercice 1 — QCM : le subjonctif I]
\textbf{1. b)} — Le discours indirect en allemand se rend généralement par le \cle{subjonctif I} (Konjunktiv I). Le conditionnel français correspond plutôt au Konjunktiv II, qui exprime l'irréel ou l'hypothèse.

\textbf{2. b)} \all{Er sagte, er sei krank.} — \all{sei} est le subjonctif I de \all{sein} à la 3\up{e} personne du singulier. Sans \all{dass}, le verbe conjugué reste en deuxième position, comme dans une phrase principale. \all{ist} serait l'indicatif (le locuteur confirme le propos, sans distance) et \all{wäre} le subjonctif II (irréel : « il serait malade »).

\textbf{3. b)} \all{habe} — radical de l'infinitif \all{hab-} + désinence \all{-e}. \all{hätte} est le subjonctif II (irréel) et \all{hat} l'indicatif présent.

\textbf{4. c)} — Dans une subordonnée introduite par \all{dass}, le verbe conjugué se place \textbf{à la fin} : \all{Er sagte, dass er krank sei.} (« Il a dit qu'il était malade. »)

\textbf{5. b)} — \all{nicht gelogen} signifie « ne pas avoir menti » (participe passé du verbe fort \all{lügen, log, hat gelogen}). La phrase signifie donc : « La journaliste affirme qu'elle n'a pas menti. »
\end{corrige}

\begin{corrige}[Exercice 2 — Vrai ou faux ?]
\textbf{1.} \all{richtig} : \all{die Falschmeldung, -en} est bien un nom féminin (comme tous les noms en \all{-ung} : \all{die Meinung, die Zeitung, die Regierung}).

\textbf{2.} \all{richtig} : \all{der Beweis, -e} ; le pluriel est \all{die Beweise}.

\textbf{3.} \all{falsch} : \all{verbreiten} (verbe faible : \all{verbreitet, verbreitete, hat verbreitet}) signifie « répandre, diffuser ». « Vérifier » se dit \all{überprüfen} (verbe à particule inséparable : \all{hat überprüft}).

\textbf{4.} \all{richtig} : le subjonctif I présent de \all{sein} est irrégulier : \all{ich sei, du seist (seiest), er/sie/es sei, wir seien, ihr seiet, sie/Sie seien}.

\textbf{5.} \all{richtig} : sans \all{dass}, on conserve l'ordre de la phrase principale, verbe en deuxième position : \all{Er sagte, er sei krank.} (« Il a dit qu'il était malade. »)

\textbf{6.} \all{falsch} : \all{glauben} signifie « croire » et \all{beweisen} « prouver ». Ce ne sont pas des synonymes : croire une information ne signifie pas la prouver. Attention au faux ami : \all{beweisen} ne signifie jamais « prévenir ».
\end{corrige}

\begin{corrige}[Exercice 3 — Vocabulaire : trouve le mot]
\textbf{1.} \all{Bevor man eine Nachricht teilt, muss man die \textbf{Quelle überprüfen}.} — \all{überprüfen} reste à l'infinitif après le modal \all{muss} et se place en fin de phrase ; la particule \all{über-} est inséparable, donc pas de \all{ge-} au participe : \all{hat überprüft}.

\textbf{2.} \all{Er hat die \textbf{Wahrheit} gesagt : er war wirklich nicht dort.} — expression figée : \all{die Wahrheit sagen} (« dire la vérité »).

\textbf{3.} \all{Fake News \textbf{verbreiten} sich sehr schnell im Internet.} — sujet pluriel, donc \all{verbreiten} ; attention au verbe pronominal \all{sich verbreiten} (« se répandre ») : le réfléchi \all{sich} suit immédiatement le verbe conjugué.

\textbf{4.} \all{Die Polizei hat einen \textbf{Beweis} gefunden : das Foto war manipuliert.} — accusatif masculin : \all{einen Beweis}.

\textbf{5.} \all{Eine \textbf{Lüge} ist das Gegenteil von der Wahrheit.} — attention à l'Umlaut : \all{die Lüge, -n} ; le verbe correspondant est \all{lügen} (« mentir »).
\end{corrige}

\begin{corrige}[Exercice 4 — Conjugaison : tableau du subjonctif I]
\begin{center}
\begin{tabular}{|l|l|l|l|l|l|}
\hline
\rowcolor{popblueL} Personne & sein & haben & sagen & verbreiten & werden \\
\hline
ich & sei & habe & sage & verbreite & werde \\
\hline
du & seist & habest & sagest & verbreitest & werdest \\
\hline
er/sie/es & sei & habe & sage & verbreite & werde \\
\hline
wir & seien & haben & sagen & verbreiten & werden \\
\hline
ihr & seiet & habet & saget & verbreitet & werdet \\
\hline
sie/Sie & seien & haben & sagen & verbreiten & werden \\
\hline
\end{tabular}
\end{center}

\textbf{Règle :} on prend le radical de l'infinitif et on ajoute les désinences \all{-e, -est, -e, -en, -et, -en}. Seul \all{sein} est irrégulier, avec le radical \all{sei-} (sans \all{-e} à la 1\up{re} et 3\up{e} personne du singulier : \all{ich sei, er sei}).

\textbf{Remarque importante :} à \all{wir} et \all{sie/Sie}, le subjonctif I est identique à l'indicatif présent ; au discours indirect, on utilise alors la \textbf{forme de remplacement} (Ersatzform), c'est-à-dire le subjonctif II : \all{sie hätten, sie verbreiteten, sie würden}. C'est aussi le cas à la 1\up{re} personne du singulier quand la forme coïncide avec l'indicatif : \all{ich sage} (ind.) = \all{ich sage} (Konj. I) $\rightarrow$ on préfère souvent \all{ich würde sagen} ou on garde \all{ich sage} si aucune confusion n'est possible.
\end{corrige}

\courssub{Je m'entraîne}

\begin{corrige}[Exercice 5 — Du discours direct au discours indirect]
\textbf{1.} \all{Der Experte sagte, dass sich die Falschmeldungen schneller verbreiteten als die Wahrheit.} / \all{Der Experte sagte, die Falschmeldungen verbreiteten sich schneller als die Wahrheit.} — Le subjonctif I présent \all{verbreiten} (3\up{e} pers. pluriel) est identique à l'indicatif : on utilise la forme de remplacement \all{verbreiteten}. Sans \all{dass}, le verbe reste en deuxième position et le pronom réfléchi \all{sich} suit le verbe.

\textbf{2.} \all{Der Experte sagte, dass viele Jugendliche alles glaubten, was sie im Netz lesen.} / \all{\ldots, viele Jugendliche glaubten alles, was sie im Netz lesen.} — Même remplacement : \all{glaubten} au lieu de \all{glauben}. La subordonnée relative \all{was sie im Netz lesen} reste à l'indicatif : elle ne rapporte pas une parole, elle décrit un fait.

\textbf{3.} \all{Der Experte sagte, dass die Quelle dieser Nachricht nicht zuverlässig sei.} / \all{\ldots, die Quelle dieser Nachricht sei nicht zuverlässig.} — \all{sei} (3\up{e} pers. sing.) diffère de l'indicatif \all{ist} : c'est la forme pure du subjonctif I, la plus fréquente dans les journaux.

\textbf{4.} \all{Der Experte sagte, dass sie den Beweis gefunden hätten.} / \all{\ldots, sie hätten den Beweis gefunden.} — Le passé du discours indirect = subjonctif I de \all{haben} ou \all{sein} + participe passé ; ici \all{haben} serait identique à l'indicatif, d'où la forme de remplacement \all{hätten}.

\textbf{5.} \all{Der Experte sagte, dass die Regierung gegen Fake News kämpfen werde.} / \all{\ldots, die Regierung werde gegen Fake News kämpfen.} — Le futur rapporté : \all{werde} + infinitif (l'infinitif se place en fin de subordonnée, après le verbe conjugué relégué).

\textbf{6.} \all{Der Experte sagte, dass er jede Information vor dem Teilen überprüfe.} / \all{\ldots, er überprüfe jede Information vor dem Teilen.} — Le \all{ich} du discours direct devient \all{er} au discours rapporté (changement de personne) ; \all{überprüfe} (3\up{e} pers. sing.) diffère de l'indicatif \all{überprüft}.
\end{corrige}

\begin{corrige}[Exercice 6 — Texte à trous : Fake News]
\all{Ein Medienexperte erklärte in einem Interview, Fake News \textbf{seien} heute ein großes Problem. Viele Menschen \textbf{hätten} nicht mehr die Zeit, jede Nachricht zu \textbf{überprüfen}. Deshalb \textbf{verbreite} sich \textbf{die Lüge} oft schneller als die Wahrheit. Er sagte auch, die \textbf{Quelle} \textbf{sei} immer das Wichtigste. Wer sie nicht prüfe, \textbf{glaube} oft alles, was er im Internet lese.}

\emph{Traduction :} « Un expert des médias a expliqué dans une interview que les fake news étaient aujourd'hui un grand problème. Beaucoup de gens n'auraient plus le temps de vérifier chaque nouvelle. C'est pourquoi le mensonge se répand souvent plus vite que la vérité. Il a dit aussi que la source était toujours le plus important. Qui ne la vérifie pas croit souvent tout ce qu'il lit sur Internet. »

\textbf{Justifications :}
\begin{itemize}
\item \all{seien} : sujet pluriel (\all{Fake News}) ; \all{sei} au singulier serait une faute d'accord.
\item \all{hätten} : forme de remplacement, car \all{haben} (subjonctif I, 3\up{e} pers. pluriel) serait identique à l'indicatif.
\item \all{überprüfen} : infinitif avec \all{zu} après \all{die Zeit haben} ; la particule inséparable \all{über-} ne se sépare pas du \all{zu} : \all{zu überprüfen}.
\item \all{verbreite} : 3\up{e} pers. sing. du subjonctif I ; elle diffère de l'indicatif \all{verbreitet}, c'est donc bien la forme attendue.
\item \all{glaube} : 3\up{e} pers. sing. du subjonctif I (indicatif : \all{glaubt}).
\end{itemize}
\end{corrige}

\begin{corrige}[Exercice 7 — Appariement : mot et définition]
\textbf{1 -- b} : \all{die Falschmeldung, -en} = information inventée pour tromper le public.

\textbf{2 -- c} : \all{der Beweis, -e} = élément qui montre qu'un fait est vrai.

\textbf{3 -- d} : \all{die Quelle, -n} = origine d'une information (journal, site, témoin). Attention au faux ami : \all{die Quelle} signifie aussi « la source » d'un cours d'eau.

\textbf{4 -- a} : \all{lügen} = dire quelque chose qui n'est pas vrai. Attention : verbe fort \all{lügen, log, hat gelogen}.

\textbf{5 -- e} : \all{überprüfen} = contrôler si une information est exacte.
\end{corrige}

\begin{corrige}[Exercice 8 — Traduction : thème guidé]
\textbf{1.} \all{Der Journalist sagte, dass sich die Falschmeldungen schneller verbreiteten als die Wahrheit.} — Forme de remplacement \all{verbreiteten} (pluriel) ; le réfléchi \all{sich} se place juste après \all{dass}.

\textbf{2.} \all{Sie behauptet, dass sie nicht gelogen habe.} (ou sans \all{dass} : \all{Sie behauptet, sie habe nicht gelogen.}) — Passé du discours indirect : \all{habe} + participe passé \all{gelogen}.

\textbf{3.} \all{Wir glauben, dass die Quelle nicht zuverlässig ist.} — On rapporte ici notre propre opinion : l'indicatif est correct ; le subjonctif I \all{sei} marquerait une distance ou un doute.

\textbf{4.} \all{Der Minister erklärte, dass die Regierung gegen Fake News kämpfen werde.} — « lutterait » renvoie à un futur dans le discours rapporté : \all{werde} + infinitif \all{kämpfen} en fin de subordonnée.
\end{corrige}

\courssub{Je me prépare au Bac}

\begin{corrige}[Exercice 9 — Compréhension écrite et questions de langue]
\textbf{A. Compréhension}

\textbf{1.} \all{In Kamerun verbreiten sich Falschmeldungen besonders schnell über WhatsApp und Facebook.} (« Au Cameroun, les fausses informations se répandent particulièrement vite via WhatsApp et Facebook. »)

\textbf{2.} \all{Viele Bürger überprüfen die Quelle nicht, bevor sie Nachrichten teilen.} (« Beaucoup de citoyens ne vérifient pas la source avant de partager des nouvelles. »)

\textbf{3.} \all{Das Problem ist besonders vor Wahlen groß, weil Politiker oder ihre Anhänger Lügen über ihre Gegner erfinden.} (« Le problème est particulièrement important avant les élections, parce que les politiciens ou leurs partisans inventent des mensonges sur leurs adversaires. »)

\textbf{4.} \all{Die Redaktion hat ein Team gebildet, das jede verdächtige Nachricht kontrolliert.} (« La rédaction a formé une équipe qui contrôle chaque nouvelle suspecte. »)

\textbf{5.} \all{Sie fordert, dass die Schulen den Schülern beibringen, wie man Informationen überprüft.} (« Elle exige que les écoles apprennent aux élèves comment vérifier les informations. ») — Remarque : après \all{fordern, dass\ldots}, l'indicatif est ici accepté car il s'agit d'une exigence formulée clairement ; le subjonctif I (\all{beibrächten} en forme de remplacement) est également possible.

\textbf{B. Questions de langue}

\textbf{1.} Exemples de verbes au subjonctif I (ou formes de remplacement) : \all{seien}, \all{teilten}, \all{hätten \ldots\ verursacht}, \all{sei}, \all{erfänden}, \all{habe \ldots\ gebildet}, \all{kontrolliere}, \all{brauche}, \all{verbreite}, \all{könne}. Ce mode est utilisé parce que le texte \textbf{rapporte les paroles} de l'expert et de la journaliste : c'est le discours indirect, qui signale que le journaliste rapporte sans confirmer ni infirmer les propos.

\textbf{2.} \all{Die Journalistin sagte, dass sie jede verdächtige Nachricht kontrollierten.} — Ici \all{sie} désigne l'équipe (pluriel) : forme de remplacement obligatoire, car \all{kontrollieren} au subjonctif I (3\up{e} pers. pluriel) serait identique à l'indicatif. Si le sujet était la journaliste seule, on écrirait : \all{dass sie jede verdächtige Nachricht kontrolliere} (3\up{e} pers. sing.).

\textbf{3.} Contraire de \all{die Wahrheit} : \all{die Lüge}. Contraire de \all{vorsichtig} (« prudent ») : \all{unvorsichtig} ou \all{leichtsinnig} (« imprudent »).

\textbf{Barème indicatif (sur 20) :} A. Compréhension : 5 $\times$ 2 = 10 points (phrase complète, information exacte, allemand correct). B. Langue : question 1 = 4 points (au moins 3 verbes relevés + explication du mode), question 2 = 3 points, question 3 = 3 points.
\end{corrige}

\begin{corrige}[Exercice 10 — Thème et version]
\textbf{A. Thème}

\textbf{1.} \all{Der Journalist sagte, dass sich die Falschmeldungen schneller verbreiteten als die Wahrheit.} — Points de vigilance : forme de remplacement \all{verbreiteten}, place de \all{sich} après \all{dass}, verbe conjugué en fin de subordonnée.

\textbf{2.} \all{Sie behauptet, dass sie nicht gelogen habe und dass sie ihre Quellen immer überprüft habe.} — Deux subordonnées en \all{dass} coordonnées ; passé du discours indirect : \all{habe} + participe passé (\all{gelogen}, \all{überprüft} — sans \all{ge-} car particule inséparable).

\textbf{3.} \all{Wir glauben, dass die Quelle nicht zuverlässig ist.} — Opinion personnelle : indicatif correct ; \all{sei} possible pour marquer la réserve.

\textbf{B. Version}

« Le porte-parole a déclaré que le gouvernement était inquiet de la propagation des fausses informations. On aurait des preuves que des groupes étrangers répandaient des mensonges. Il a dit que les citoyens devraient vérifier chaque information, et il a ajouté que la vérité triompherait toujours. »

\textbf{Remarque de méthode :} le subjonctif I allemand se rend en français soit par l'indicatif dans une subordonnée introduite par « que » (discours rapporté neutre), soit par le \textbf{conditionnel journalistique} (« on aurait », « le gouvernement serait inquiet ») pour souligner que les propos ne sont pas confirmés. Les deux stratégies sont acceptées au Bac, à condition d'être cohérentes dans tout le texte.

\textbf{Barème indicatif (sur 20) :} Thème : 3 phrases $\times$ 4 = 12 points (mode, place du verbe, vocabulaire). Version : 8 points (fidélité du sens, qualité du français, rendu du subjonctif I).
\end{corrige}

\begin{corrige}[Exercice 11 — Expression écrite type Bac]
\textbf{Barème indicatif (sur 20) :}
\begin{itemize}
\item Respect de la tâche et du plan demandé (problème, deux paroles rapportées, opinion et solution) : 5 points.
\item Contraintes grammaticales : au moins 3 verbes au subjonctif I correctement formés : 4 points ; au moins 2 subordonnées en \all{dass} avec le verbe à la fin : 3 points.
\item Vocabulaire de la leçon (au moins 4 mots correctement employés) : 4 points.
\item Correction linguistique générale (orthographe, majuscules, Umlaute, accords, ponctuation) : 4 points.
\end{itemize}

\textbf{Production modèle :}

\all{Fake News sind Falschmeldungen, die sich besonders über WhatsApp und Facebook schnell verbreiten. Ein Experte aus Yaoundé erklärte, dass viele Bürger die Quelle einer Nachricht nicht überprüften. Er sagte auch, manche Lügen hätten schon Konflikte verursacht. Eine Schülerin meinte, sie glaube nicht mehr alles, was sie im Netz lese. Ich finde, dass jeder die Wahrheit suchen muss: Man soll jede Information prüfen, bevor man sie teilt. Die Schulen müssen den Schülern beibringen, wie man Fake News erkennt. Nur so kann man gegen dieses Problem kämpfen.} (environ 95 mots)

\emph{Traduction :} « Les fake news sont de fausses informations qui se répandent très vite, surtout via WhatsApp et Facebook. Un expert de Yaoundé a expliqué que beaucoup de citoyens ne vérifiaient pas la source d'une nouvelle. Il a dit aussi que certains mensonges avaient déjà provoqué des conflits. Une élève a dit qu'elle ne croyait plus tout ce qu'elle lisait sur le Net. Je pense que chacun doit chercher la vérité : on devrait vérifier toute information avant de la partager. Les écoles doivent apprendre aux élèves à reconnaître les fake news. C'est la seule façon de lutter contre ce problème. »

\textbf{Points de vigilance :}
\begin{itemize}
\item Subjonctif I et non subjonctif II : \all{sie glaube}, \all{hätten verursacht} (forme de remplacement), et non \all{sie würde glauben}.
\item Sans \all{dass}, le verbe reste en deuxième position : \all{sie glaube nicht mehr alles}.
\item Verbe conjugué à la fin dans toute subordonnée en \all{dass} : \all{dass viele Bürger die Quelle nicht überprüften}.
\item Majuscules à tous les noms (\all{Fake News, Quelle, Lügen, Wahrheit, Schulen, Schülern}) et Umlaut de \all{Lüge}.
\item Ne pas traduire « se répandre » par \all{spreizen} : le verbe correct est \all{sich verbreiten}.
\item Datif pluriel : \all{den Schülern} (avec \all{-n} au datif pluriel).
\end{itemize}
\end{corrige}

\newpage

\coursec{Fiche bilan}

\begin{popfigure}
\begin{tikzpicture}[mindmap, grow cyclic, every node/.style={concept, font=\small},
  root concept/.append style={concept color=popblue, font=\bfseries},
  level 1/.append style={level distance=3.4cm, sibling angle=51},
  level 1 concept/.append style={font=\footnotesize}]
\node[root concept, align=center]{Fake News\\ AL17}
  child[concept color=poporange]{ node[align=center]{Texte\\ \all{Fake News im Netz}} }
  child[concept color=popgreen]{ node[align=center]{Lexique\\ \all{die Lüge, die Wahrheit, verbreiten}} }
  child[concept color=poppurple]{ node[align=center]{Konjunktiv I\\ discours indirect\\ \all{er sagte, er habe...}} }
  child[concept color=poppink]{ node[align=center]{Subordonnée\\ \all{dass} + verbe final} }
  child[concept color=popgold]{ node[align=center]{Compréhension\\ globale puis détaillée} }
  child[concept color=popblue!70]{ node[align=center]{Expression\\ résumé, production} }
  child[concept color=popgreen!70]{ node[align=center]{Méthode\\ commentaire de texte} };
\end{tikzpicture}
\legende{Carte mentale de la leçon AL17 : la cybercriminalité et les fake news}
\end{popfigure}

\begin{bilan}
\textbf{1. Lexique clé (à connaître avec article et pluriel)}
\begin{itemize}
  \item \all{die Falschmeldung, -en} : la fausse information, la fake news
  \item \all{die Lüge, -n} : le mensonge ; \all{lügen (log, hat gelogen)} : mentir
  \item \all{die Wahrheit, -en} : la vérité ; \all{wahr} : vrai
  \item \all{verbreiten (verbreitete, hat verbreitet)} : diffuser, répandre
  \item \all{glauben (glaubte, hat geglaubt)} : croire ; \all{jemandem glauben} : croire quelqu'un (datif)
  \item \all{der Beweis, -e} : la preuve ; \all{beweisen (bewies, hat bewiesen)} : prouver
  \item \all{die Nachricht, -en} : l'information, la nouvelle ; \all{die Quelle, -n} : la source
  \item \all{prüfen / überprüfen} : vérifier ; \all{teilen} : partager (sur les réseaux)
  \item \all{das Gerücht, -e} : la rumeur ; \all{täuschen} : tromper
\end{itemize}

\textbf{2. Grammaire à connaître par cœur}
\begin{center}
\begin{tabularx}{\linewidth}{|X|X|}
\hline
\rowcolor{popblueL} Règle & Exemple \\
\hline
Discours indirect = \cle{Konjunktiv I} : on rapporte les paroles de quelqu'un sans les garantir. & \all{Er sagte, er \textbf{habe} die Nachricht nicht gelesen.} « Il a dit qu'il n'avait pas lu l'information. » \\
\hline
Konjunktiv I se forme sur le radical de l'infinitif : \all{ich sei, du seiest, er sei, wir seien, ihr seiet, sie seien}. & \all{Sie behauptet, sie \textbf{sei} unschuldig.} « Elle affirme être innocente. » \\
\hline
Si la forme du KI = indicatif, on utilise le \cle{Konjunktiv II} de remplacement. & \all{Sie sagten, sie \textbf{hätten} nichts gewusst.} (et non \all{haben}) \\
\hline
Subordonnée en \all{dass} : le verbe conjugué va à la \textbf{fin} de la proposition. & \all{Ich glaube, dass die Meldung falsch \textbf{ist}.} « Je crois que l'information est fausse. » \\
\hline
Discours indirect sans \all{dass} : verbe en position 2 ; avec \all{dass} : verbe en fin. & \all{Er sagte, er habe keine Zeit. / Er sagte, dass er keine Zeit habe.} \\
\hline
\end{tabularx}
\end{center}

\textbf{3. Méthodes express}
\begin{itemize}
  \item \textbf{Comprendre un texte} : 1) lire le titre et la source ; 2) lecture globale (qui ? quoi ? où ?) ; 3) lecture détaillée avec le glossaire ; 4) répondre par des phrases complètes en allemand.
  \item \textbf{Résumé} : une phrase par idée principale, au discours indirect (Konjunktiv I), sans citer le texte mot à mot.
  \item \textbf{Traduire au discours indirect} : verbe de parole (\all{sagen, behaupten, meinen, berichten}) + \all{dass} + verbe au Konjunktiv I en fin de phrase.
\end{itemize}
\end{bilan}

\begin{aretenir}[Checklist avant le Bac]
\begin{itemize}
  \item $\square$ Je connais l'article et le pluriel des noms du thème : \all{die Falschmeldung, -en ; die Lüge, -n ; die Wahrheit, -en ; der Beweis, -e ; das Gerücht, -e}.
  \item $\square$ Je sais conjuguer \all{sein} au Konjunktiv I : \all{ich sei, du seiest, er sei, wir seien, ihr seiet, sie seien}.
  \item $\square$ Je sais former le Konjunktiv I des autres verbes : \all{er habe, er gehe, er komme, er wisse}.
  \item $\square$ Je remplace par le Konjunktiv II quand le KI ressemble à l'indicatif : \all{sie hätten, sie gingen}.
  \item $\square$ Je mets le verbe conjugué à la fin dans toute subordonnée en \all{dass}.
  \item $\square$ Je sais rapporter une parole avec et sans \all{dass} (verbe en position 2 ou en fin).
  \item $\square$ Je distingue \all{glauben} (+ datif pour la personne) et \all{beweisen} (+ accusatif).
  \item $\square$ Je n'oublie pas la majuscule à tous les noms : \all{die Wahrheit, der Beweis, das Netz}.
  \item $\square$ Je sais écrire un résumé de 4 à 5 phrases au discours indirect.
  \item $\square$ Je connais les verbes de parole : \all{sagen, behaupten, berichten, meinen, erklären}.
  \item $\square$ Je méfie des faux amis : \all{die Nachricht} = la nouvelle/l'information, pas « le message moral » ; \all{bekommen} = recevoir, pas « devenir ».
  \item $\square$ Je relis ma production : place du verbe, cas, orthographe des Umlaute (ä, ö, ü) et du ß.
\end{itemize}
\end{aretenir}

\findecours

\end{document}
